% Complex Numbers — Pure Mathematics with Statistics Upper Sixth — Cours signé OnBuch+
% Official MINESEC syllabus. Compiler avec : tectonic PMS02-complex-numbers.tex
\newcommand{\DOCMATIERE}{Pure Mathematics with Statistics}
\newcommand{\DOCNIVEAU}{Upper Sixth}
\newcommand{\DOCMODULE}{Upper Sixth — Pure mathematics with statistics}
\newcommand{\DOCLECON}{2}
\newcommand{\DOCTITRE}{Complex Numbers}
% OnBuch+ — préambule « cours signés » ENRICHI (compilé avec tectonic / XeLaTeX)
% Système visuel « Pop » d'OnBuch+ : fond crème (jamais blanc), bordures encre
% épaisses, ombres « dures » décalées, titres Archivo Black en capitales.
% Version MATHÉMATIQUES : graphes (pgfplots/tikz), boîtes pédagogiques
% (définition, propriété, méthode, attention, le savais-tu, exercice résolu,
% exercice, TP, bilan), page de garde et signature OnBuch+ ancrée en bas.
%
% Macros attendues dans main.tex AVANT \input{preamble} :
%   \DOCMATIERE  \DOCNIVEAU  \DOCMODULE  \DOCLECON (numéro)  \DOCTITRE
\documentclass[11pt]{article}

\usepackage{fontspec}
\usepackage[english]{babel}
\usepackage[a4paper,margin=2cm,top=3.4cm,bottom=2.8cm,headheight=1.4cm,footskip=1.3cm]{geometry}
\usepackage{amsmath,amssymb,amsfonts}
\usepackage{mathtools} % \xmapsto, \coloneqq... souvent utilisés par les modèles
\usepackage{cancel} % simplifications barrées (bilans de forces)
% \abs{z} (module, valeur absolue) et \norm{u} (norme), utilisés par les modèles
\providecommand{\abs}{}\let\abs\relax\DeclarePairedDelimiter{\abs}{\lvert}{\rvert}
\providecommand{\norm}{}\let\norm\relax\DeclarePairedDelimiter{\norm}{\lVert}{\rVert}
\usepackage[table]{xcolor} % [table] : fournit aussi \rowcolors (alternance de lignes)
\usepackage{pagecolor}
\usepackage{tikz}
\usetikzlibrary{arrows.meta,positioning,calc,shapes.geometric,shapes.misc,shapes.arrows,decorations.pathmorphing,decorations.pathreplacing,decorations.markings,patterns,shadows,mindmap,trees,fit,backgrounds,matrix,intersections,angles,quotes}
\usepackage{pgfplots}
\usepackage[version=4]{mhchem} % équations nucléaires \ce{^{235}_{92}U}
\usepackage[european,siunitx]{circuitikz} % schémas électriques
\pgfplotsset{compat=1.18}
\usepgfplotslibrary{fillbetween,groupplots} % aires entre courbes (calcul intégral)
\usepackage{enumitem}
\usepackage{fancyhdr}
% [most] charge toutes les bibliothèques tcolorbox (listings, minted, raster,
% documentation...) et gonfle inutilement la mémoire TeX sur un document long
% (~300 boîtes) : on ne charge que ce qui est réellement utilisé.
\usepackage[skins,breakable]{tcolorbox}
\usepackage{graphicx}
\usepackage{colortbl}
\usepackage{booktabs}
\usepackage{tabularx}
\usepackage{diagbox} % case d'en-tête diagonale des tableaux à double entrée
% Colonnes extensibles centrée / gauche / droite (C, L, R), souvent utilisées par les modèles
\newcolumntype{C}{>{\centering\arraybackslash}X}
\newcolumntype{L}{>{\raggedright\arraybackslash}X}
\newcolumntype{R}{>{\raggedleft\arraybackslash}X}
\usepackage{array}
\usepackage{multicol}
\usepackage{siunitx}
% parse-numbers=false : accepte aussi \SI{2,5 \times 10^{-3}}{...}
\sisetup{locale=UK,output-decimal-marker={.},group-minimum-digits=4,parse-numbers=false}
\DeclareSIUnit\ohm{\ensuremath{\symup{\Omega}}} % U+2126 absent de la police
\DeclareSIUnit\division{div} % réglages d'oscilloscope (ms/div, V/div)
\DeclareSIUnit\tr{tr} % tours (vitesse de rotation, ex. \SI{1500}{\tr\per\minute})
\usepackage{qrcode}
% \degree : commande fréquemment utilisée par les modèles pour un angle ou une
% température en dehors de \SI{}{} (non fournie par les paquets chargés).
\providecommand{\degree}{^{\circ}}
% \qed (fin de démonstration), utilisé par les modèles sans amsthm
\providecommand{\qed}{\hfill$\square$}
% \barre : double barre des valeurs interdites dans un tableau de variations (tkz-tab)
\providecommand{\barre}{\ensuremath{\|}}
% environnement proof (amsthm non chargé) utilisé par les modèles
\ifdefined\proof\else\newenvironment{proof}[1][Proof]{\par\noindent\textit{#1.}\ }{\qed\par}\fi
\usepackage{lastpage}
\usepackage{hyperref}
\hypersetup{hidelinks}

% ============================================================
% Palette « Pop » OnBuch+ — mêmes hex que la classe P (plus_ui.dart)
% ============================================================
\definecolor{poporange}{HTML}{F59321}
\definecolor{popdark}{HTML}{E07A0C}
\definecolor{popcream}{HTML}{FFF6E8}
\definecolor{popink}{HTML}{1C1714}
\definecolor{poppurple}{HTML}{7A5AE0}
\definecolor{popgreen}{HTML}{2E9E6A}
\definecolor{poppink}{HTML}{E0457B}
\definecolor{popblue}{HTML}{2F7DE1}
\definecolor{popgold}{HTML}{C98A12}
\definecolor{popmuted}{HTML}{8A7F76}
\definecolor{popline}{HTML}{EADFD2}
\definecolor{popgray}{HTML}{8A7F76}
\colorlet{popgrey}{popgray}
% Alias tolérants (le modèle nomme parfois les couleurs différemment)
\colorlet{popwhite}{white}
\colorlet{popblack}{popink}
\colorlet{popred}{poppink}
\colorlet{popyellow}{popgold}
% Teintes claires pour fonds de tableaux / zones
\colorlet{popblueL}{popblue!10!white}
\colorlet{poporangeL}{poporange!14!white}
\colorlet{popgreenL}{popgreen!12!white}
\colorlet{poppurpleL}{poppurple!12!white}
\colorlet{poppinkL}{poppink!12!white}
\colorlet{popgoldL}{popgold!14!white}

\pagecolor{popcream}

% ============================================================
% Typographie
% ============================================================
\defaultfontfeatures{Ligatures=TeX}
\setmainfont{PlusJakartaSans-Medium.ttf}[Path=../../../fonts/,BoldFont=PlusJakartaSans-Bold.ttf]
% Maths en sans-serif (Fira Math) avec chiffres et lettres droites de Plus
% Jakarta, pour que les formules se fondent dans le texte.
\usepackage[math-style=ISO,bold-style=ISO]{unicode-math}
\setmathfont{FiraMath-Regular.otf}[Path=../../../fonts/]
\setmathfont{PlusJakartaSans-Medium.ttf}[Path=../../../fonts/,range={up/num,up/{Latin,latin},\mathup}]
% (icomma retiré : décimales avec point en anglais)
% Caractères mathématiques Unicode absents des polices de texte (Plus Jakarta,
% Archivo Black) : rendus par leur équivalent mathématique, en texte comme en
% formule — sinon ils s'affichent en carré vide (ex. « Arithmétique dans ℤ »).
\usepackage{newunicodechar}
\newunicodechar{ℝ}{\ensuremath{\mathbb{R}}}
\newunicodechar{ℤ}{\ensuremath{\mathbb{Z}}}
\newunicodechar{ℂ}{\ensuremath{\mathbb{C}}}
\newunicodechar{ℕ}{\ensuremath{\mathbb{N}}}
\newunicodechar{ℚ}{\ensuremath{\mathbb{Q}}}
\newunicodechar{𝔻}{\ensuremath{\mathbb{D}}}
\newunicodechar{α}{\ensuremath{\alpha}}
\newunicodechar{β}{\ensuremath{\beta}}
\newunicodechar{γ}{\ensuremath{\gamma}}
\newunicodechar{δ}{\ensuremath{\delta}}
\newunicodechar{ε}{\ensuremath{\varepsilon}}
\newunicodechar{θ}{\ensuremath{\theta}}
\newunicodechar{λ}{\ensuremath{\lambda}}
\newunicodechar{μ}{\ensuremath{\mu}}
\newunicodechar{σ}{\ensuremath{\sigma}}
\newunicodechar{τ}{\ensuremath{\tau}}
\newunicodechar{φ}{\ensuremath{\varphi}}
\newunicodechar{ψ}{\ensuremath{\psi}}
\newunicodechar{ω}{\ensuremath{\omega}}
\newunicodechar{Σ}{\ensuremath{\Sigma}}
% majuscules produites par \MakeUppercase dans les titres (α-aminés -> Α)
\newunicodechar{Α}{\ensuremath{\alpha}}
\newunicodechar{Β}{\ensuremath{\beta}}
\newunicodechar{Γ}{\ensuremath{\Gamma}}
\newunicodechar{Δ}{\ensuremath{\Delta}}
\newunicodechar{Ε}{\ensuremath{\varepsilon}}
\newunicodechar{Θ}{\ensuremath{\Theta}}
\newunicodechar{Λ}{\ensuremath{\Lambda}}
\newunicodechar{Μ}{\ensuremath{\mu}}
\newunicodechar{Τ}{\ensuremath{\tau}}
\newunicodechar{Φ}{\ensuremath{\Phi}}
\newunicodechar{Ψ}{\ensuremath{\Psi}}
\newunicodechar{Ω}{\ensuremath{\Omega}}
\newunicodechar{↦}{\ensuremath{\mapsto}}
\newunicodechar{∈}{\ensuremath{\in}}
\newunicodechar{∉}{\ensuremath{\notin}}
\newunicodechar{∧}{\ensuremath{\wedge}}
\newunicodechar{⇔}{\ensuremath{\Leftrightarrow}}
\newunicodechar{⟺}{\ensuremath{\Longleftrightarrow}}
\newunicodechar{⇒}{\ensuremath{\Rightarrow}}
\newunicodechar{⟹}{\ensuremath{\Longrightarrow}}
\newunicodechar{∘}{\ensuremath{\circ}}
\newunicodechar{∪}{\ensuremath{\cup}}
\newunicodechar{∩}{\ensuremath{\cap}}
\newunicodechar{≡}{\ensuremath{\equiv}}
\newunicodechar{⊂}{\ensuremath{\subset}}
\newunicodechar{⊥}{\ensuremath{\perp}}
\newunicodechar{∃}{\ensuremath{\exists}}
\newunicodechar{∀}{\ensuremath{\forall}}
\newunicodechar{∛}{\ensuremath{\sqrt[3]{\,}}}
\newunicodechar{∜}{\ensuremath{\sqrt[4]{\,}}}
\newunicodechar{⃗}{\ensuremath{{}^{\to}}}
\newunicodechar{ⁿ}{\ifmmode^{n}\else\textsuperscript{\ensuremath{n}}\fi}
\newunicodechar{ˣ}{\ifmmode^{x}\else\textsuperscript{\ensuremath{x}}\fi}
\newunicodechar{ᵇ}{\ifmmode^{b}\else\textsuperscript{\ensuremath{b}}\fi}
\newunicodechar{ʳ}{\ifmmode^{r}\else\textsuperscript{\ensuremath{r}}\fi}
\newunicodechar{ᵘ}{\ifmmode^{u}\else\textsuperscript{\ensuremath{u}}\fi}
\newunicodechar{ʸ}{\ifmmode^{y}\else\textsuperscript{\ensuremath{y}}\fi}
\newunicodechar{ᵃ}{\ifmmode^{a}\else\textsuperscript{\ensuremath{a}}\fi}
\newunicodechar{ᵉ}{\ifmmode^{e}\else\textsuperscript{\ensuremath{e}}\fi}
\newunicodechar{ᵏ}{\ifmmode^{k}\else\textsuperscript{\ensuremath{k}}\fi}
\newunicodechar{ᵖ}{\ifmmode^{p}\else\textsuperscript{\ensuremath{p}}\fi}
\newunicodechar{ᴺ}{\ifmmode^{N}\else\textsuperscript{\ensuremath{N}}\fi}
\newunicodechar{ᶜ}{\ifmmode^{c}\else\textsuperscript{\ensuremath{c}}\fi}
\newunicodechar{ʰ}{\ifmmode^{h}\else\textsuperscript{\ensuremath{h}}\fi}
\newunicodechar{ᵠ}{\ifmmode^{q}\else\textsuperscript{\ensuremath{q}}\fi}
\newunicodechar{ₙ}{\ifmmode_{n}\else\textsubscript{\ensuremath{n}}\fi}
\newunicodechar{ₐ}{\ifmmode_{a}\else\textsubscript{\ensuremath{a}}\fi}
\newunicodechar{ᵢ}{\ifmmode_{i}\else\textsubscript{\ensuremath{i}}\fi}
\newunicodechar{ᵣ}{\ifmmode_{r}\else\textsubscript{\ensuremath{r}}\fi}
\newunicodechar{ₖ}{\ifmmode_{k}\else\textsubscript{\ensuremath{k}}\fi}
\newunicodechar{ᵦ}{\ifmmode_{\beta}\else\textsubscript{\ensuremath{\beta}}\fi}
\newunicodechar{ₓ}{\ifmmode_{x}\else\textsubscript{\ensuremath{x}}\fi}
\newfontfamily\popdisplay{ArchivoBlack-Regular.ttf}[Path=../../../fonts/]
\newfontfamily\poplogo{SpaceGrotesk-Bold.ttf}[Path=../../../fonts/]
\newfontfamily\popsemi{PlusJakartaSans-SemiBold.ttf}[Path=../../../fonts/]
\newfontfamily\popxbold{PlusJakartaSans-ExtraBold.ttf}[Path=../../../fonts/]
\newfontfamily\popxboldls{PlusJakartaSans-ExtraBold.ttf}[Path=../../../fonts/,LetterSpace=3.0]

\setlist[enumerate]{leftmargin=1.7em,itemsep=5pt,topsep=4pt}
\setlist[itemize]{leftmargin=1.7em,itemsep=4pt,topsep=4pt,label={\color{poporange}\textbullet}}
\setlength{\parindent}{0pt}
\setlength{\parskip}{5pt}
\renewcommand{\arraystretch}{1.25}

% pgfplots : style Pop par défaut
\pgfplotsset{
  every axis/.append style={axis line style={popink,line width=1pt},tick style={popink},
    grid style={popline,line width=0.5pt},label style={font=\small},tick label style={font=\footnotesize}},
  popcurve/.style={poporange,line width=1.8pt,no markers,smooth},
}
% Même style disponible dans un \draw TikZ simple (hors pgfplots)
\tikzset{popcurve/.style={poporange,line width=1.8pt}}
\tikzset{popgrid/.style={popline,very thin}}
\colorlet{popcurve}{poporange} % le nom du style est parfois utilisé comme couleur

% ============================================================
% Pill / squiggle
% ============================================================
\newcommand{\poppill}[3]{%
  \tikz[baseline=-0.55ex]\node[draw=popink,line width=1.2pt,rounded corners=2.6pt,%
    fill=#2,inner xsep=6.5pt,inner ysep=3pt]{%
    \popxboldls\fontsize{7}{7}\selectfont\color{#3}\MakeUppercase{#1}};%
}
\newcommand{\popsquiggle}[1]{%
  \tikz[baseline=-0.4ex]{
    \draw[#1,line width=2.6pt,line cap=round]
      plot [smooth] coordinates {(0,0.09) (0.34,0.01) (0.68,0.21) (1.02,0.01) (1.36,0.10)};
  }%
}
% Mot-clé surligné (marqueur orange sous le texte)
\newcommand{\cle}[1]{{\popxbold\color{popdark}#1}}

% ============================================================
% En-tête / pied
% ============================================================
\pagestyle{fancy}
\fancyhf{}
\renewcommand{\headrulewidth}{0pt}
\renewcommand{\footrulewidth}{0pt}
\lhead{\raisebox{-0.15ex}{\poplogo\fontsize{13}{13}\selectfont\color{popink}OnBuch\color{poporange}+}}
\rhead{\poppill{\DOCMATIERE\ \textbullet\ \DOCNIVEAU\ \textbullet\ Chapter \DOCLECON}{white}{popink}}
\lfoot{\popxbold\fontsize{7.6}{7.6}\selectfont\color{popmuted}\MakeUppercase{Signed course OnBuch+}\ \textbullet\ {\popsemi\DOCTITRE}}
\rfoot{\tikz[baseline=-0.6ex]\node[rounded corners=8.5pt,fill=popink,minimum height=17pt,minimum width=17pt,inner xsep=5pt,inner sep=0]{\popxbold\fontsize{8}{8}\selectfont\color{popcream}\thepage\,/\,\pageref*{LastPage}};}

% ============================================================
% Titres
% ============================================================
\newcommand{\chaptitre}[2]{%
  \begin{tcolorbox}[enhanced,colback=poporange,colframe=popink,boxrule=2.6pt,arc=11pt,
      drop shadow=popink,left=15pt,right=15pt,top=12pt,bottom=11pt,before skip=3pt,after skip=18pt]
    \poppill{#2}{popink}{popcream}\par\vspace{9pt}
    {\raggedright\hyphenpenalty=10000\exhyphenpenalty=10000\popdisplay\fontsize{22}{24}\selectfont\color{popcream}\MakeUppercase{#1}\par}
  \end{tcolorbox}
}
% Section numérotée : pastille orange avec le numéro + titre Archivo
\newcounter{popsec}
\newcommand{\coursec}[1]{%
  \stepcounter{popsec}\setcounter{popsub}{0}%
  \par\vspace{16pt}\noindent
  \begin{minipage}{\linewidth}
  \tikz[baseline=-0.7ex]\node[draw=popink,line width=1.6pt,fill=poporange,rounded corners=4pt,minimum size=22pt,inner sep=0]{\popdisplay\fontsize{12}{12}\selectfont\color{popcream}\thepopsec};%
  \hspace{8pt}{\popdisplay\fontsize{15}{17}\selectfont\color{popink}\MakeUppercase{#1}}\par
  \vspace{4pt}\noindent{\color{poporange}\rule{36pt}{3.6pt}}
  \end{minipage}\par\nopagebreak\vspace{8pt}
}
\newcounter{popsub}
\newcommand{\courssub}[1]{%
  \stepcounter{popsub}%
  \par\vspace{9pt}\noindent{\popxbold\fontsize{11.5}{13}\selectfont\color{popink}{\color{poporange}\thepopsec.\thepopsub}\hspace{6pt}#1}\par\nopagebreak\vspace{4pt}
}

% ============================================================
% Boîtes pédagogiques — même recette : fond blanc, bordure épaisse colorée,
% ombre dure encre, badge plein. Titre optionnel [#1].
% ============================================================
\tcbset{popbase/.style={breakable,enhanced,colback=white,boxrule=2.3pt,arc=9pt,
  shadow={3pt}{-3pt}{0pt}{popink},left=14pt,right=14pt,top=11pt,bottom=11pt,before skip=11pt,after skip=15pt,
  fontupper=\popsemi\fontsize{10.6}{14.5}\selectfont\color{popink},
  fontlower=\popsemi\fontsize{10.6}{14.5}\selectfont\color{popink}}}
% Coche dessinée (le glyphe ✓ n'existe pas dans les polices du document)
\newcommand{\popcheck}[1]{\tikz[baseline=-0.5ex]\draw[#1,line width=1.6pt,line cap=round,line join=round](-0.13,0.0)--(-0.04,-0.09)--(0.13,0.1);}
\providecommand{\coche}{\popcheck{popgreen}}
\providecommand{\resultat}[1]{\par\smallskip\noindent\fcolorbox{popgreen}{popgreenL}{\parbox{\dimexpr\linewidth-2\fboxsep-2\fboxrule\relax}{\textbf{Result.} #1}}\par\smallskip}
\newcommand{\popboxhead}[3]{\poppill{#1}{#2}{white}\ifx\relax#3\relax\else\hspace{6pt}{\popxbold\fontsize{10.6}{12}\selectfont\color{popink}#3}\fi\par\vspace{7pt}}

\newtcolorbox{aretenir}[1][]{popbase,colframe=popgreen,before upper={\popboxhead{Key points}{popgreen}{#1}}}
\newtcolorbox{exemplebox}[1][]{popbase,colframe=poporange,before upper={\popboxhead{Exemple}{poporange}{#1}}}
\newtcolorbox{definition}[1][]{popbase,colframe=popblue,before upper={\popboxhead{Definition}{popblue}{#1}}}
\newtcolorbox{propriete}[1][]{popbase,colframe=poppurple,before upper={\popboxhead{Property}{poppurple}{#1}}}
\newtcolorbox{methode}[1][]{popbase,colframe=popgold,before upper={\popboxhead{Method}{popgold}{#1}}}
\newtcolorbox{attention}[1][]{popbase,colframe=poppink,before upper={\popboxhead{Warning}{poppink}{#1}}}
\newtcolorbox{experience}[1][]{popbase,colframe=popblue,colback=popblueL,before upper={\popboxhead{Experiment}{popblue}{#1}}}
% « Le savais-tu ? » : carte violette pleine (contexte, culture, Cameroun)
\newtcolorbox{savaistu}[1][]{popbase,colback=poppurple,colframe=popink,
  fontupper=\popsemi\fontsize{10.4}{14.2}\selectfont\color{white},
  before upper={\poppill{Did you know?}{white}{poppurple}\ifx\relax#1\relax\else\hspace{6pt}{\popxbold\color{white}#1}\fi\par\vspace{7pt}}}
% Exercice résolu : énoncé en haut, \tcblower puis la solution sur fond crème
\newcounter{popexo}\newcounter{popexr}
\newtcolorbox{exoresolu}[1][]{popbase,colframe=poporange,colbacklower=popcream,
  segmentation style={popink,line width=1.2pt,dashed},
  before upper={\stepcounter{popexr}\popboxhead{Worked example \thepopexr}{poporange}{#1}},
  before lower={\poppill{Solution}{popink}{popcream}\par\vspace{6pt}}}
% Exercice (non résolu en place) — la correction est dans la partie Corrigés
\newtcolorbox{exercice}[1][]{popbase,colframe=popink,
  before upper={\stepcounter{popexo}\popboxhead{Exercise \thepopexo}{popink}{#1}}}
% Corrigé
\newtcolorbox{corrige}[1][]{popbase,colframe=popgreen,colback=popgreenL,
  before upper={\popboxhead{Solution}{popgreen}{#1}}}
% Objectifs / prérequis / bilan
\newtcolorbox{objectifs}{popbase,colframe=popink,colback=poporangeL,
  before upper={\popboxhead{Chapter objectives}{popink}{}}}
\newtcolorbox{prerequis}{popbase,colframe=popink,colback=popblueL,
  before upper={\popboxhead{Prerequisites}{popblue}{}}}
\newtcolorbox{bilan}{popbase,colframe=popink,colback=white,boxrule=2.8pt,
  before upper={\popboxhead{Fiche bilan}{poporange}{L'essentiel à connaître pour l'examen}}}
% Figure encadrée avec légende
\newtcolorbox{popfigure}[1][]{enhanced,breakable=false,colback=white,colframe=popline,boxrule=1.4pt,arc=8pt,
  left=10pt,right=10pt,top=10pt,bottom=8pt,before skip=10pt,after skip=12pt,halign=center,
  lower separated=false,halign lower=center,
  fontlower=\popsemi\fontsize{9}{11.5}\selectfont\color{popmuted},
  #1}
\newcommand{\legende}[1]{\par\vspace{6pt}{\popsemi\fontsize{9}{11.5}\selectfont\color{popmuted}#1}}

% ============================================================
% Signature OnBuch+
% ============================================================
\newcommand{\poplogoinline}[1]{{\poplogo\fontsize{#1}{#1}\selectfont\color{popink}OnBuch\color{poporange}+}}
\newcommand{\signatureonbuch}{%
  \begin{tcolorbox}[enhanced,colback=popink,colframe=popink,boxrule=2.2pt,arc=10pt,
      drop shadow=poporange,left=14pt,right=14pt,top=10pt,bottom=10pt,before skip=0pt,after skip=0pt]
    \tikz[baseline=-0.7ex]\node[circle,fill=popgreen,draw=popcream,line width=1.3pt,minimum size=22pt,inner sep=0]{\popcheck{white}};%
    \hspace{9pt}\begin{minipage}[c]{0.62\linewidth}
      {\popxbold\fontsize{12}{13}\selectfont\color{popcream}Signed\ \textbullet\ {\poplogo OnBuch\color{poporange}+}}\par\vspace{2pt}
      {\raggedright\popsemi\fontsize{8}{10.5}\selectfont\color{popline}\MakeUppercase{OnBuch+ teaching team}\\\MakeUppercase{Follows the official MINESEC syllabus}\par}
    \end{minipage}\hfill
    {\poplogo\fontsize{22}{22}\selectfont\color{popcream}OnBuch\color{poporange}+}
  \end{tcolorbox}%
}

% ============================================================
% Page de garde
% #1 titre · #2 matière · #3 niveau · #4 module · #5 n° leçon · #6 durée ·
% #7 description / objectifs (texte ou itemize)
% ============================================================
\newcommand{\pagegarde}[7]{%
  \thispagestyle{empty}
  \newgeometry{a4paper,margin=1.7cm,top=1.6cm,bottom=1.4cm}%
  \begin{tikzpicture}[remember picture,overlay]
    \fill[poporange!18] ([xshift=-3.2cm,yshift=-3.6cm]current page.north east) circle (4.6cm);
    \fill[poppurple!14] ([xshift=2.4cm,yshift=7.6cm]current page.south west) circle (3.8cm);
  \end{tikzpicture}%
  {\poplogo\fontsize{30}{30}\selectfont\color{popink}OnBuch\color{poporange}+}\hfill
  \raisebox{0.6ex}{\poppill{Signed course \textbullet\ Premium}{popink}{popcream}}\par
  \vspace{1pt}\popsquiggle{poppurple}\par
  \vspace{8pt}
  \begin{tcolorbox}[enhanced,colback=poporange,colframe=popink,boxrule=2.8pt,arc=13pt,
      drop shadow=popink,left=18pt,right=18pt,top=12pt,bottom=13pt,after skip=10pt]
    \poppill{#2 \textbullet\ #3}{popink}{popcream}\hspace{5pt}\poppill{#4}{white}{popink}\par\vspace{8pt}
    {\popdisplay\fontsize{14}{14}\selectfont\color{popink}CHAPTER #5}\par\vspace{4pt}
    {\raggedright\hyphenpenalty=10000\exhyphenpenalty=10000\popdisplay\fontsize{25}{27}\selectfont\color{popcream}\MakeUppercase{#1}\par}
    \vspace{8pt}
    {\popxbold\fontsize{9.5}{9.5}\selectfont\color{popink}\MakeUppercase{Suggested time: #6}}
  \end{tcolorbox}
  \begin{tcolorbox}[enhanced,colback=white,colframe=popink,boxrule=2.2pt,arc=10pt,
      drop shadow=popink,left=16pt,right=16pt,top=10pt,bottom=10pt,after skip=8pt]
    \poppill{In this chapter}{poppurple}{white}\par\vspace{5pt}
    \popsemi\fontsize{9.3}{12.6}\selectfont\color{popink}#7
  \end{tcolorbox}
  \noindent\begin{minipage}[t]{0.487\linewidth}
    \begin{tcolorbox}[enhanced,colback=popgreen,colframe=popink,boxrule=2.2pt,arc=10pt,
        drop shadow=popink,left=11pt,right=11pt,top=10pt,bottom=10pt,height=3.1cm,valign=center]
      \tcbox[colback=white,colframe=popink,boxrule=1.3pt,arc=5pt,boxsep=0pt,nobeforeafter,
        left=3pt,right=3pt,top=3pt,bottom=3pt,valign=center]{\qrcode[height=1.9cm]{https://play.google.com/store/apps/details?id=com.luvvix.onbuch&hl=en}}%
      \hspace{8pt}\begin{minipage}[c]{0.5\linewidth}
        {\popdisplay\fontsize{11}{12.5}\selectfont\color{white}\MakeUppercase{Download OnBuch}}\par\vspace{4pt}
        {\raggedright\popsemi\fontsize{8.4}{11}\selectfont\color{white}Free \textbullet\ Google Play\\AI tutor, past papers, results\par}
      \end{minipage}
    \end{tcolorbox}
  \end{minipage}\hfill
  \begin{minipage}[t]{0.487\linewidth}
    \begin{tcolorbox}[enhanced,colback=poppurple,colframe=popink,boxrule=2.2pt,arc=10pt,
        drop shadow=popink,left=11pt,right=11pt,top=10pt,bottom=10pt,height=3.1cm,valign=center]
      \tikz[baseline=-0.7ex]\node[circle,fill=white,draw=popink,line width=1.3pt,minimum size=1.6cm,inner sep=0]{\popdisplay\fontsize{24}{24}\selectfont\color{poppurple}+};%
      \hspace{8pt}\begin{minipage}[c]{0.6\linewidth}
        {\popdisplay\fontsize{11}{12.5}\selectfont\color{white}\MakeUppercase{Go OnBuch+}}\par\vspace{4pt}
        {\raggedright\popsemi\fontsize{8.4}{11}\selectfont\color{white}All signed courses, unlimited AI tutor, PDF sheets — OnBuch+ tab in the app\par}
      \end{minipage}
    \end{tcolorbox}
  \end{minipage}
  \vspace{8pt}
  \popcontenu
  \vfill
  \signatureonbuch
  \restoregeometry
  \newpage
}

% Rangée « Ce cours contient » (page de garde)
\newcommand{\poptile}[3]{% #1 couleur #2 gros texte #3 légende
  \begin{tcolorbox}[enhanced,colback=#1,colframe=popink,boxrule=2pt,arc=9pt,shadow={2.5pt}{-2.5pt}{0pt}{popink},
      left=6pt,right=6pt,top=9pt,bottom=9pt,halign=center,valign=center,height=1.75cm,width=\linewidth,nobeforeafter]
    {\popdisplay\fontsize{12.5}{13}\selectfont\color{white}\MakeUppercase{#2}}\par\vspace{3pt}
    {\centering\popsemi\fontsize{7.8}{9.5}\selectfont\color{white}#3\par}
  \end{tcolorbox}}
\newcommand{\popcontenu}{%
  \par\noindent{\popxboldls\fontsize{7.5}{7.5}\selectfont\color{popmuted}THIS SIGNED COURSE CONTAINS}\par\vspace{7pt}
  \noindent\begin{minipage}[t]{0.235\linewidth}\poptile{popblue}{Course}{complete, illustrated, step by step}\end{minipage}\hfill
  \begin{minipage}[t]{0.235\linewidth}\poptile{poporange}{Methods}{and worked examples}\end{minipage}\hfill
  \begin{minipage}[t]{0.235\linewidth}\poptile{poppink}{Exercises}{exam style + solutions}\end{minipage}\hfill
  \begin{minipage}[t]{0.235\linewidth}\poptile{popgreen}{Summary sheet}{the essentials to remember}\end{minipage}\par}

% Sommaire
\newtcolorbox{sommaire}{enhanced,colback=white,colframe=popink,boxrule=2.2pt,arc=10pt,
  drop shadow=popink,left=18pt,right=18pt,top=13pt,bottom=13pt,before skip=6pt,after skip=20pt,
  before upper={\poppill{Contents}{poppurple}{white}\par\vspace{9pt}\popsemi\fontsize{10.6}{14.5}\selectfont\color{popink}}}

% Signature de fin de cours — poussée tout en bas de la dernière page
\newcommand{\findecours}{%
  \par\vfill\nopagebreak
  \begin{center}\popsquiggle{poporange}\end{center}\vspace{4pt}
  \signatureonbuch
}
\AtBeginDocument{\def\llbracket{\mathopen{[\mkern-3.5mu[}}\def\rrbracket{\mathclose{]\mkern-3.5mu]}}} % intervalles d'entiers (glyphe ⟦ absent de la police)
% Glyphes absents de Fira Math : redéfinis à partir de glyphes disponibles
\AtBeginDocument{%
  \def\setminus{\mathbin{\backslash}}\let\smallsetminus\setminus
  \def\vdots{\mathord{\vbox{\baselineskip3.2pt\lineskiplimit0pt\kern4pt\hbox{.}\hbox{.}\hbox{.}}}}%
  \def\ddots{\mathinner{\mkern1mu\raise6.5pt\hbox{.}\mkern2mu\raise3.6pt\hbox{.}\mkern2mu\raise0.7pt\hbox{.}\mkern1mu}}%
  \def\triangle{\mathord{\scalebox{0.9}{$\symup{\Delta}$}}}%
  \def\nabla{\mathord{\raisebox{\depth}{\scalebox{1}[-1]{$\symup{\Delta}$}}}}%
  \def\checkmark{\popcheck{popdark}}%
  \def\star{\mathbin{\ast}}\def\bigstar{\mathord{\ast}}%
  \def\diamond{\mathbin{\scalebox{0.6}{$\Diamond$}}}%
  \def\bigcup{\mathop{\mathchoice{\raisebox{-0.3em}{\scalebox{1.7}{$\cup$}}}{\scalebox{1.25}{$\cup$}}{\cup}{\cup}}\displaylimits}%
  \def\bigcap{\mathop{\mathchoice{\raisebox{-0.3em}{\scalebox{1.7}{$\cap$}}}{\scalebox{1.25}{$\cap$}}{\cap}{\cap}}\displaylimits}%
  \def\lesseqqgtr{\mathrel{\lessgtr}}\def\bigoplus{\mathop{\oplus}\displaylimits}%
}
\providecommand{\ssi}{if and only if} % abréviation fréquente
\AtBeginDocument{\providecommand{\wideparen}{\overparen}} % arc de cercle

% Fonctions usuelles que les modèles utilisent sans que amsmath les fournisse
\providecommand{\arsinh}{\operatorname{arsinh}}\providecommand{\arcosh}{\operatorname{arcosh}}
\providecommand{\artanh}{\operatorname{artanh}}\providecommand{\arsech}{\operatorname{arsech}}
\providecommand{\arcsch}{\operatorname{arcsch}}\providecommand{\arcoth}{\operatorname{arcoth}}
\providecommand{\arcsinh}{\operatorname{arsinh}}\providecommand{\arccosh}{\operatorname{arcosh}}
\providecommand{\arctanh}{\operatorname{artanh}}\providecommand{\sech}{\operatorname{sech}}
\providecommand{\csch}{\operatorname{csch}}\providecommand{\arcsec}{\operatorname{arcsec}}
\providecommand{\arccsc}{\operatorname{arccsc}}\providecommand{\arccot}{\operatorname{arccot}}
\providecommand{\sgn}{\operatorname{sgn}}\providecommand{\lcm}{\operatorname{lcm}}
\providecommand{\Var}{\operatorname{Var}}\providecommand{\Cov}{\operatorname{Cov}}

%% FIGFIX-BEGIN v5 — correctifs globaux des figures (docs/onbuch-generation/tools/figfix)
\usepackage{adjustbox}\usepackage{etoolbox}\tcbuselibrary{raster}
% toute figure TikZ/pgfplots plus large que la ligne est réduite à la largeur disponible
\BeforeBeginEnvironment{tikzpicture}{\begin{adjustbox}{max width=\linewidth}}
\AfterEndEnvironment{tikzpicture}{\end{adjustbox}}
% légendes pgfplots lisibles par-dessus les courbes
\pgfplotsset{every axis/.append style={legend style={fill=white,fill opacity=0.92,text opacity=1,draw=black!15,inner sep=3pt}}}
% étiquettes d'arêtes (syntaxe edge["..."]) avec fond blanc
\tikzset{every edge quotes/.append style={fill=white,inner sep=1.2pt}}
%% FIGFIX-END
\begin{document}

\pagegarde{Complex Numbers}{Pure Mathematics with Statistics}{Upper Sixth}{Upper Sixth — Pure mathematics with statistics}{2}{19 periods}{This chapter introduces complex numbers, the extension of the real number system built on the imaginary unit i. Students learn to compute with complex numbers in Cartesian, polar and exponential forms, represent them on the Argand diagram, apply De Moivre's theorem, extract roots, and interpret loci geometrically — all core skills for the GCE Advanced Level examination.
\vspace{4pt}
\begin{multicols}{2}\small
\begin{itemize}
\item By the end of this chapter I can explain the historical need for complex numbers and define the imaginary unit i with its powers.
\item By the end of this chapter I can add, subtract, multiply and divide complex numbers in Cartesian form.
\item By the end of this chapter I can use the complex conjugate and its properties to simplify expressions and solve equations.
\item By the end of this chapter I can represent a complex number on an Argand diagram and interpret its modulus and argument geometrically.
\item By the end of this chapter I can convert a complex number between Cartesian, polar (trigonometric) and exponential forms.
\item By the end of this chapter I can apply De Moivre's theorem to compute powers and to relate multiple angles to powers of cos x and sin x.
\end{itemize}\end{multicols}}

\begin{sommaire}
\begin{multicols}{2}
\begin{enumerate}
\item Birth of Complex Numbers and the Imaginary Unit i
\item Algebra of Complex Numbers and the Conjugate
\item Equations and Roots in ℂ
\item The Argand Diagram, Modulus and Argument
\item Polar and Exponential Forms
\item De Moivre's Theorem and Trigonometric Applications
\item Roots of Complex Numbers: nth Roots and Cube Roots
\item Loci in the Complex Plane and Synthesis
\item Integration activity
\item Methods and worked examples
\item Exercises
\item Solutions to the exercises
\item Summary sheet
\end{enumerate}
\end{multicols}
\end{sommaire}

\begin{objectifs}
\begin{itemize}
\item By the end of this chapter I can explain the historical need for complex numbers and define the imaginary unit i with its powers.
\item By the end of this chapter I can add, subtract, multiply and divide complex numbers in Cartesian form.
\item By the end of this chapter I can use the complex conjugate and its properties to simplify expressions and solve equations.
\item By the end of this chapter I can represent a complex number on an Argand diagram and interpret its modulus and argument geometrically.
\item By the end of this chapter I can convert a complex number between Cartesian, polar (trigonometric) and exponential forms.
\item By the end of this chapter I can apply De Moivre's theorem to compute powers and to relate multiple angles to powers of cos x and sin x.
\item By the end of this chapter I can find square roots and nth roots (including cube roots) of complex numbers.
\item By the end of this chapter I can solve polynomial equations with complex roots, including those with real coefficients.
\item By the end of this chapter I can describe and sketch loci defined by conditions on |z| and arg(z).
\end{itemize}
\end{objectifs}

\begin{prerequis}
\begin{itemize}
\item Real number algebra: expanding, factorising and solving quadratic equations (Lower Sixth, PMS01).
\item Trigonometric ratios, standard angles and identities such as cos²x + sin²x = 1 (Lower Sixth trigonometry).
\item Radian measure and the unit circle.
\item Basic coordinate geometry: distance between two points, equations of circles and perpendicular bisectors.
\item Surds and rationalising denominators (useful technique transferred to conjugates).
\end{itemize}
\end{prerequis}

\newpage

\coursec{Birth of Complex Numbers and the Imaginary Unit i}

An engineer in Douala analysing the alternating current supplied by ENEO models the impedance of a circuit by a number whose square is negative — yet no real number has a negative square. Likewise, when a student at the University of Yaoundé I asks a calculator for the square root of $-9$, the machine replies \texttt{ERROR}, as if such a number could not exist. For centuries, mathematicians like Cardano and Bombelli faced the same wall while solving cubic equations. Their bold answer was to invent a new number, $i$, whose square is $-1$. This chapter opens the door to that invisible but powerful world: the complex numbers, which today govern electricity, signal processing and even the mobile phone networks connecting Bamenda to Buea.

\courssub{Historical Motivation: From Cubic Equations to the Imaginary Unit}

\begin{popfigure}
\begin{tikzpicture}[scale=0.9, every node/.style={font=\small}]
% Timeline axis
\draw[->, thick] (0,0) -- (14,0) node[right] {Time};
% Tick marks and labels
\foreach \x/\year/\name in {
  1/1545/Cardano,
  3.5/1572/Bombelli,
  6/1748/Euler,
  8.5/1797/Wessel,
  11/1806/Argand,
  13/1831/Gauss
}{
  \draw (\x,0.1) -- (\x,-0.1) node[below=2pt] {\year};
  \node[above=4pt, align=center] at (\x,0.6) {\name};
}
% Descriptions
\node[align=center, text width=2.2cm] at (1,1.6) {Cardano publishes \emph{Ars Magna}; meets $\sqrt{-121}$ while solving $x^{3}=15x+4$.};
\node[align=center, text width=2.2cm] at (3.5,1.6) {Bombelli develops rules for ``imaginary'' numbers; computes $(2+i)^{3}$.};
\node[align=center, text width=2.2cm] at (6,1.6) {Euler introduces the symbol $i$ for $\sqrt{-1}$; $e^{i\pi}+1=0$.};
\node[align=center, text width=2.2cm] at (8.5,1.6) {Wessel gives the geometric interpretation (vectors).};
\node[align=center, text width=2.2cm] at (11,1.6) {Argand publishes the complex plane independently.};
\node[align=center, text width=2.2cm] at (13,1.6) {Gauss gives a rigorous foundation; coins ``complex number''.};
\end{tikzpicture}
\legende{Historical timeline of complex numbers: from Cardano's cubic to Gauss's rigorous foundation.}
\end{popfigure}

The need for complex numbers arose not so much from quadratic equations like $x^{2}+1=0$ (which were simply declared ``impossible'') as from \textbf{cubic equations}. In 1545, Gerolamo Cardano published \emph{Ars Magna}, in which the formula for solving $x^{3}+px+q=0$ sometimes required manipulating square roots of negative numbers — even when the final answer was a perfectly real number. For instance, solving $x^{3}=15x+4$ leads to
\[
x=\sqrt[3]{2+\sqrt{-121}}+\sqrt[3]{2-\sqrt{-121}}.
\]
Cardano called such expressions ``sophistic'' and ``useless''. It was Rafael Bombelli (1572) who dared to treat $\sqrt{-1}$ as a formal symbol obeying $(\sqrt{-1})^{2}=-1$. He showed that the two cube roots above are $2+i$ and $2-i$, whose sum gives the real solution $x=4$ — and indeed $4^{3}=64=15\times 4+4$. Later, Leonhard Euler introduced the notation $i=\sqrt{-1}$ and discovered the profound identity $e^{i\pi}+1=0$. The geometric representation by Caspar Wessel (1797), Jean-Robert Argand (1806) and Carl Friedrich Gauss turned these ``imaginary'' numbers into points of a plane, finally granting them full mathematical citizenship.

\begin{savaistu}[Why ``imaginary''?]
The term ``imaginary'' was coined by Descartes in 1637 as a derogatory label — he considered such numbers ``imagined'' rather than ``real''. Gauss later preferred the word ``complex'' because a number $a+ib$ has two ``parts''. Today, complex numbers are as real as any other mathematical object; they are indispensable in physics, engineering and pure mathematics.
\end{savaistu}

\courssub{The Imaginary Unit $i$ and its Algebra}

\begin{definition}[Imaginary Unit]
The \cle{imaginary unit} $i$ is the number defined by the property
\[
i^{2}=-1.
\]
Equivalently, we write $i=\sqrt{-1}$. Any real multiple of $i$ (i.e. $bi$ with $b\in\mathbb{R}$) is called a \cle{purely imaginary number}.
\end{definition}

From $i^{2}=-1$ we immediately obtain the first powers:
\[
i^{1}=i,\qquad i^{2}=-1,\qquad i^{3}=i^{2}\cdot i=-i,\qquad i^{4}=i^{2}\cdot i^{2}=(-1)(-1)=1.
\]
Because $i^{4}=1$, the powers of $i$ repeat every four steps: for any integer $n$,
\[
i^{n+4}=i^{n}\cdot i^{4}=i^{n}.
\]

\begin{propriete}[Periodicity of Powers of $i$]
The sequence $(i^{n})_{n\in\mathbb{N}}$ is periodic with period $4$:
\[
i^{4k}=1,\qquad i^{4k+1}=i,\qquad i^{4k+2}=-1,\qquad i^{4k+3}=-i\qquad(k\in\mathbb{N}).
\]
The proof (a direct application of $i^{4}=1$) is immediate and may be required as a short justification in the examination.
\end{propriete}

\begin{popfigure}
\begin{tikzpicture}[scale=1.2]
\draw[thick, ->] (-1.6,0) -- (1.6,0) node[right] {$\Re$};
\draw[thick, ->] (0,-1.6) -- (0,1.6) node[above] {$\Im$};
\draw[thick, popblue] (0,0) circle (1);
\foreach \angle/\lbl in {90/i, 180/-1, 270/-i, 0/1}{
  \draw[poporange, thick, ->] (0,0) -- (\angle:1);
  \node[popdark] at (\angle:1.3) {$\lbl$};
}
\draw[popgreen, thick, ->] (0.5,0) arc (0:90:0.5);
\node[popgreen, font=\footnotesize] at (45:0.75) {$\times i$};
\draw[popgreen, thick, ->] (0,0.5) arc (90:180:0.5);
\node[popgreen, font=\footnotesize] at (135:0.75) {$\times i$};
\draw[popgreen, thick, ->] (-0.5,0) arc (180:270:0.5);
\node[popgreen, font=\footnotesize] at (225:0.75) {$\times i$};
\draw[popgreen, thick, ->] (0,-0.5) arc (270:360:0.5);
\node[popgreen, font=\footnotesize] at (315:0.75) {$\times i$};
\end{tikzpicture}
\legende{The cycle of powers of $i$ on the unit circle: each multiplication by $i$ rotates by $90^{\circ}$ counter‑clockwise.}
\end{popfigure}

\begin{methode}[Evaluating $i^{n}$ for any integer $n$]
\begin{enumerate}
\item Divide $n$ by $4$ and find the remainder $r$ with $0\le r\le 3$. (For negative $n$, add multiples of $4$ to the exponent until it becomes non‑negative, using $i^{4}=1$.)
\item Then $i^{n}=i^{r}$ using the table:
\begin{center}
\begin{tabular}{|c|c|c|c|c|}
\hline
\rowcolor{popblueL} $r$ & $0$ & $1$ & $2$ & $3$ \\
\hline
$i^{r}$ & $1$ & $i$ & $-1$ & $-i$ \\
\hline
\end{tabular}
\end{center}
\end{enumerate}
\end{methode}

\begin{exoresolu}[Quick evaluation of a large power]
Evaluate $i^{2025}$.
\tcblower
Divide $2025$ by $4$: $2025 = 4\times 506 + 1$, so the remainder is $1$.
Hence $i^{2025}=i^{1}=i$.
\end{exoresolu}

\begin{attention}[The square‑root trap]
It is tempting to write $\sqrt{-a}\,\sqrt{-b}=\sqrt{(-a)(-b)}=\sqrt{ab}$ for $a,b>0$. This is \textbf{false}, because the rule $\sqrt{x}\,\sqrt{y}=\sqrt{xy}$ is valid only for $x,y\ge 0$ in $\mathbb{R}$. The correct procedure is:
\[
\sqrt{-a}=i\sqrt{a},\qquad \sqrt{-b}=i\sqrt{b}\quad\Longrightarrow\quad
\sqrt{-a}\,\sqrt{-b}=i\sqrt{a}\cdot i\sqrt{b}=i^{2}\sqrt{ab}=-\sqrt{ab}.
\]
Always convert $\sqrt{-a}$ to $i\sqrt{a}$ \emph{before} multiplying.
\end{attention}

\courssub{Complex Numbers: Definition and Basic Properties}

\begin{definition}[Complex Number]
A \cle{complex number} is any expression of the form
\[
z = a + ib,
\]
where $a,b\in\mathbb{R}$ and $i^{2}=-1$. The real number $a$ is called the \cle{real part} of $z$, denoted $\Re(z)$ or $\operatorname{Re}(z)$; the real number $b$ is called the \cle{imaginary part} of $z$, denoted $\Im(z)$ or $\operatorname{Im}(z)$. The form $a+ib$ is called the \cle{Cartesian} (or \cle{algebraic}) \cle{form} of $z$.
\end{definition}

Thus every complex number corresponds uniquely to an ordered pair $(a,b)\in\mathbb{R}^{2}$. Two complex numbers are equal if and only if their real parts and their imaginary parts are respectively equal:
\[
a+ib = c+id \;\Longleftrightarrow\; a=c \ \text{and}\ b=d.
\]

Special cases:
\begin{itemize}
\item If $b=0$, then $z=a\in\mathbb{R}$; we say $z$ is \cle{purely real}.
\item If $a=0$, then $z=ib$ with $b\in\mathbb{R}$; we say $z$ is \cle{purely imaginary}.
\item The number $0 = 0+0i$ is both purely real and purely imaginary.
\end{itemize}

The set of all complex numbers is denoted $\mathbb{C}$:
\[
\mathbb{C} = \{ a+ib \mid a,b\in\mathbb{R} \}.
\]
It contains the real numbers as a subset (those with $b=0$), so we have the chain of inclusions
\[
\mathbb{N} \subset \mathbb{Z} \subset \mathbb{Q} \subset \mathbb{R} \subset \mathbb{C}.
\]

\begin{popfigure}
\begin{tikzpicture}[scale=1.1]
% Nested ellipses for number sets
\draw[popblue, thick] (0,0) ellipse (4.5 and 3.5);
\node[popblue] at (3.2,2.6) {$\mathbb{C}$};
\draw[popgreen, thick] (0,-0.2) ellipse (3.4 and 2.6);
\node[popgreen] at (2.4,1.9) {$\mathbb{R}$};
\draw[poporange, thick] (0,-0.4) ellipse (2.4 and 1.8);
\node[poporange] at (1.7,1.2) {$\mathbb{Q}$};
\draw[poppurple, thick] (0,-0.6) ellipse (1.5 and 1.1);
\node[poppurple] at (1.05,0.35) {$\mathbb{Z}$};
\draw[poppink, thick] (0,-0.8) ellipse (0.85 and 0.55);
\node[poppink] at (0,-0.8) {$\mathbb{N}$};
% Sample elements
\node[popdark] at (-3.3,2.4) {$3+2i$};
\node[popdark] at (-2.6,1.3) {$\sqrt{2}$};
\node[popdark] at (-1.8,0.1) {$-\frac{5}{3}$};
\node[popdark] at (-0.9,-1.6) {$-7$};
\node[popdark] at (0.9,-2.2) {$4$};
\end{tikzpicture}
\legende{Nested number sets: $\mathbb{N}\subset\mathbb{Z}\subset\mathbb{Q}\subset\mathbb{R}\subset\mathbb{C}$. Each larger set contains the numbers needed to solve more equations.}
\end{popfigure}

\begin{aretenir}[Key Points of Section 1]
\begin{itemize}
\item The imaginary unit $i$ satisfies $i^{2}=-1$; it was introduced to solve equations like $x^{2}+1=0$ and to make the cubic formulas work.
\item Powers of $i$ cycle every $4$: $i,\, -1,\, -i,\, 1$. To compute $i^{n}$, find the remainder of $n$ upon division by $4$.
\item A complex number has the form $z=a+ib$ with $a,b\in\mathbb{R}$; $a=\Re(z)$, $b=\Im(z)$.
\item Equality: $a+ib=c+id \iff a=c$ and $b=d$.
\item The hierarchy of number sets: $\mathbb{N}\subset\mathbb{Z}\subset\mathbb{Q}\subset\mathbb{R}\subset\mathbb{C}$.
\item \textbf{Warning:} $\sqrt{-a}\,\sqrt{-b}\neq\sqrt{ab}$ in general; always write $\sqrt{-a}=i\sqrt{a}$ first.
\end{itemize}
\end{aretenir}

\courssub{Worked Examples and Common Pitfalls}

\begin{exoresolu}[Solving a simple quadratic in $\mathbb{C}$]
Solve $x^{2}+4=0$ in $\mathbb{C}$.
\tcblower
$x^{2}=-4 \implies x=\pm\sqrt{-4}=\pm i\sqrt{4}=\pm 2i$.
Check: $(2i)^{2}=4i^{2}=4(-1)=-4$ and $(-2i)^{2}=4i^{2}=-4$. Both are solutions: $S=\{-2i,\,2i\}$.
\end{exoresolu}

\begin{exoresolu}[Simplifying an expression with powers of $i$]
Simplify $\dfrac{i^{17}+i^{23}}{i^{6}}$.
\tcblower
Reduce each exponent modulo $4$:
\begin{align*}
17 &= 4\times 4+1 &\Longrightarrow\quad i^{17}&=i,\\
23 &= 4\times 5+3 &\Longrightarrow\quad i^{23}&=-i,\\
6  &= 4\times 1+2 &\Longrightarrow\quad i^{6}&=-1.
\end{align*}
Hence
\[
\frac{i^{17}+i^{23}}{i^{6}}=\frac{i+(-i)}{-1}=\frac{0}{-1}=0.
\]
\end{exoresolu}

\begin{exoresolu}[Using the equality of two complex numbers]
Find the real numbers $x$ and $y$ such that $(x+2y)+i(3x-y)=5+5i$.
\tcblower
By equality of the real and imaginary parts:
\[
\begin{cases}
x+2y=5\\
3x-y=5
\end{cases}
\]
From the second equation, $y=3x-5$. Substituting into the first: $x+2(3x-5)=5$, so $7x=15$, giving $x=\dfrac{15}{7}$ and $y=3\times\dfrac{15}{7}-5=\dfrac{45-35}{7}=\dfrac{10}{7}$.
Check: $\dfrac{15}{7}+2\times\dfrac{10}{7}=\dfrac{35}{7}=5$ and $3\times\dfrac{15}{7}-\dfrac{10}{7}=\dfrac{35}{7}=5$. \quad$\blacksquare$
\end{exoresolu}

\begin{exercice}[Practice: Powers of $i$ and Complex Forms]
\begin{enumerate}
\item Evaluate: (a) $i^{2024}$ \quad (b) $i^{-13}$ \quad (c) $i^{100}+i^{101}+i^{102}+i^{103}$.
\item Write each number in the form $a+ib$ and identify $\Re(z)$ and $\Im(z)$:
\begin{enumerate}
\item $5$ \quad (b) $-3i$ \quad (c) $\sqrt{-18}$ \quad (d) $(2+i)(1+i)$.
\end{enumerate}
\item True or False? Justify.
\begin{enumerate}
\item $\sqrt{-9}\cdot\sqrt{-4}=6$.
\item Every real number is a complex number.
\item The imaginary part of $3-4i$ is $-4i$.
\end{enumerate}
\item Find real numbers $x$ and $y$ such that $2x+iy=4-3i$.
\end{enumerate}
\end{exercice}

\begin{corrige}[Exercise 1]
\begin{enumerate}
\item (a) $2024=4\times506+0$, remainder $0 \Rightarrow i^{2024}=1$.
(b) $-13 = 4\times(-4)+3$, so $i^{-13}=i^{3}=-i$. (Equivalently $i^{-13}=1/i^{13}=1/i=-i$, since $i\times(-i)=-i^{2}=1$.)
(c) The four consecutive powers give $i^{100}=1$, $i^{101}=i$, $i^{102}=-1$, $i^{103}=-i$; the sum is $1+i-1-i=0$.
\item (a) $5=5+0i$: $\Re(z)=5$, $\Im(z)=0$.
(b) $-3i=0-3i$: $\Re(z)=0$, $\Im(z)=-3$.
(c) $\sqrt{-18}=i\sqrt{18}=3\sqrt{2}\,i$: $\Re(z)=0$, $\Im(z)=3\sqrt{2}$.
(d) $(2+i)(1+i)=2+2i+i+i^{2}=1+3i$: $\Re(z)=1$, $\Im(z)=3$.
\item (a) False: $\sqrt{-9}\,\sqrt{-4}=3i\cdot 2i=6i^{2}=-6\neq 6$.
(b) True: any $a\in\mathbb{R}$ can be written $a+0i\in\mathbb{C}$.
(c) False: the imaginary part is the \emph{real} number $-4$, not $-4i$.
\item Equating parts: $2x=4$ and $y=-3$, so $x=2$, $y=-3$.
\end{enumerate}
\end{corrige}

\begin{attention}[GCE Multiple‑Choice Trap]
In a multiple‑choice question you may be asked: ``What is $i^{2025}$?'' with options $i$, $-1$, $-i$, $1$. The quick method: $2025 = 4\times 506+1$, so the answer is $i$. Never attempt to multiply $i$ by itself $2025$ times! Also remember: the imaginary part of $a+ib$ is the \textbf{real} number $b$, not $ib$.
\end{attention}

\begin{savaistu}[Complex Numbers in Cameroon's Technology]
The mobile phone networks that connect Yaoundé, Douala, Bamenda and Buea rely on complex numbers for signal modulation, error‑correcting codes and antenna design. Alternating current, such as that distributed by ENEO, is analysed using \emph{phasors} — complex numbers representing amplitude and phase. Medical imaging techniques such as MRI use complex Fourier transforms to reconstruct images. The ``imaginary'' unit $i$ is therefore very real in everyday Cameroonian life!
\end{savaistu}

\coursec{Algebra of Complex Numbers and the Conjugate}

In the previous section, we discovered how the need to solve equations like $x^2+1=0$ led mathematicians to introduce the \cle{imaginary unit} $i$, defined by $i^2=-1$. We saw that every complex number can be written uniquely in the \cle{standard (or Cartesian) form} $z=a+ib$, where $a,b\in\mathbb{R}$. Now we must learn how to \emph{calculate} with these numbers: add, subtract, multiply, divide, and manipulate their conjugates. These algebraic skills are the bedrock of every later topic — polar form, De Moivre's theorem, loci, and the solution of polynomial equations in $\mathbb{C}$. Master them now, and the rest of the chapter will flow naturally.

\courssub{Addition and Subtraction of Complex Numbers}

Adding or subtracting complex numbers is completely analogous to adding vectors in the plane: we simply combine the real parts and the imaginary parts separately.

\begin{definition}[Addition and Subtraction in $\mathbb{C}$]
Let $z_1=a+ib$ and $z_2=c+id$ with $a,b,c,d\in\mathbb{R}$. Then
\[
z_1\pm z_2 = (a\pm c) + i(b\pm d).
\]
\end{definition}

\begin{exemplebox}[Addition and Subtraction]
Let $z_1=3+4i$ and $z_2=-2+7i$. Compute $z_1+z_2$ and $z_1-z_2$.
\begin{align*}
z_1+z_2 &= (3+(-2)) + i(4+7) = 1+11i,\\
z_1-z_2 &= (3-(-2)) + i(4-7) = 5-3i.
\end{align*}
\end{exemplebox}

\begin{aretenir}[Key Point]
The set $\mathbb{C}$ is \emph{closed} under addition and subtraction: the sum or difference of two complex numbers is always a complex number. The operations are commutative and associative, just as in $\mathbb{R}$.
\end{aretenir}

\courssub{Multiplication of Complex Numbers}

Multiplication uses the distributive law (FOIL) together with the defining rule $i^2=-1$.

\begin{definition}[Multiplication in $\mathbb{C}$]
For $z_1=a+ib$ and $z_2=c+id$,
\[
z_1z_2 = (a+ib)(c+id) = ac + iad + ibc + i^2bd = (ac-bd) + i(ad+bc).
\]
\end{definition}

\begin{exemplebox}[Multiplication]
Multiply $(2+3i)$ by $(4-i)$.
\begin{align*}
(2+3i)(4-i) &= 2\cdot4 + 2(-i) + 3i\cdot4 + 3i(-i)\\
&= 8 -2i +12i -3i^2\\
&= 8 +10i -3(-1)\\
&= 11+10i.
\end{align*}
\end{exemplebox}

\begin{attention}[Classic Trap: Squaring a Complex Number]
Many students write $(a+ib)^2 = a^2 + b^2$. \textbf{This is false!} The correct expansion is
\[
(a+ib)^2 = a^2 + 2iab + i^2b^2 = (a^2-b^2) + 2iab.
\]
For example, $(2+3i)^2 = 4 + 12i + 9i^2 = -5+12i$, \emph{not} $4+9=13$.
\end{attention}

\courssub{The Complex Conjugate}

The \cle{complex conjugate} is a simple but extraordinarily powerful operation: it reflects a complex number across the real axis.

\begin{popfigure}
\begin{tikzpicture}[scale=1.2]
% Axes
\draw[->] (-3.5,0) -- (3.5,0) node[right] {$\mathrm{Re}$};
\draw[->] (0,-3) -- (0,3) node[above] {$\mathrm{Im}$};
% Grid lines
\foreach \x in {-3,-2,-1,1,2,3} \draw[gray!30] (\x,-2.5) -- (\x,2.5);
\foreach \y in {-2,-1,1,2} \draw[gray!30] (-3,\y) -- (3,\y);
% Points
\coordinate (z) at (2,1.5);
\coordinate (zbar) at (2,-1.5);
% Vectors
\draw[popblue, thick, ->] (0,0) -- (z) node[midway, above left] {$z=a+ib$};
\draw[poporange, thick, ->] (0,0) -- (zbar) node[midway, below left] {$\bar{z}=a-ib$};
% Dashed reflection line
\draw[dashed, popmuted] (2,1.5) -- (2,-1.5);
% Labels
\node[popblue] at (2.5,1.5) {$z$};
\node[poporange] at (2.5,-1.5) {$\bar{z}$};
\node[popmuted] at (3.2,0) {Reflection across $\mathrm{Re}$ axis};
% Real axis label
\node[below] at (3.5,0) {Real axis};
\end{tikzpicture}
\legende{The conjugate $\bar{z}$ is the mirror image of $z$ in the real axis.}
\end{popfigure}

\begin{definition}[Complex Conjugate]
For $z=a+ib$ ($a,b\in\mathbb{R}$), the \cle{complex conjugate} of $z$ is
\[
\bar{z} = a-ib.
\]
On the Argand diagram, $\bar{z}$ is the reflection of $z$ across the real axis.
\end{definition}

\begin{propriete}[Fundamental Identity: $z\bar{z}$ is Real and Non-Negative]
\[
z\bar{z} = (a+ib)(a-ib) = a^2 - i^2b^2 = a^2+b^2 \in \mathbb{R}_{\ge 0}.
\]
This product is exactly the \emph{square of the modulus} of $z$ (denoted $|z|^2$), which we will study in Section 4. \emph{Proof is examinable.}
\end{propriete}

\begin{exemplebox}[Conjugates]
Find $\bar{z}$ for each $z$:
\begin{enumerate}
\item $z=5-2i \quad\Rightarrow\quad \bar{z}=5+2i$
\item $z=-3 \quad\Rightarrow\quad \bar{z}=-3$ (real numbers are self-conjugate)
\item $z=4i \quad\Rightarrow\quad \bar{z}=-4i$ (purely imaginary numbers change sign)
\end{enumerate}
\end{exemplebox}

\courssub{Division of Complex Numbers}

Division is the only operation that does not immediately yield the standard form $a+ib$. The trick — standard in GCE examinations — is to multiply numerator and denominator by the \emph{conjugate of the denominator}, turning the denominator into a real number.

\begin{methode}[Division Algorithm: $\dfrac{z_1}{z_2}$ in Standard Form]
Let $z_1=a+ib$, $z_2=c+id$ ($z_2\neq 0$).
\begin{enumerate}
\item Write the quotient as a fraction: $\dfrac{a+ib}{c+id}$.
\item Multiply numerator and denominator by $\bar{z_2}=c-id$:
      \[
      \frac{a+ib}{c+id} \cdot \frac{c-id}{c-id}.
      \]
\item Expand the numerator using FOIL; the denominator becomes $c^2+d^2$ (a real number) because $(c+id)(c-id)=c^2+d^2$.
\item Separate the real and imaginary parts:
      \[
      \frac{(ac+bd) + i(bc-ad)}{c^2+d^2}
      = \frac{ac+bd}{c^2+d^2} + i\,\frac{bc-ad}{c^2+d^2}.
      \]
\end{enumerate}
\end{methode}

\begin{popfigure}
\begin{tikzpicture}[node distance=1.8cm, every node/.style={font=\small}]
% Nodes
\node[draw, rounded corners, fill=popblueL] (start) {Quotient $\frac{a+ib}{c+id}$};
\node[draw, rounded corners, fill=popgreenL, below of=start] (identify) {Identify denominator $c+id$};
\node[draw, rounded corners, fill=poporange, below of=identify] (conjugate) {Write its conjugate $c-id$};
\node[draw, rounded corners, fill=poppink, below of=conjugate] (multiply) {Multiply top \& bottom by $c-id$};
\node[draw, rounded corners, fill=popgold, below of=multiply] (denom) {Denominator $\to c^2+d^2 \in \mathbb{R}$};
\node[draw, rounded corners, fill=poppurple, below of=denom] (expand) {Expand numerator: $(a+ib)(c-id)$};
\node[draw, rounded corners, fill=popblueL, below of=expand] (separate) {Separate: $\frac{ac+bd}{c^2+d^2} + i\frac{bc-ad}{c^2+d^2}$};
\node[draw, rounded corners, fill=popgreenL, below of=separate] (end) {Result in standard form $x+iy$};
% Arrows
\draw[->, thick, popdark] (start) -- (identify);
\draw[->, thick, popdark] (identify) -- (conjugate);
\draw[->, thick, popdark] (conjugate) -- (multiply);
\draw[->, thick, popdark] (multiply) -- (denom);
\draw[->, thick, popdark] (denom) -- (expand);
\draw[->, thick, popdark] (expand) -- (separate);
\draw[->, thick, popdark] (separate) -- (end);
\end{tikzpicture}
\legende{Flowchart for dividing complex numbers.}
\end{popfigure}

\begin{exemplebox}[Division]
Express $\dfrac{3+4i}{1-2i}$ in the form $a+ib$.
\begin{align*}
\frac{3+4i}{1-2i} &= \frac{3+4i}{1-2i} \cdot \frac{1+2i}{1+2i}\\
&= \frac{(3+4i)(1+2i)}{1^2+2^2}\\
&= \frac{3+6i+4i+8i^2}{5}\\
&= \frac{3+10i-8}{5}\\
&= \frac{-5+10i}{5}\\
&= -1+2i.
\end{align*}
\end{exemplebox}

\begin{exoresolu}[GCE-Style Division]
Simplify $\dfrac{(2-i)^2}{3+i}$ and express the result in the form $a+ib$.
\tcblower
\textbf{Solution.}
First compute the numerator:
\[
(2-i)^2 = 4 -4i + i^2 = 3-4i.
\]
Now divide:
\[
\frac{3-4i}{3+i} = \frac{3-4i}{3+i} \cdot \frac{3-i}{3-i}
= \frac{(3-4i)(3-i)}{9+1}
= \frac{9-3i-12i+4i^2}{10}
= \frac{9-15i-4}{10}
= \frac{5-15i}{10}
= \frac{1}{2} - \frac{3}{2}i.
\]
\textbf{Answer:} $\dfrac{1}{2} - \dfrac{3}{2}i$.
\end{exoresolu}

\courssub{Properties of Complex Conjugates}

The conjugate operation interacts beautifully with all algebraic operations. These properties are frequently tested in multiple-choice and short-answer GCE questions.

\begin{propriete}[Algebraic Properties of the Conjugate]
For all $z_1,z_2\in\mathbb{C}$:
\begin{enumerate}
\item $\overline{z_1\pm z_2} = \bar{z_1}\pm\bar{z_2}$
\item $\overline{z_1z_2} = \bar{z_1}\bar{z_2}$
\item $\overline{\left(\dfrac{z_1}{z_2}\right)} = \dfrac{\bar{z_1}}{\bar{z_2}}\quad (z_2\neq 0)$
\item $\overline{\bar{z}} = z$ (involution)
\item $z+\bar{z} = 2\mathrm{Re}(z)$ \quad and \quad $z-\bar{z} = 2i\,\mathrm{Im}(z)$
\item $z\in\mathbb{R} \iff z=\bar{z}$
\item $z$ is purely imaginary $\iff z=-\bar{z}$
\end{enumerate}
\emph{Proofs of (1)--(3) are straightforward using Cartesian forms and are examinable.}
\end{propriete}

\begin{exemplebox}[Using Conjugate Properties]
Let $z_1=2+3i$, $z_2=1-i$. Verify $\overline{z_1z_2}=\bar{z_1}\bar{z_2}$.
\begin{align*}
z_1z_2 &= (2+3i)(1-i) = 2-2i+3i-3i^2 = 5+i\\
\overline{z_1z_2} &= 5-i.\\
\bar{z_1}\bar{z_2} &= (2-3i)(1+i) = 2+2i-3i-3i^2 = 5-i. \quad\checkmark
\end{align*}
\end{exemplebox}

\begin{aretenir}[Quick Extraction of Real and Imaginary Parts]
From $z+\bar{z}=2\mathrm{Re}(z)$ and $z-\bar{z}=2i\mathrm{Im}(z)$, we obtain
\[
\mathrm{Re}(z) = \frac{z+\bar{z}}{2},\qquad \mathrm{Im}(z) = \frac{z-\bar{z}}{2i}.
\]
These formulas are extremely useful when $z$ is given in polar or exponential form (Sections 5 and 6).
\end{aretenir}

\courssub{Simplification of Expressions (GCE-Style Drill)}

GCE papers often ask you to simplify expressions that mix products, quotients, powers, and conjugates. The golden rule: \textbf{always convert everything to standard form $a+ib$ at the end}.

\begin{exoresolu}[Mixed Expression]
Simplify $\dfrac{\overline{(1+2i)(3-i)}}{(2+i)^2}$ and write the result as $a+ib$.
\tcblower
\textbf{Step 1: Simplify the numerator using conjugate properties.}
\[
\overline{(1+2i)(3-i)} = \overline{1+2i}\;\overline{3-i} = (1-2i)(3+i).
\]
Compute:
\[
(1-2i)(3+i) = 3+i-6i-2i^2 = 3-5i+2 = 5-5i.
\]

\textbf{Step 2: Simplify the denominator.}
\[
(2+i)^2 = 4+4i+i^2 = 3+4i.
\]

\textbf{Step 3: Divide.}
\[
\frac{5-5i}{3+4i} = \frac{5-5i}{3+4i}\cdot\frac{3-4i}{3-4i}
= \frac{15-20i-15i+20i^2}{9+16}
= \frac{15-35i-20}{25}
= \frac{-5-35i}{25}
= -\frac{1}{5} - \frac{7}{5}i.
\]

\textbf{Answer:} $-\dfrac{1}{5} - \dfrac{7}{5}i$.
\end{exoresolu}

\begin{attention}[Common Mistakes in Mixed Expressions]
\begin{itemize}
\item \textbf{Conjugating only part of a product:} $\overline{z_1z_2} \neq z_1\bar{z_2}$. You must conjugate \emph{each factor}: $\overline{z_1z_2}=\bar{z_1}\bar{z_2}$.
\item \textbf{Forgetting to rationalise the final denominator:} Even if you used conjugate properties correctly, the final answer \emph{must} be in the form $a+ib$ with $a,b\in\mathbb{R}$.
\item \textbf{Sign errors when expanding $(c+id)(c-id)$:} Remember $(c+id)(c-id)=c^2+d^2$, \emph{not} $c^2-d^2$.
\end{itemize}
\end{attention}

\courssub{Consolidation Exercises}

\begin{exercice}[Exercise 2.1]
Express each of the following in the form $a+ib$ ($a,b\in\mathbb{R}$):
\begin{enumerate}
\item $(3-2i) + (-1+5i)$
\item $(4+i)(2-3i)$
\item $\dfrac{2+3i}{1-i}$
\item $\dfrac{(1+i)^2}{2-i}$
\item $\overline{(2-i)(3+4i)}$
\item $\dfrac{\overline{3-4i}}{1+2i}$
\item $\dfrac{1}{2+3i} + \dfrac{1}{2-3i}$
\item $(1+i)^4$ \emph{(Hint: square twice)}
\end{enumerate}
\end{exercice}

\begin{corrige}[Exercise 2.1]
\begin{enumerate}
\item $2+3i$
\item $11-10i$
\item $-\frac{1}{2}+\frac{5}{2}i$
\item $-\frac{2}{5}+\frac{4}{5}i$
\item $10-5i$
\item $\frac{11}{5}-\frac{2}{5}i$
\item $\frac{4}{13}$ (real)
\item $-4$
\end{enumerate}
\end{corrige}

\begin{exercice}[Exercise 2.2 — GCE Style]
Given that $z = x+iy$ satisfies $\dfrac{z}{z^*} = 3+4i$, where $z^*$ denotes the conjugate of $z$, find the ratio $x:y$.
\end{exercice}

\begin{corrige}[Exercise 2.2]
Let $z=x+iy$, then $z^*=x-iy$.
\[
\frac{x+iy}{x-iy} = 3+4i
\implies x+iy = (3+4i)(x-iy) = 3x-3iy+4ix-4i^2y = (3x+4y) + i(4x-3y).
\]
Equate real and imaginary parts:
\begin{align*}
x &= 3x+4y \implies -2x = 4y \implies x = -2y,\\
y &= 4x-3y \implies 4y = 4x \implies x = y.
\end{align*}
The only solution to both is $x=y=0$, but $z\neq 0$ for the quotient to be defined. \textbf{Re-check:} The second equation gives $y=4x-3y \Rightarrow 4y=4x \Rightarrow x=y$. Combined with $x=-2y$, we get $y=-2y \Rightarrow 3y=0 \Rightarrow y=0$, contradiction. \textbf{Wait — there is no non-zero $z$ satisfying this equation!} Let's verify: $\left|\frac{z}{z^*}\right| = \frac{|z|}{|z^*|} = 1$, but $|3+4i|=5\neq 1$. Hence \textbf{no such complex number exists}. This is a classic GCE trap: always check the modulus first!
\end{corrige}

\begin{savaistu}[Historical Note]
The symbol $\bar{z}$ for the conjugate was introduced by Cauchy in 1821. Earlier, Gauss used the notation $z'$ and called it the \emph{norm} of $z$ (what we now call the modulus squared). The geometric interpretation as reflection in the real axis is due to Argand (1806) and Wessel (1799), whose work was largely ignored until the 1830s. In Cameroon's GCE examinations, both $\bar{z}$ and $z^*$ are accepted notations for the conjugate — be comfortable with both.
\end{savaistu}

\begin{aretenir}[Summary of Section 2]
\begin{itemize}
\item Addition/Subtraction: combine real and imaginary parts separately.
\item Multiplication: use FOIL and $i^2=-1$; $(a+ib)(c+id)=(ac-bd)+i(ad+bc)$.
\item Conjugate: $\overline{a+ib}=a-ib$; geometric reflection in the real axis.
\item Key identity: $z\bar{z}=a^2+b^2=|z|^2 \in \mathbb{R}_{\ge 0}$.
\item Division: multiply numerator and denominator by the conjugate of the denominator.
\item Conjugate properties: linear over $+,-,\times,\div$; $\bar{\bar{z}}=z$; $z+\bar{z}=2\mathrm{Re}(z)$, $z-\bar{z}=2i\mathrm{Im}(z)$.
\item Always give final answers in standard form $a+ib$.
\end{itemize}
\end{aretenir}

\coursec{Equations and Roots in \texorpdfstring{$\mathbb{C}$}{C}}

In the previous section we learned to add, multiply and divide complex numbers, and we discovered the \cle{complex conjugate} $\bar{z}$, which behaves beautifully under all four operations. We now put this algebra to work. The whole \emph{reason} complex numbers were invented was to solve equations: the equation $x^{2}+1=0$ has no solution in $\mathbb{R}$, but it has two solutions in $\mathbb{C}$. In this section we learn to solve linear and quadratic equations in $\mathbb{C}$, and we prove one of the most elegant theorems of the GCE syllabus: for polynomials with \emph{real} coefficients, complex roots always come in \cle{conjugate pairs}.

\courssub{The Idea of a Root of a Complex Number (Lesson 29)}

\begin{definition}[Root of a complex number]
Let $w$ be a complex number and $n$ a positive integer. An \cle{$n$th root} of $w$ is any complex number $z$ such that
\[ z^{n}=w. \]
In particular, a \emph{square root} of $w$ is a solution of $z^{2}=w$.
\end{definition}

In $\mathbb{R}$, a negative number has no square root. In $\mathbb{C}$ the situation is transformed: \textbf{every} non-zero complex number has exactly two square roots, three cube roots, and in general $n$ distinct $n$th roots. For example, the equation $z^{2}=-9$ has the two solutions $z=3i$ and $z=-3i$, since $(\pm 3i)^{2}=9i^{2}=-9$.

\begin{exemplebox}[First encounter with square roots in $\mathbb{C}$]
Solve $z^{2}=-16$. We seek $z$ with $z^{2}=-16=16i^{2}$, so $z^{2}-(4i)^{2}=0$, i.e. $(z-4i)(z+4i)=0$. Hence $z=4i$ or $z=-4i$. Notice the two roots are \emph{opposites} of each other.
\end{exemplebox}

Finding the square roots of a \emph{general} complex number such as $3+4i$ requires a method (solving for the real and imaginary parts, or using the polar form); we develop it fully in the lessons on the square root of a complex number and on $n$th roots. For now, keep the concept in mind: \emph{solving $z^{2}=w$ is always possible in $\mathbb{C}$}.

\courssub{Linear Equations with Complex Coefficients (Lesson 30)}

A linear equation $az+b=0$ with $a,b\in\mathbb{C}$, $a\neq 0$, is solved exactly as in $\mathbb{R}$: isolate $z$ and divide, using the conjugate technique of Lesson 27.

\begin{exoresolu}[A linear equation in $\mathbb{C}$]
Solve $(2+i)z+3-i=0$, giving the answer in the form $x+iy$.
\tcblower
\textbf{Solution.} $z=\dfrac{-(3-i)}{2+i}=\dfrac{-3+i}{2+i}$. Multiply numerator and denominator by the conjugate of the denominator:
\[ z=\frac{(-3+i)(2-i)}{(2+i)(2-i)}=\frac{-6+3i+2i-i^{2}}{4+1}=\frac{-5+5i}{5}=-1+i. \]
\emph{Check:} $(2+i)(-1+i)+3-i=(-2+2i-i+i^{2})+3-i=(-3+i)+3-i=0$. \checkmark
\end{exoresolu}

The key point: in $\mathbb{C}$ a linear equation $az+b=0$ ($a\neq 0$) \emph{always} has exactly one solution, $z=-b/a$, because division by any non-zero complex number is always possible.

\courssub{Quadratic Equations with Real Coefficients}

Consider $az^{2}+bz+c=0$ with $a,b,c\in\mathbb{R}$, $a\neq 0$, and discriminant $\Delta=b^{2}-4ac$. Completing the square (valid in $\mathbb{C}$, since only the four operations are used) gives
\[ a\left(z+\frac{b}{2a}\right)^{2}=\frac{b^{2}-4ac}{4a}=\frac{\Delta}{4a}
\qquad\Longleftrightarrow\qquad \left(z+\frac{b}{2a}\right)^{2}=\frac{\Delta}{4a^{2}}. \]
Everything now depends on the sign of $\Delta$.

\begin{methode}[Solving $az^{2}+bz+c=0$ with real coefficients]
\begin{enumerate}
\item Compute $\Delta=b^{2}-4ac$.
\item If $\Delta>0$: two distinct real roots $z=\dfrac{-b\pm\sqrt{\Delta}}{2a}$.
\item If $\Delta=0$: one (double) real root $z=-\dfrac{b}{2a}$.
\item If $\Delta<0$: write $\sqrt{\Delta}=i\sqrt{|\Delta|}$ and apply the usual formula:
\[ z=\frac{-b\pm i\sqrt{|\Delta|}}{2a}, \]
giving \textbf{two conjugate complex roots}.
\end{enumerate}
\end{methode}

\begin{attention}[A negative discriminant is not the end of the story]
In Lower Sixth you learned that $\Delta<0$ means ``no real solution''. In $\mathbb{C}$ this becomes: \emph{two conjugate complex solutions}. Never write ``no solution'' at Advanced Level — write ``no \textbf{real} solution; the complex solutions are \dots''.
\end{attention}

\begin{exoresolu}[Solving $z^{2}-2z+5=0$]
Solve $z^{2}-2z+5=0$ in $\mathbb{C}$.
\tcblower
\textbf{Solution.} $\Delta=(-2)^{2}-4(1)(5)=4-20=-16<0$. Hence $\sqrt{\Delta}=i\sqrt{16}=4i$ and
\[ z=\frac{2\pm 4i}{2}=1\pm 2i. \]
The roots are $z_{1}=1+2i$ and $z_{2}=1-2i=\overline{z_{1}}$: a conjugate pair, as promised. \emph{Check:} $(1+2i)^{2}-2(1+2i)+5=(1+4i-4)-2-4i+5=0$. \checkmark
\end{exoresolu}

\begin{popfigure}
\begin{center}
\begin{tikzpicture}[scale=1.5]
\draw[->,gray] (-1.2,0) -- (2.8,0) node[right] {\small Re};
\draw[->,gray] (0,-2.8) -- (0,2.8) node[above] {\small Im};
\draw[dashed,popmuted] (1,2) -- (1,0) -- (1,-2);
\draw[dashed,popmuted] (1,2) -- (0,2) (1,-2) -- (0,-2);
\fill[popblue] (1,2) circle (1.8pt) node[above right,popdark] {\small $z_1=1+2i$};
\fill[poporange] (1,-2) circle (1.8pt) node[below right,popdark] {\small $z_2=1-2i$};
\draw[<->,poppurple,thick] (1.35,1.7) -- (1.35,-1.7) node[midway,right,poppurple] {\small conjugates};
\foreach \x in {1,2} \draw (\x,0.06) -- (\x,-0.06) node[below,gray] {\scriptsize $\x$};
\foreach \y in {-2,2} \draw (0.06,\y) -- (-0.06,\y) node[left,gray] {\scriptsize $\y$};
\end{tikzpicture}
\end{center}
\legende{The roots of $z^{2}-2z+5=0$ plotted on an Argand diagram: they are reflections of each other in the real axis — the geometric meaning of conjugation.}
\end{popfigure}

\courssub{The Conjugate Root Theorem}

Why did the roots above come out as conjugates? It is not a coincidence.

\begin{propriete}[Conjugate Root Theorem — proof examinable]
Let $P(z)=a_{n}z^{n}+a_{n-1}z^{n-1}+\dots+a_{1}z+a_{0}$ be a polynomial with \textbf{real} coefficients. If $\alpha\in\mathbb{C}$ satisfies $P(\alpha)=0$, then $P(\bar{\alpha})=0$. In words: \emph{complex roots of a real polynomial occur in conjugate pairs}.
\end{propriete}

\textbf{Proof.} Suppose $P(\alpha)=0$, i.e. $a_{n}\alpha^{n}+\dots+a_{1}\alpha+a_{0}=0$. Take the conjugate of both sides and use the properties of Lesson 28 ($\overline{u+v}=\bar{u}+\bar{v}$, $\overline{uv}=\bar{u}\,\bar{v}$, hence $\overline{\alpha^{k}}=\bar{\alpha}^{k}$):
\[ \overline{a_{n}\alpha^{n}+\dots+a_{1}\alpha+a_{0}}=\bar{0}
\;\Longrightarrow\; \overline{a_{n}}\,\bar{\alpha}^{n}+\dots+\overline{a_{1}}\,\bar{\alpha}+\overline{a_{0}}=0. \]
Since each coefficient $a_{k}$ is \emph{real}, $\overline{a_{k}}=a_{k}$, so
\[ a_{n}\bar{\alpha}^{n}+a_{n-1}\bar{\alpha}^{n-1}+\dots+a_{1}\bar{\alpha}+a_{0}=0, \]
which says exactly $P(\bar{\alpha})=0$. $\blacksquare$

\begin{attention}[The hypothesis ``real coefficients'' is essential]
When the coefficients are complex, the theorem fails: the roots of a quadratic with complex coefficients are generally \textbf{not} conjugates of each other. For example $z^{2}-(3+i)z+(2+2i)=0$ has roots $2$ and $1+i$ (check!), which are not conjugates.
\end{attention}

\begin{savaistu}[The Fundamental Theorem of Algebra]
The Conjugate Root Theorem is a first taste of a deeper result, the \emph{Fundamental Theorem of Algebra} (first proved rigorously by Gauss around 1799): every polynomial of degree $n$ with complex coefficients has exactly $n$ roots in $\mathbb{C}$, counted with multiplicity. This is why $\mathbb{C}$ is called \emph{algebraically closed} — in $\mathbb{C}$, no polynomial equation is ever ``impossible''.
\end{savaistu}

\courssub{Quadratic Equations with Complex Coefficients}

The quadratic formula $z=\dfrac{-b\pm\sqrt{\Delta}}{2a}$ remains valid when $a,b,c\in\mathbb{C}$ — completing the square used only field operations. The new difficulty is computing $\sqrt{\Delta}$ when $\Delta$ is itself complex. For GCE purposes, $\sqrt{\Delta}$ is found \emph{by inspection} (recognising a perfect square) or by the square-root method of Lesson 32.

\begin{exoresolu}[A quadratic with complex coefficients]
Solve $z^{2}-(3+i)z+(2+2i)=0$.
\tcblower
\textbf{Solution.} Compute the discriminant, expanding $(3+i)^{2}$ in full:
\[ \Delta=(3+i)^{2}-4(2+2i)=(9+6i+i^{2})-8-8i=(8+6i)-8-8i=-2i. \]
We need $\sqrt{-2i}$. By inspection, $(1-i)^{2}=1-2i+i^{2}=-2i$ \checkmark, so $\sqrt{\Delta}=\pm(1-i)$. Then
\[ z=\frac{(3+i)\pm(1-i)}{2} \;\Longrightarrow\; z_{1}=\frac{3+i+1-i}{2}=\frac{4}{2}=2,\qquad z_{2}=\frac{3+i-1+i}{2}=\frac{2+2i}{2}=1+i. \]
\emph{Check:} sum of roots $=2+(1+i)=3+i=-b/a$ \checkmark; product $=2(1+i)=2+2i=c/a$ \checkmark. Note the roots are \emph{not} conjugates — the coefficients are not real, so the Conjugate Root Theorem does not apply.
\end{exoresolu}

\begin{attention}[Careful arithmetic with $\Delta$]
Most errors in this topic come from expanding $(p+qi)^{2}$ incorrectly. Remember $(p+qi)^{2}=p^{2}-q^{2}+2pqi$: the cross term carries the $i$, and $i^{2}=-1$ flips the sign of $q^{2}$. Write the expansion in full before simplifying, and always verify your candidate square root by squaring it again.
\end{attention}

\courssub{Factorisation, Sum and Product of Roots}

\begin{propriete}[Factorisation and Vieta's relations in $\mathbb{C}$]
Every quadratic $az^{2}+bz+c$ ($a\neq 0$, coefficients real or complex) splits over $\mathbb{C}$ as
\[ az^{2}+bz+c=a(z-z_{1})(z-z_{2}), \]
where $z_{1},z_{2}$ are its roots. Expanding and comparing coefficients gives the \cle{sum and product of roots}:
\[ z_{1}+z_{2}=-\frac{b}{a},\qquad z_{1}z_{2}=\frac{c}{a}. \]
These relations remain valid in $\mathbb{C}$.
\end{propriete}

Indeed $a(z-z_{1})(z-z_{2})=az^{2}-a(z_{1}+z_{2})z+az_{1}z_{2}$, and matching with $az^{2}+bz+c$ yields both formulas. For a real quadratic with $\Delta<0$, the conjugate pair $z_{1}=p+qi$, $z_{2}=p-qi$ gives the nice real values $z_{1}+z_{2}=2p$ and $z_{1}z_{2}=p^{2}+q^{2}=|z_{1}|^{2}$.

\begin{exemplebox}[Using sum and product]
The roots of $z^{2}-2z+5=0$ are $1\pm 2i$: sum $=2=-(-2)/1$ \checkmark, product $=(1+2i)(1-2i)=1+4=5$ \checkmark. Conversely, a quadratic with roots $3i$ and $-3i$ is $z^{2}-(0)z+(3i)(-3i)=z^{2}+9=0$.
\end{exemplebox}

\begin{methode}[``Show that $\alpha$ is a root, hence factorise completely'' (GCE favourite)]
Given a real polynomial $P$ and a complex number $\alpha$:
\begin{enumerate}
\item \textbf{Verify} $P(\alpha)=0$ by direct substitution (compute powers of $\alpha$ carefully).
\item \textbf{Conjugate Root Theorem:} since the coefficients are real, $\bar{\alpha}$ is also a root.
\item \textbf{Build the quadratic factor} $(z-\alpha)(z-\bar{\alpha})=z^{2}-2\operatorname{Re}(\alpha)\,z+|\alpha|^{2}$, which has \emph{real} coefficients.
\item \textbf{Divide} $P$ by this quadratic (long division or inspection) to obtain the remaining factor, then solve it.
\end{enumerate}
\end{methode}

\begin{exoresolu}[Exam-style: factorising a cubic over $\mathbb{C}$]
Let $P(z)=z^{3}-z^{2}+3z+5$. Show that $1+2i$ is a root of $P$, and hence factorise $P$ completely over $\mathbb{C}$.
\tcblower
\textbf{Solution.} \emph{Step 1:} compute powers: $(1+2i)^{2}=1+4i-4=-3+4i$; $(1+2i)^{3}=(1+2i)(-3+4i)=-3+4i-6i+8i^{2}=-11-2i$. Then
\[ P(1+2i)=(-11-2i)-(-3+4i)+3(1+2i)+5=-11-2i+3-4i+3+6i+5=0. \checkmark \]
\emph{Step 2:} coefficients are real, so $1-2i$ is also a root.

\emph{Step 3:} $(z-1-2i)(z-1+2i)=(z-1)^{2}+4=z^{2}-2z+5$.

\emph{Step 4:} divide: $z^{3}-z^{2}+3z+5=(z^{2}-2z+5)(z+1)$ (check: $(z^{2}-2z+5)(z+1)=z^{3}-z^{2}+3z+5$ \checkmark). Hence
\[ P(z)=(z+1)\bigl(z-(1+2i)\bigr)\bigl(z-(1-2i)\bigr), \]
with roots $-1$, $1+2i$, $1-2i$.
\end{exoresolu}

\begin{popfigure}
\begin{center}
\begin{tikzpicture}[
  node distance=0.9cm,
  every node/.style={align=center},
  start/.style={rectangle, rounded corners, draw=popdark, fill=popblueL, text width=4.6cm, inner sep=5pt},
  dec/.style={diamond, draw=popdark, fill=popgold!30, aspect=2.2, inner sep=1pt, text width=2.6cm},
  res/.style={rectangle, rounded corners, draw=popdark, fill=popgreenL, text width=3.6cm, inner sep=5pt},
  arr/.style={->, thick, popdark}]
\node[start, align=center] (s){\textbf{Solve } $az^{2}+bz+c=0$\\ compute $\Delta=b^{2}-4ac$};
\node[dec, below=of s] (d) {Sign of $\Delta$?};
\node[res, below left=1.1cm and -1.2cm of d, align=center] (l){$\Delta>0$:\\ two real roots\\ $\dfrac{-b\pm\sqrt{\Delta}}{2a}$};
\node[res, below=1.6cm of d, align=center] (m){$\Delta=0$:\\ double root\\ $z=-\dfrac{b}{2a}$};
\node[res, below right=1.1cm and -1.2cm of d, align=center] (r){$\Delta<0$:\\ conjugate pair\\ $\dfrac{-b\pm i\sqrt{|\Delta|}}{2a}$};
\draw[arr] (s) -- (d);
\draw[arr] (d) -| (l);
\draw[arr] (d) -- (m);
\draw[arr] (d) -| (r);
\end{tikzpicture}
\end{center}
\legende{Decision tree for solving $az^{2}+bz+c=0$ with real coefficients, according to the sign of the discriminant.}
\end{popfigure}

\courssub{Consolidation Exercises}

\begin{exercice}[Quadratics in $\mathbb{C}$]
Solve in $\mathbb{C}$: \textbf{(a)} $z^{2}+4z+13=0$; \textbf{(b)} $2z^{2}-2z+1=0$; \textbf{(c)} $z^{2}-(5+i)z+(8+i)=0$ (find $\sqrt{\Delta}$ by inspection).
\end{exercice}

\begin{corrige}[Exercise 1]
\textbf{(a)} $\Delta=16-52=-36$, $\sqrt{\Delta}=6i$: $z=\dfrac{-4\pm 6i}{2}=-2\pm 3i$.
\textbf{(b)} $\Delta=4-8=-4$, $\sqrt{\Delta}=2i$: $z=\dfrac{2\pm 2i}{4}=\dfrac{1\pm i}{2}$.
\textbf{(c)} $\Delta=(5+i)^{2}-4(8+i)=(24+10i)-(32+4i)=-8+6i$. Inspection: $(1+3i)^{2}=1+6i-9=-8+6i$ \checkmark. Hence $z=\dfrac{(5+i)\pm(1+3i)}{2}$, giving $z_{1}=3+2i$, $z_{2}=2-i$.
\end{corrige}

\begin{exercice}[Factorising over $\mathbb{C}$]
Show that $2+i$ is a root of $P(z)=z^{3}-3z^{2}+z+5$, and factorise $P$ completely over $\mathbb{C}$.
\end{exercice}

\begin{corrige}[Exercise 2]
$(2+i)^{2}=3+4i$; $(2+i)^{3}=(2+i)(3+4i)=2+11i$. Then $P(2+i)=(2+11i)-3(3+4i)+(2+i)+5=2+11i-9-12i+2+i+5=0$ \checkmark. Real coefficients $\Rightarrow$ $2-i$ is a root. Quadratic factor: $(z-2-i)(z-2+i)=(z-2)^{2}+1=z^{2}-4z+5$. Division: $P(z)=(z^{2}-4z+5)(z+1)$. Complete factorisation: $P(z)=(z+1)\bigl(z-(2+i)\bigr)\bigl(z-(2-i)\bigr)$.
\end{corrige}

\begin{aretenir}[Key points of this section]
\begin{itemize}
\item Every non-zero complex number has square roots in $\mathbb{C}$: solving $z^{2}=w$ is always possible.
\item A linear equation $az+b=0$ ($a\neq 0$) has the unique solution $z=-b/a$; divide using the conjugate.
\item Real quadratic with $\Delta<0$: two \textbf{conjugate} roots $z=\dfrac{-b\pm i\sqrt{|\Delta|}}{2a}$ — never ``no solution''.
\item \textbf{Conjugate Root Theorem:} real-coefficient polynomials have their non-real roots in conjugate pairs (proof examinable).
\item Complex coefficients: the quadratic formula still works, but $\sqrt{\Delta}$ must be found in $\mathbb{C}$, and the roots are generally \emph{not} conjugates.
\item $az^{2}+bz+c=a(z-z_{1})(z-z_{2})$ with $z_{1}+z_{2}=-b/a$, $z_{1}z_{2}=c/a$ — valid in $\mathbb{C}$.
\item GCE technique: verify $P(\alpha)=0$, deduce $\bar{\alpha}$, form the real quadratic factor, divide, finish.
\end{itemize}
\end{aretenir}

\coursec{The Argand Diagram, Modulus and Argument}

In the previous section we learned how to solve quadratic equations in $\mathbb{C}$ and how to manipulate complex numbers algebraically. We now turn to a geometric viewpoint that will transform our understanding of complex numbers. By representing a complex number as a point or a vector in a plane, concepts like modulus and argument acquire a vivid visual meaning, and operations like multiplication and division become simple geometric transformations. This geometric representation, known as the \cle{Argand diagram}, is the gateway to the polar and exponential forms that will simplify powers and roots enormously.

\courssub{The Argand Diagram: Geometric Representation}

A complex number $z = a + ib$ is completely determined by the ordered pair of real numbers $(a, b)$. It is therefore natural to associate $z$ with a point $M$ in a Cartesian coordinate system.

\begin{definition}[The Argand Diagram]
The \cle{Argand diagram} (or complex plane) is a Cartesian plane where:
\begin{itemize}
    \item The horizontal axis ($x$-axis) is called the \textbf{real axis} and represents the real part $\operatorname{Re}(z) = a$.
    \item The vertical axis ($y$-axis) is called the \textbf{imaginary axis} and represents the imaginary part $\operatorname{Im}(z) = b$.
\end{itemize}
The complex number $z = a + ib$ is represented by the point $M(a, b)$ or by the position vector $\overrightarrow{OM}$.
\end{definition}

\begin{popfigure}
\begin{tikzpicture}[scale=1.2, >=stealth]
    % Axes
    \draw[->, thick, popdark] (-3.5,0) -- (3.5,0) node[right] {$\operatorname{Re}(z)$};
    \draw[->, thick, popdark] (0,-3) -- (0,3) node[above] {$\operatorname{Im}(z)$};
    % Ticks
    \foreach \x in {-3,-2,-1,1,2,3} \draw (\x,0.05) -- (\x,-0.05) node[below, font=\small] {\x};
    \foreach \y in {-2,-1,1,2} \draw (0.05,\y) -- (-0.05,\y) node[left, font=\small] {\y};
    % Point M
    \coordinate (O) at (0,0);
    \coordinate (M) at (2.5,1.5);
    \draw[popblue, thick, ->] (O) -- (M) node[midway, above left, popblue] {$\overrightarrow{OM}$};
    \fill[popblue] (M) circle (2.5pt) node[above right] {$M(a,b)$};
    % Projections
    \draw[dashed, popmuted] (M) -- (2.5,0) node[below, font=\small] {$a$};
    \draw[dashed, popmuted] (M) -- (0,1.5) node[left, font=\small] {$b$};
    % Modulus and argument labels
    \draw[poporange, thick] (O) -- (M) node[midway, below right, poporange] {$|z|$};
    \draw[poppurple, thick, ->] (0.5,0) arc (0:30.96:0.5) node[midway, right, poppurple] {$\theta$};
\end{tikzpicture}
\legende{The Argand diagram: $z = a + ib$ corresponds to point $M(a,b)$ or vector $\overrightarrow{OM}$.}
\end{popfigure}

\begin{aretenir}[Key Geometric Correspondences]
\begin{itemize}
    \item $z = 0$ corresponds to the origin $O$.
    \item Real numbers ($b=0$) lie on the real axis.
    \item Pure imaginary numbers ($a=0$) lie on the imaginary axis.
    \item The complex conjugate $\bar{z} = a - ib$ is the reflection of $M$ across the real axis.
    \item The opposite $-z = -a - ib$ is the point symmetric to $M$ with respect to the origin (central symmetry).
\end{itemize}
\end{aretenir}

\begin{savaistu}[Historical Note]
The geometric representation of complex numbers was independently proposed by the Norwegian surveyor Caspar Wessel (1797), the Swiss mathematician Jean-Robert Argand (1806), and the French mathematician Adrien-Quentin Buée (1806). It was later popularised by Gauss, who coined the term ``complex number'' and established the notation $i$ for $\sqrt{-1}$. This geometric insight turned complex numbers from an algebraic curiosity into a fundamental tool for analysis and physics.
\end{savaistu}

\courssub{The Modulus of a Complex Number}

The distance from the origin $O$ to the point $M(a,b)$ is a fundamental quantity.

\begin{definition}[Modulus]
The \cle{modulus} (or absolute value) of a complex number $z = a + ib$ is the length of the vector $\overrightarrow{OM}$ or the distance $OM$. It is denoted by $|z|$ and defined by:
\[
|z| = \sqrt{a^2 + b^2}.
\]
\end{definition}

\begin{propriete}[Properties of the Modulus]
For any complex numbers $z, z_1, z_2$:
\begin{enumerate}
    \item $|z| \ge 0$, and $|z| = 0 \iff z = 0$.
    \item $|z| = |\bar{z}| = |-z|$.
    \item \textbf{Multiplicative property:} $|z_1 z_2| = |z_1|\,|z_2|$.
    \item \textbf{Quotient property:} $\left|\dfrac{z_1}{z_2}\right| = \dfrac{|z_1|}{|z_2|}$ (for $z_2 \neq 0$).
    \item \textbf{Triangle inequality:} $|z_1 + z_2| \le |z_1| + |z_2|$.
\end{enumerate}
\end{propriete}

\begin{exemplebox}[Computing Moduli]
Find the modulus of $z_1 = 3 + 4i$, $z_2 = -5i$, and $z_3 = \dfrac{1+i}{1-i}$.
\begin{align*}
|z_1| &= \sqrt{3^2 + 4^2} = \sqrt{25} = 5. \\
|z_2| &= \sqrt{0^2 + (-5)^2} = 5. \\
|z_3| &= \frac{|1+i|}{|1-i|} = \frac{\sqrt{1^2+1^2}}{\sqrt{1^2+(-1)^2}} = \frac{\sqrt{2}}{\sqrt{2}} = 1.
\end{align*}
\end{exemplebox}

\courssub{The Argument of a Complex Number}

While the modulus tells us \emph{how far} the point is from the origin, the argument tells us \emph{in which direction}.

\begin{definition}[Argument]
Let $z = a + ib \neq 0$. An \cle{argument} of $z$, denoted $\arg(z)$, is any real number $\theta$ such that:
\[
\cos\theta = \frac{a}{|z|}, \qquad \sin\theta = \frac{b}{|z|}.
\]
Geometrically, $\theta$ is the oriented angle $(\overrightarrow{Ox}, \overrightarrow{OM})$, measured in radians, positive in the anticlockwise direction.
\end{definition}

Since adding $2\pi$ to an angle does not change its cosine and sine, the argument is defined \emph{modulo $2\pi$}: if $\theta$ is an argument, then $\theta + 2k\pi$ ($k \in \mathbb{Z}$) are also arguments. To have a unique value, we define the \cle{principal argument}.

\begin{definition}[Principal Argument]
The \cle{principal argument} of $z \neq 0$, denoted $\operatorname{Arg}(z)$ (capital A), is the unique argument $\theta$ satisfying:
\[
-\pi < \theta \le \pi.
\]
\end{definition}

\begin{attention}[Important Warnings]
\begin{itemize}
    \item \textbf{Argument of zero is undefined.} Never write $\arg(0)$ or $\operatorname{Arg}(0)$; the point $O$ has no well-defined direction.
    \item \textbf{$\tan\theta = \frac{b}{a}$ is not enough.} The equation $\tan\theta = \frac{b}{a}$ has two solutions in $(-\pi, \pi]$ differing by $\pi$. You \emph{must} check the quadrant of the point $M(a,b)$ to choose the correct one.
\end{itemize}
\end{attention}

\courssub{Finding the Principal Argument: Quadrant by Quadrant}

Let $z = a + ib \neq 0$ and $\alpha = \arctan\left|\frac{b}{a}\right| \in \left[0, \frac{\pi}{2}\right)$ (the acute angle with the real axis). The principal argument $\operatorname{Arg}(z)$ depends on the signs of $a$ and $b$:

\begin{popfigure}
\begin{center}
\begin{tabular}{|c|c|c|c|}
\hline
\rowcolor{popblueL}
\textbf{Quadrant} & \textbf{Signs of $(a,b)$} & $\boldsymbol{\tan\theta = b/a}$ & \textbf{$\operatorname{Arg}(z)$} \\
\hline
I & $a>0,\ b>0$ & $>0$ & $\alpha = \arctan(b/a)$ \\
\hline
II & $a<0,\ b>0$ & $<0$ & $\pi - \alpha = \pi - \arctan(|b/a|)$ \\
\hline
III & $a<0,\ b<0$ & $>0$ & $-\pi + \alpha = -\pi + \arctan(|b/a|)$ \\
\hline
IV & $a>0,\ b<0$ & $<0$ & $-\alpha = -\arctan(|b/a|)$ \\
\hline
\end{tabular}
\end{center}
\legende{Determining the principal argument according to the quadrant of $M(a,b)$.}
\end{popfigure}

\begin{exemplebox}[Finding Principal Arguments]
Find $\operatorname{Arg}(z)$ for:
\begin{enumerate}
    \item $z_1 = 1 + i\sqrt{3}$: $a>0, b>0$ (QI). $\alpha = \arctan(\sqrt{3}/1) = \pi/3$. $\operatorname{Arg}(z_1) = \pi/3$.
    \item $z_2 = -2 + 2i$: $a<0, b>0$ (QII). $\alpha = \arctan(2/2) = \pi/4$. $\operatorname{Arg}(z_2) = \pi - \pi/4 = 3\pi/4$.
    \item $z_3 = -1 - i\sqrt{3}$: $a<0, b<0$ (QIII). $\alpha = \arctan(\sqrt{3}) = \pi/3$. $\operatorname{Arg}(z_3) = -\pi + \pi/3 = -2\pi/3$.
    \item $z_4 = 3 - 3i$: $a>0, b<0$ (QIV). $\alpha = \arctan(1) = \pi/4$. $\operatorname{Arg}(z_4) = -\pi/4$.
    \item $z_5 = -5$: Negative real axis. $\operatorname{Arg}(z_5) = \pi$.
    \item $z_6 = 4i$: Positive imaginary axis. $\operatorname{Arg}(z_6) = \pi/2$.
\end{enumerate}
\end{exemplebox}

\begin{propriete}[Properties of the Argument]
For $z, z_1, z_2 \neq 0$:
\begin{enumerate}
    \item $\arg(\bar{z}) = -\arg(z) \pmod{2\pi}$.
    \item $\arg(-z) = \arg(z) + \pi \pmod{2\pi}$.
    \item \textbf{Product:} $\arg(z_1 z_2) = \arg(z_1) + \arg(z_2) \pmod{2\pi}$.
    \item \textbf{Quotient:} $\arg\left(\dfrac{z_1}{z_2}\right) = \arg(z_1) - \arg(z_2) \pmod{2\pi}$.
    \item \textbf{Power:} $\arg(z^n) = n\arg(z) \pmod{2\pi}$.
\end{enumerate}
\textbf{Note:} For the principal argument, these equalities hold modulo $2\pi$; one may need to add or subtract $2\pi$ to bring the result back into $(-\pi, \pi]$.
\end{propriete}

\courssub{Geometric Interpretation of Operations}

The Argand diagram gives a beautiful geometric meaning to algebraic operations.

\paragraph{Addition and Subtraction: The Parallelogram Rule}
Let $z_1 = a_1 + ib_1$ and $z_2 = a_2 + ib_2$ be represented by points $M_1, M_2$ and vectors $\overrightarrow{OM_1}, \overrightarrow{OM_2}$.
\[
z_1 + z_2 = (a_1+a_2) + i(b_1+b_2) \quad \longleftrightarrow \quad \overrightarrow{OM_1} + \overrightarrow{OM_2} = \overrightarrow{OM_3}
\]
where $M_3$ is the fourth vertex of the parallelogram with vertices $O, M_1, M_3, M_2$.
\[
z_1 - z_2 = (a_1-a_2) + i(b_1-b_2) \quad \longleftrightarrow \quad \overrightarrow{OM_1} - \overrightarrow{OM_2} = \overrightarrow{M_2M_1}.
\]
Hence $|z_1 - z_2|$ is exactly the \textbf{distance} between the points $M_1$ and $M_2$.

\begin{popfigure}
\begin{tikzpicture}[scale=1.1, >=stealth]
    \draw[->, thick, popdark] (-1,0) -- (5,0) node[right] {$\operatorname{Re}$};
    \draw[->, thick, popdark] (0,-1) -- (0,4) node[above] {$\operatorname{Im}$};
    \coordinate (O) at (0,0);
    \coordinate (M1) at (3,1);
    \coordinate (M2) at (1,2.5);
    \coordinate (M3) at (4,3.5); % M1+M2
    % Vectors
    \draw[popblue, thick, ->] (O) -- (M1) node[midway, below right, popblue] {$\vec{z_1}$};
    \draw[poporange, thick, ->] (O) -- (M2) node[midway, left, poporange] {$\vec{z_2}$};
    \draw[popgreen, thick, ->] (O) -- (M3) node[midway, above right, popgreen] {$\vec{z_1}+\vec{z_2}$};
    % Parallelogram
    \draw[dashed, popmuted] (M1) -- (M3);
    \draw[dashed, popmuted] (M2) -- (M3);
    % Points
    \fill[popblue] (M1) circle (2pt) node[below right] {$M_1(z_1)$};
    \fill[poporange] (M2) circle (2pt) node[above left] {$M_2(z_2)$};
    \fill[popgreen] (M3) circle (2pt) node[above right] {$M_3(z_1+z_2)$};
    \fill[popdark] (O) circle (2pt) node[below left] {$O$};
    % Difference vector
    \draw[poppurple, thick, ->] (M2) -- (M1) node[midway, above, poppurple] {$\vec{z_1}-\vec{z_2}$};
\end{tikzpicture}
\legende{Parallelogram rule: $z_1+z_2$ is the diagonal; $z_1-z_2$ is the vector from $M_2$ to $M_1$.}
\end{popfigure}

\paragraph{Conjugation and Opposite}
\begin{itemize}
    \item \textbf{Conjugate $\bar{z}$:} Reflection of $M(a,b)$ across the real axis $\to M'(a,-b)$. Modulus unchanged, argument changes sign.
    \item \textbf{Opposite $-z$}: Central symmetry about the origin $\to M'(-a,-b)$. Modulus unchanged, argument increased by $\pi$ (or decreased by $\pi$).
\end{itemize}

\courssub{Square Roots of a Complex Number (Lesson 32)}

Every non-zero complex number has exactly two square roots, which are opposites of each other. To find them, we solve $(x+iy)^2 = a+ib$ with $x,y \in \mathbb{R}$.

\begin{methode}[Finding the Square Roots of $a+ib$]
Given $z = a+ib \neq 0$, we seek $w = x+iy$ such that $w^2 = z$.
\begin{enumerate}
    \item Expand: $(x+iy)^2 = (x^2 - y^2) + i(2xy) = a + ib$.
    \item Equate real and imaginary parts:
    \[
    \begin{cases}
    x^2 - y^2 = a & (1)\\
    2xy = b & (2)
    \end{cases}
    \]
    \item Use the modulus: $|w^2| = |z| \implies (x^2+y^2)^2 = a^2+b^2 \implies x^2+y^2 = \sqrt{a^2+b^2} = |z|$. \quad (3)
    \item Solve the system (1) and (3) for $x^2$ and $y^2$:
    \[
    x^2 = \frac{|z| + a}{2}, \qquad y^2 = \frac{|z| - a}{2}.
    \]
    \item Determine the signs of $x$ and $y$ using equation (2): \textbf{$x$ and $y$ have the same sign if $b>0$, opposite signs if $b<0$}.
    \item The two square roots are $\pm(x+iy)$.
\end{enumerate}
\end{methode}

\begin{exoresolu}[Square Roots of $3+4i$ and $7-24i$]
\textbf{Statement:} Find the square roots of $3+4i$ and $7-24i$.
\tcblower
\textbf{Detailed Solution:}

\textbf{1. Square roots of $3+4i$.}
Here $a=3, b=4$. $|z| = \sqrt{3^2+4^2} = 5$.
\[
x^2 = \frac{5+3}{2} = 4 \implies x = \pm 2, \qquad
y^2 = \frac{5-3}{2} = 1 \implies y = \pm 1.
\]
Since $b=4 > 0$, $x$ and $y$ have the \textbf{same sign}.
The two square roots are $\mathbf{2+i}$ and $\mathbf{-2-i}$.
Check: $(2+i)^2 = 4 + 4i + i^2 = 3+4i$. ✓

\textbf{2. Square roots of $7-24i$.}
Here $a=7, b=-24$. $|z| = \sqrt{7^2+(-24)^2} = \sqrt{49+576} = \sqrt{625} = 25$.
\[
x^2 = \frac{25+7}{2} = 16 \implies x = \pm 4, \qquad
y^2 = \frac{25-7}{2} = 9 \implies y = \pm 3.
\]
Since $b=-24 < 0$, $x$ and $y$ have \textbf{opposite signs}.
The two square roots are $\mathbf{4-3i}$ and $\mathbf{-4+3i}$.
Check: $(4-3i)^2 = 16 - 24i + 9i^2 = 7-24i$. ✓
\end{exoresolu}

\begin{attention}[Common Mistakes in Square Roots]
\begin{itemize}
    \item Forgetting that $x^2+y^2 = |z|$ (not $|z|^2$).
    \item Taking $x = \pm\sqrt{\dots}$ and $y = \pm\sqrt{\dots}$ independently, yielding four candidates instead of two. \textbf{Always use the sign of $b$ (equation $2xy=b$) to pair the signs correctly.}
    \item Writing $\sqrt{3+4i} = 2+i$. The radical symbol $\sqrt{\cdot}$ denotes the \emph{principal} square root (usually the one with non-negative real part, or argument in $(-\pi/2, \pi/2]$). It is safer to write: ``The square roots of $3+4i$ are $2+i$ and $-2-i$.''
\end{itemize}
\end{attention}

\courssub{Synthesis: Modulus–Argument Form}

The modulus $r = |z|$ and an argument $\theta = \arg(z)$ give the \cle{polar (trigonometric) form}:
\[
z = r(\cos\theta + i\sin\theta), \qquad r \ge 0,\ \theta \in \mathbb{R}.
\]
This form makes the geometric properties of multiplication and division obvious:
\[
z_1 z_2 = r_1 r_2 \bigl(\cos(\theta_1+\theta_2) + i\sin(\theta_1+\theta_2)\bigr),
\quad
\frac{z_1}{z_2} = \frac{r_1}{r_2} \bigl(\cos(\theta_1-\theta_2) + i\sin(\theta_1-\theta_2)\bigr).
\]
We will exploit this fully in the next section with the exponential form $z = r e^{i\theta}$ and De Moivre's theorem.

\begin{aretenir}[Summary of Key Points]
\begin{itemize}
    \item \textbf{Argand diagram:} $z = a+ib \leftrightarrow M(a,b)$ or $\overrightarrow{OM}$.
    \item \textbf{Modulus:} $|z| = \sqrt{a^2+b^2} = OM$. Properties: $|z_1z_2|=|z_1||z_2|$, $|z_1/z_2|=|z_1|/|z_2|$, triangle inequality.
    \item \textbf{Argument:} $\arg(z)$ is the oriented angle $(\overrightarrow{Ox}, \overrightarrow{OM})$. Principal argument $\operatorname{Arg}(z) \in (-\pi, \pi]$. Undefined for $z=0$.
    \item \textbf{Quadrant rule:} Use signs of $a,b$ to find $\operatorname{Arg}(z)$ from $\alpha = \arctan|b/a|$.
    \item \textbf{Geometry:} Addition = parallelogram; $|z_1-z_2| = M_1M_2$; $\bar{z}$ = reflection in real axis; $-z$ = central symmetry.
    \item \textbf{Square roots:} Solve $x^2-y^2=a,\ 2xy=b,\ x^2+y^2=|z|$. Two opposite roots $\pm(x+iy)$. Sign of $x,y$ determined by sign of $b$.
\end{itemize}
\end{aretenir}

\begin{exercice}[Practice Exercises]
\begin{enumerate}
    \item Plot the following complex numbers on an Argand diagram and find their modulus and principal argument:
    $z_1 = -2\sqrt{3} + 2i$, \quad $z_2 = -1 - i$, \quad $z_3 = 4e^{i2\pi/3}$.
    \item Given $z_1 = 2(\cos\frac{\pi}{6} + i\sin\frac{\pi}{6})$ and $z_2 = 3(\cos\frac{-\pi}{3} + i\sin\frac{-\pi}{3})$, find the modulus and principal argument of $z_1z_2$ and $z_1/z_2$ without converting to algebraic form.
    \item Find the square roots of:
    \begin{enumerate}
        \item $5 + 12i$
        \item $-8 - 6i$
        \item $-7 + 24i$
    \end{enumerate}
    \item Let $z_1 = 1 + i\sqrt{3}$ and $z_2 = \sqrt{3} - i$. Represent $z_1, z_2, z_1+z_2, z_1-z_2$ on the same Argand diagram. Verify geometrically that $|z_1+z_2|^2 + |z_1-z_2|^2 = 2(|z_1|^2+|z_2|^2)$ (parallelogram law).
    \item If $|z| = 2$ and $\arg(z) = \frac{2\pi}{3}$, find the modulus and principal argument of $iz$, $\frac{1}{z}$, and $\bar{z}$.
\end{enumerate}
\end{exercice}

\begin{corrige}[Exercise 1]
\begin{enumerate}
    \item 
    \begin{itemize}
        \item $z_1 = -2\sqrt{3} + 2i$: $|z_1| = \sqrt{12+4} = 4$. QII: $\alpha = \arctan(2/(2\sqrt{3})) = \pi/6$. $\operatorname{Arg}(z_1) = \pi - \pi/6 = 5\pi/6$.
        \item $z_2 = -1 - i$: $|z_2| = \sqrt{2}$. QIII: $\alpha = \arctan(1) = \pi/4$. $\operatorname{Arg}(z_2) = -\pi + \pi/4 = -3\pi/4$.
        \item $z_3 = 4e^{i2\pi/3}$: $|z_3| = 4$, $\operatorname{Arg}(z_3) = 2\pi/3$ (already in polar form).
    \end{itemize}
    \item $|z_1z_2| = 2\times 3 = 6$, $\arg(z_1z_2) = \pi/6 + (-\pi/3) = -\pi/6 \in (-\pi,\pi]$. $|z_1/z_2| = 2/3$, $\arg(z_1/z_2) = \pi/6 - (-\pi/3) = \pi/2$.
    \item 
    \begin{enumerate}
        \item $5+12i$: $|z|=13$. $x^2=(13+5)/2=9 \Rightarrow x=\pm3$, $y^2=(13-5)/2=4 \Rightarrow y=\pm2$. $b>0 \Rightarrow$ same sign. Roots: $\pm(3+2i)$.
        \item $-8-6i$: $|z|=10$. $x^2=(10-8)/2=1 \Rightarrow x=\pm1$, $y^2=(10+8)/2=9 \Rightarrow y=\pm3$. $b<0 \Rightarrow$ opposite signs. Roots: $\pm(1-3i)$.
        \item $-7+24i$: $|z|=25$. $x^2=(25-7)/2=9 \Rightarrow x=\pm3$, $y^2=(25+7)/2=16 \Rightarrow y=\pm4$. $b>0 \Rightarrow$ same sign. Roots: $\pm(3+4i)$.
    \end{enumerate}
    \item $z_1=1+i\sqrt{3}$ ($|z_1|=2, \arg=\pi/3$), $z_2=\sqrt{3}-i$ ($|z_2|=2, \arg=-\pi/6$). 
    $z_1+z_2 = (1+\sqrt{3}) + i(\sqrt{3}-1)$, $z_1-z_2 = (1-\sqrt{3}) + i(\sqrt{3}+1)$.
    $|z_1+z_2|^2 = (1+\sqrt{3})^2+(\sqrt{3}-1)^2 = 1+3+2\sqrt{3}+3+1-2\sqrt{3}=8$.
    $|z_1-z_2|^2 = (1-\sqrt{3})^2+(\sqrt{3}+1)^2 = 1+3-2\sqrt{3}+3+1+2\sqrt{3}=8$.
    LHS $= 8+8=16$. RHS $= 2(4+4)=16$. Verified.
    \item $|z|=2, \arg(z)=2\pi/3$.
    $iz$: $|iz|=2$, $\arg(iz)=2\pi/3+\pi/2=7\pi/6 \equiv -5\pi/6$ (principal).
    $1/z$: $|1/z|=1/2$, $\arg(1/z)=-2\pi/3$ (principal).
    $\bar{z}$: $|\bar{z}|=2$, $\arg(\bar{z})=-2\pi/3$ (principal).
\end{enumerate}
\end{corrige}

\coursec{Polar and Exponential Forms}

In the previous section we learned to \emph{locate} a complex number in the Argand diagram using two pieces of information: its distance $r = |z|$ from the origin and the angle $\theta = \arg(z)$ that the line $Oz$ makes with the positive real axis. We now exploit a beautiful idea: these two numbers $r$ and $\theta$ do not merely \emph{describe} the position of $z$ — they can be used to \emph{write} $z$ itself. This gives two new ways of writing a complex number, the \cle{polar form} and the \cle{exponential form}, which turn the hardest operations of complex algebra (multiplication, division, powers, roots) into simple arithmetic on $r$ and $\theta$.

\courssub{The Polar (Trigonometric) Form of a Complex Number}

\textbf{From geometry to algebra.}\par
Let $z = a + ib$ be a non-zero complex number, represented by the point $M(a,b)$ in the Argand diagram. Let $r = |z| = \sqrt{a^{2}+b^{2}}$ and $\theta = \arg(z)$. By the very definition of cosine and sine on the circle of radius $r$ centred at $O$, the coordinates of $M$ satisfy
\[
a = r\cos\theta \qquad \text{and} \qquad b = r\sin\theta .
\]
Substituting into $z = a + ib$ and factorising $r$:
\[
z = r\cos\theta + i\,r\sin\theta = r(\cos\theta + i\sin\theta).
\]

\begin{definition}[Polar form of a complex number]
Let $z$ be a non-zero complex number with modulus $r = |z| > 0$ and argument $\theta = \arg(z)$. The \cle{polar form} (or \cle{trigonometric form}) of $z$ is
\[
\boxed{\,z = r(\cos\theta + i\sin\theta)\,}.
\]
The number $0$ has no polar form, since its argument is undefined. By convention, unless otherwise stated, $\theta$ denotes the \cle{principal argument} $\operatorname{Arg}(z) \in {]{-\pi},\,\pi]}$.
\end{definition}

\begin{popfigure}
\begin{tikzpicture}[scale=2.2, >=stealth]
\draw[->, popmuted] (-0.5,0) -- (2.6,0) node[below left, popink]{\small Re};
\draw[->, popmuted] (0,-0.5) -- (0,2.2) node[below left, popink]{\small Im};
\draw[popline, dashed] (0,0) circle (1);
\coordinate (O) at (0,0);
\coordinate (M) at (1.5,1.299);
\draw[very thick, popblue] (O) -- (M) node[midway, above left, popdark]{$r$};
\draw[poporange, thick] (M) -- (1.5,0) node[midway, right, popdark]{\small $b = r\sin\theta$};
\draw[poporange, thick] (O) -- (1.5,0) node[midway, below, popdark]{\small $a = r\cos\theta$};
\fill[popdark] (M) circle (0.025) node[above right, popdark]{$M:\ z$};
\draw[->, popgreen, thick] (0.55,0) arc (0:40.9:0.55) node[midway, right, popgreen]{$\theta$};
\node[popdark, align=center] at (1.9,1.95) {\small $z = a + ib$};
\node[popdark, align=center] at (1.9,1.72) {\small $z = r(\cos\theta + i\sin\theta)$};
\node[popdark, align=center] at (1.9,1.49) {\small $z = re^{i\theta}$};
\end{tikzpicture}
\legende{One point, three names: the Cartesian, polar and exponential forms of $z$.}
\end{popfigure}

\begin{methode}[Converting $a + ib$ to polar form]
\begin{enumerate}
\item Compute the modulus $r = \sqrt{a^{2}+b^{2}}$ (always positive).
\item Find $\theta$ from the system $\cos\theta = \dfrac{a}{r}$, $\sin\theta = \dfrac{b}{r}$.
\item Use the \emph{signs} of $a$ and $b$ to place $\theta$ in the correct quadrant (a calculator only gives $\arctan(b/a)$, which may be wrong by $\pi$). Recall the CAST rule:
  \begin{itemize}
  \item $a>0,\ b>0$ (Q1): $\theta = \arctan(b/a)$
  \item $a<0,\ b>0$ (Q2): $\theta = \pi + \arctan(b/a)$
  \item $a<0,\ b<0$ (Q3): $\theta = -\pi + \arctan(b/a)$ (or $\pi + \arctan(b/a)$)
  \item $a>0,\ b<0$ (Q4): $\theta = \arctan(b/a)$ (negative)
  \end{itemize}
\item Write $z = r(\cos\theta + i\sin\theta)$ with $\theta \in {]{-\pi},\,\pi]}$.
\end{enumerate}
\end{methode}

It is essential to memorise the standard angles, which appear constantly at the GCE:
\[
\cos\tfrac{\pi}{6}=\tfrac{\sqrt3}{2},\quad \sin\tfrac{\pi}{6}=\tfrac12;\qquad
\cos\tfrac{\pi}{4}=\sin\tfrac{\pi}{4}=\tfrac{\sqrt2}{2};\qquad
\cos\tfrac{\pi}{3}=\tfrac12,\quad \sin\tfrac{\pi}{3}=\tfrac{\sqrt3}{2}.
\]

\begin{exemplebox}[Cartesian to polar]
Write $z = -1 + i\sqrt{3}$ in polar form.
\[
r = \sqrt{(-1)^{2}+(\sqrt3)^{2}} = \sqrt{1+3} = 2.
\]
Then $\cos\theta = -\tfrac12$ and $\sin\theta = \tfrac{\sqrt3}{2}$. Since $a<0$ and $b>0$, $\theta$ lies in the second quadrant: $\theta = \dfrac{2\pi}{3}$. Hence
\[
z = 2\left(\cos\frac{2\pi}{3} + i\sin\frac{2\pi}{3}\right).
\]
Conversely, polar to Cartesian is direct substitution: $2\left(\cos\frac{\pi}{4}+i\sin\frac{\pi}{4}\right) = 2\left(\frac{\sqrt2}{2}+i\frac{\sqrt2}{2}\right) = \sqrt2 + i\sqrt2$.
\end{exemplebox}

\begin{exemplebox}[Quadrant determination practice]
Write $z = -\sqrt{3} - i$ in polar form.
\[
r = \sqrt{3+1} = 2,\quad \cos\theta = -\frac{\sqrt3}{2},\quad \sin\theta = -\frac12.
\]
Both negative $\Rightarrow$ third quadrant. The reference angle is $\frac{\pi}{6}$, so $\theta = -\frac{5\pi}{6}$ (principal value). Thus
\[
z = 2\left(\cos\left(-\frac{5\pi}{6}\right) + i\sin\left(-\frac{5\pi}{6}\right)\right).
\]
\end{exemplebox}

\begin{attention}[The modulus must be positive]
A polar form requires $r > 0$. If a calculation produces something like $-2(\cos\frac{\pi}{6}+i\sin\frac{\pi}{6})$, this is \emph{not} a polar form. Absorb the minus sign into the angle by adding $\pi$:
\[
-2\left(\cos\tfrac{\pi}{6}+i\sin\tfrac{\pi}{6}\right) = 2\left(\cos\tfrac{7\pi}{6}+i\sin\tfrac{7\pi}{6}\right),
\]
because $\cos(\theta+\pi) = -\cos\theta$ and $\sin(\theta+\pi) = -\sin\theta$. Always check that your final $r$ is positive and $\theta \in {]{-\pi},\,\pi]}$.
\end{attention}

\courssub{Multiplication, Division and Powers in Polar Form}

Here lies the true power of the polar form. Take $z_1 = r_1(\cos\theta_1 + i\sin\theta_1)$ and $z_2 = r_2(\cos\theta_2 + i\sin\theta_2)$. Multiplying:
\begin{align*}
z_1 z_2 &= r_1 r_2\bigl[(\cos\theta_1\cos\theta_2 - \sin\theta_1\sin\theta_2) + i(\sin\theta_1\cos\theta_2 + \cos\theta_1\sin\theta_2)\bigr] \\
&= r_1 r_2\bigl[\cos(\theta_1+\theta_2) + i\sin(\theta_1+\theta_2)\bigr],
\end{align*}
where we used the addition formulae $\cos(A+B) = \cos A\cos B - \sin A\sin B$ and $\sin(A+B) = \sin A\cos B + \cos A\sin B$.

\begin{propriete}[Multiplication and division in polar form]
For non-zero complex numbers $z_1 = r_1(\cos\theta_1 + i\sin\theta_1)$ and $z_2 = r_2(\cos\theta_2 + i\sin\theta_2)$:
\[
z_1 z_2 = r_1 r_2\bigl(\cos(\theta_1+\theta_2) + i\sin(\theta_1+\theta_2)\bigr), \qquad
\frac{z_1}{z_2} = \frac{r_1}{r_2}\bigl(\cos(\theta_1-\theta_2) + i\sin(\theta_1-\theta_2)\bigr).
\]
In words: \cle{moduli multiply (or divide), arguments add (or subtract)}. The proof of the product formula, shown above, is examinable; the quotient formula follows by multiplying numerator and denominator by the conjugate of $z_2$.
\end{propriete}

Geometrically, multiplying by $z_2$ \emph{rotates} the plane through the angle $\theta_2$ and \emph{scales} distances by the factor $r_2$. This is why engineers and physicists love complex numbers: rotation becomes multiplication.

Applying the product rule repeatedly with the same number gives a preview of De Moivre's theorem (studied in the next section):
\[
\bigl[r(\cos\theta + i\sin\theta)\bigr]^{n} = r^{n}\bigl(\cos n\theta + i\sin n\theta\bigr) \qquad (n \in \mathbb{N}).
\]

\begin{exemplebox}[A power made easy]
Compute $(1+i)^{8}$. In Cartesian form this would require the binomial theorem with nine terms. In polar form: $1+i = \sqrt2\left(\cos\frac{\pi}{4}+i\sin\frac{\pi}{4}\right)$, so
\[
(1+i)^{8} = (\sqrt2)^{8}\left(\cos\frac{8\pi}{4}+i\sin\frac{8\pi}{4}\right) = 16(\cos 2\pi + i\sin 2\pi) = 16.
\]
\end{exemplebox}

\begin{experience}[Geometric interpretation of multiplication]
On graph paper, plot $z_1 = 1+i$ and $z_2 = \sqrt{3}+i$. Convert both to polar form, compute the product $z_1z_2$ using the polar rule, and plot the result. Measure the angle and distance from the origin. You will see that the argument of the product is the sum of the arguments, and the modulus is the product of the moduli. This visual confirmation is worth a thousand algebraic proofs.
\end{experience}

\courssub{Euler's Formula and the Exponential Form}

\textbf{A remarkable observation.}\par
Define, for real $\theta$, the function $f(\theta) = \cos\theta + i\sin\theta$. The multiplication rule above says $f(\theta_1)f(\theta_2) = f(\theta_1+\theta_2)$: $f$ \emph{transforms addition into multiplication}. But this is exactly the characteristic law of an exponential function, $e^{x_1}e^{x_2} = e^{x_1+x_2}$. Moreover, differentiating formally, $f'(\theta) = -\sin\theta + i\cos\theta = i(\cos\theta + i\sin\theta) = i\,f(\theta)$, so $f$ satisfies the differential equation $f' = i f$ with $f(0)=1$ — the same equation satisfied by $e^{i\theta}$. The great Swiss mathematician Leonhard Euler turned this observation into a definition.

\begin{propriete}[Euler's formula]
For every real number $\theta$,
\[
\boxed{\,e^{i\theta} = \cos\theta + i\sin\theta\,}.
\]
Consequently every non-zero complex number can be written in \cle{exponential form}
\[
z = re^{i\theta}, \qquad r = |z| > 0,\quad \theta = \arg(z).
\]
(The proof, via power series, is beyond the syllabus; the formula itself must be known and used fluently.)
\end{propriete}

\begin{attention}[$e^{i\theta}$ is not an ordinary exponential]
The symbol $e^{i\theta}$ is a \emph{complex} number of modulus $1$; it is not the real exponential function. In particular you may \textbf{not} apply $\ln$ to it as if it were a positive real number, and $e^{i\theta}$ is never $0$, never negative — it simply travels around the unit circle as $\theta$ varies.
\end{attention}

\begin{popfigure}
\begin{tikzpicture}[scale=2.4, >=stealth]
\draw[->, popmuted] (-1.6,0) -- (1.6,0) node[below left, popink]{\small Re};
\draw[->, popmuted] (0,-1.4) -- (0,1.5) node[below left, popink]{\small Im};
\draw[popblue, thick] (0,0) circle (1);
\coordinate (M) at (0.5,0.866);
\draw[popdark, thick] (0,0) -- (M);
\fill[poporange] (M) circle (0.03) node[above right, popdark]{$e^{i\theta}$};
\draw[->, popgreen, thick] (0.45,0) arc (0:60:0.45) node[midway, right, popgreen]{$\theta$};
\draw[popmuted, dashed] (M) -- (0.5,0) node[below, popdark]{\small $\cos\theta$};
\fill[popink] (-1,0) circle (0.035) node[below left, popink]{$e^{i\pi} = -1$};
\fill[popmuted] (1,0) circle (0.025) node[below right, popmuted]{$1$};
\fill[popmuted] (0,1) circle (0.025) node[above right, popmuted]{$i$};
\fill[popmuted] (0,-1) circle (0.025) node[below right, popmuted]{$-i$};
\end{tikzpicture}
\legende{As $\theta$ varies, $e^{i\theta} = \cos\theta + i\sin\theta$ describes the unit circle. At $\theta = \pi$ it passes through $-1$.}
\end{popfigure}

Setting $\theta = \pi$ in Euler's formula gives $\cos\pi + i\sin\pi = -1 + 0i$, hence the most celebrated identity in all of mathematics.

\begin{aretenir}[Euler's identity]
\[
e^{i\pi} + 1 = 0.
\]
One short equation unites the five fundamental constants $0, 1, i, \pi, e$, together with the three basic operations (addition, multiplication, exponentiation). It is often called the most beautiful formula in mathematics.
\end{aretenir}

\begin{savaistu}[Euler and the language of physics]
Leonhard Euler (1707--1783) was the most prolific mathematician in history. The notation $e^{i\theta}$ that he introduced is today the working language of electrical engineering (alternating current is written $U = U_0 e^{i\omega t}$), of signal processing behind every mobile phone, and of quantum mechanics. The radio masts you see across Douala and Yaound\'e are designed with calculations written in exponential form.
\end{savaistu}

\courssub{Algebraic Rules of the Exponential Form}

All the polar rules become, in exponential notation, the familiar index laws — which is precisely why the notation $e^{i\theta}$ was chosen.

\begin{propriete}[Rules of computation with $e^{i\theta}$]
For real $\theta, \theta_1, \theta_2$ and non-zero $z_1 = r_1 e^{i\theta_1}$, $z_2 = r_2 e^{i\theta_2}$:
\[
e^{i\theta_1}\,e^{i\theta_2} = e^{i(\theta_1+\theta_2)}, \qquad
\frac{1}{e^{i\theta}} = e^{-i\theta}, \qquad
\frac{z_1}{z_2} = \frac{r_1}{r_2}\,e^{i(\theta_1-\theta_2)}, \qquad
\overline{re^{i\theta}} = re^{-i\theta}, \qquad
\left(re^{i\theta}\right)^{n} = r^{n}e^{in\theta}.
\]
Also $|e^{i\theta}| = 1$ for every real $\theta$.
\end{propriete}

These are exactly the polar rules rewritten: the conjugation rule, for instance, is just $\overline{r(\cos\theta+i\sin\theta)} = r(\cos\theta - i\sin\theta) = r(\cos(-\theta) + i\sin(-\theta)) = re^{-i\theta}$, since cosine is even and sine is odd.

\begin{popfigure}
\begin{tikzpicture}[>=stealth, node distance=1.1cm]
\node[draw=popblue, very thick, rounded corners, fill=popblueL, align=center, text width=3.4cm] (cart) {\textbf{Cartesian}\\ $z = a + ib$};
\node[draw=popgreen, very thick, rounded corners, fill=popgreenL, align=center, text width=3.4cm, right=3.2cm of cart] (pol) {\textbf{Polar}\\ $z = r(\cos\theta + i\sin\theta)$};
\node[draw=poporange, very thick, rounded corners, fill=white, align=center, text width=3.4cm, right=3.2cm of pol] (exp) {\textbf{Exponential}\\ $z = re^{i\theta}$};
\draw[->, very thick, popdark] (cart) -- (pol) node[midway, above, align=center, text width=3.1cm]{\small $r=\sqrt{a^2+b^2}$; $\cos\theta=\frac{a}{r}$, $\sin\theta=\frac{b}{r}$};
\draw[->, very thick, popdark] (pol) -- (cart) node[midway, below, align=center]{\small $a = r\cos\theta$, $b = r\sin\theta$};
\draw[->, very thick, popdark] (pol) -- (exp) node[midway, above, align=center]{\small Euler: $\cos\theta + i\sin\theta = e^{i\theta}$};
\draw[->, very thick, popdark] (exp) -- (pol) node[midway, below, align=center]{\small $e^{i\theta} = \cos\theta + i\sin\theta$};
\end{tikzpicture}
\legende{Conversion flowchart between the three forms of a complex number.}
\end{popfigure}

\begin{exoresolu}[Computing with the exponential form]
Given $z_1 = 2e^{i\pi/3}$ and $z_2 = \sqrt{2}\,e^{-i\pi/4}$, determine in exponential form and then in Cartesian form: \textbf{(a)} $z_1 z_2$; \textbf{(b)} $\dfrac{z_1}{z_2}$; \textbf{(c)} $z_1^{3}$.
\tcblower
\textbf{(a)} Moduli multiply, arguments add:
\[
z_1 z_2 = 2\sqrt{2}\,e^{i\left(\frac{\pi}{3}-\frac{\pi}{4}\right)} = 2\sqrt{2}\,e^{i\pi/12}.
\]
\textbf{(b)} Moduli divide, arguments subtract:
\[
\frac{z_1}{z_2} = \frac{2}{\sqrt2}\,e^{i\left(\frac{\pi}{3}+\frac{\pi}{4}\right)} = \sqrt{2}\,e^{i\,7\pi/12}.
\]
\textbf{(c)} Power rule:
\[
z_1^{3} = 2^{3}e^{i\pi} = 8(\cos\pi + i\sin\pi) = -8.
\]
For the Cartesian form of (a): $\cos\frac{\pi}{12} = \frac{\sqrt6+\sqrt2}{4}$ and $\sin\frac{\pi}{12} = \frac{\sqrt6-\sqrt2}{4}$, so
\[
z_1z_2 = 2\sqrt2\left(\frac{\sqrt6+\sqrt2}{4} + i\,\frac{\sqrt6-\sqrt2}{4}\right) = \frac{2\sqrt3+2}{2} + i\,\frac{2\sqrt3-2}{2} = (\sqrt3+1) + i(\sqrt3-1).
\]
\end{exoresolu}

\begin{methode}[GCE strategy: choosing the most convenient form]
\begin{enumerate}
\item \textbf{Addition and subtraction}: use the Cartesian form $a+ib$ (real parts and imaginary parts combine separately).
\item \textbf{Multiplication, division, powers, roots}: convert to polar or exponential form first; the computation reduces to arithmetic on $r$ and $\theta$.
\item \textbf{Conjugation and modulus}: both forms work, but $\overline{re^{i\theta}} = re^{-i\theta}$ and $|re^{i\theta}| = r$ are instantaneous.
\item \textbf{Final answer}: give the form requested by the question; if none is requested, $a+ib$ with exact values is the safest.
\end{enumerate}
\end{methode}

\begin{attention}[Frequent errors at the GCE]
\begin{itemize}
\item Writing $z = r(\cos\theta - i\sin\theta)$ when $\theta$ is in a lower quadrant instead of using the negative angle $\theta$ itself: the sign belongs to the \emph{angle}, not to the formula.
\item Forgetting to check the quadrant: $\arctan$ alone cannot distinguish $\frac{\pi}{3}$ from $\frac{4\pi}{3}$.
\item Mixing forms inside one computation, e.g.\ adding $re^{i\theta}$ to $a+ib$ directly. Convert first.
\item Claiming $e^{i\theta} > 0$ or solving $e^{i\theta} = 2$: impossible, since $|e^{i\theta}| = 1$.
\item Forgetting to reduce the argument to the principal range ${]{-\pi},\,\pi]}$ in the final answer.
\end{itemize}
\end{attention}

\begin{exercice}[Practice]
\begin{enumerate}
\item Write in polar and in exponential form: (a) $z = 1 - i$; (b) $z = -2\sqrt3 - 2i$; (c) $z = -5i$.
\item Given $z_1 = 3e^{i\pi/6}$ and $z_2 = 2e^{i\pi/2}$, compute $z_1z_2$, $\dfrac{z_1}{z_2}$ and $z_2^{4}$, giving each answer in exponential and in Cartesian form.
\item Show that $(1+i\sqrt3)^{6}$ is a real number and find its value.
\item Simplify $\dfrac{e^{i\theta}-e^{-i\theta}}{2i}$ and $\dfrac{e^{i\theta}+e^{-i\theta}}{2}$. Which familiar functions do you recognise?
\item Express $\cos^3\theta$ in terms of $\cos3\theta$ and $\cos\theta$ using Euler's relations. (Hint: start from $2\cos\theta = e^{i\theta}+e^{-i\theta}$ and cube both sides.)
\end{enumerate}
\end{exercice}

\begin{corrige}[Exercise 1]
\textbf{1.} (a) $r = \sqrt2$, $\cos\theta = \frac{1}{\sqrt2}$, $\sin\theta = -\frac{1}{\sqrt2}$, so $\theta = -\frac{\pi}{4}$: $z = \sqrt2\left(\cos\frac{\pi}{4} - i\sin\frac{\pi}{4}\right) = \sqrt2\,e^{-i\pi/4}$.
(b) $r = \sqrt{12+4} = 4$; $\cos\theta = -\frac{\sqrt3}{2}$, $\sin\theta = -\frac12$, third quadrant: $\theta = -\frac{5\pi}{6}$ (or $\frac{7\pi}{6}$). Hence $z = 4e^{-5i\pi/6}$.
(c) $r = 5$, $\theta = -\frac{\pi}{2}$: $z = 5e^{-i\pi/2}$.

\textbf{2.} $z_1z_2 = 6e^{i(\pi/6+\pi/2)} = 6e^{2i\pi/3} = 6\left(-\frac12 + i\frac{\sqrt3}{2}\right) = -3 + 3i\sqrt3$.
$\dfrac{z_1}{z_2} = \frac32 e^{i(\pi/6-\pi/2)} = \frac32 e^{-i\pi/3} = \frac32\left(\frac12 - i\frac{\sqrt3}{2}\right) = \frac34 - \frac{3\sqrt3}{4}i$.
$z_2^{4} = 16e^{2i\pi} = 16$.

\textbf{3.} $1+i\sqrt3 = 2e^{i\pi/3}$, so $(1+i\sqrt3)^{6} = 2^{6}e^{2i\pi} = 64(\cos 2\pi + i\sin 2\pi) = 64 \in \mathbb{R}$.

\textbf{4.} Since $e^{i\theta} = \cos\theta + i\sin\theta$ and $e^{-i\theta} = \cos\theta - i\sin\theta$, subtracting gives $e^{i\theta}-e^{-i\theta} = 2i\sin\theta$ and adding gives $e^{i\theta}+e^{-i\theta} = 2\cos\theta$. Hence
\[
\frac{e^{i\theta}-e^{-i\theta}}{2i} = \sin\theta, \qquad \frac{e^{i\theta}+e^{-i\theta}}{2} = \cos\theta.
\]
These are \emph{Euler's relations}, which will be our main tool for expressing powers of $\cos x$ and $\sin x$ in terms of multiple angles later in this chapter.

\textbf{5.} $2\cos\theta = e^{i\theta}+e^{-i\theta}$. Cubing: $8\cos^3\theta = (e^{i\theta}+e^{-i\theta})^3 = e^{3i\theta} + 3e^{i\theta} + 3e^{-i\theta} + e^{-3i\theta} = (e^{3i\theta}+e^{-3i\theta}) + 3(e^{i\theta}+e^{-i\theta}) = 2\cos3\theta + 6\cos\theta$. Hence $\cos^3\theta = \frac{1}{4}\cos3\theta + \frac{3}{4}\cos\theta$.
\end{corrige}

\begin{aretenir}[Key points of this section]
\begin{itemize}
\item Polar form: $z = r(\cos\theta + i\sin\theta)$ with $r = |z| > 0$, $\theta = \arg(z) \in {]{-\pi},\,\pi]}$.
\item Euler's formula: $e^{i\theta} = \cos\theta + i\sin\theta$; exponential form $z = re^{i\theta}$; and $e^{i\pi} + 1 = 0$.
\item Moduli multiply/divide, arguments add/subtract; $\overline{re^{i\theta}} = re^{-i\theta}$; $(re^{i\theta})^{n} = r^{n}e^{in\theta}$.
\item Convert to exponential form for products, quotients and powers; stay in Cartesian form for sums.
\item Euler's relations: $\cos\theta = \frac{e^{i\theta}+e^{-i\theta}}{2}$, $\sin\theta = \frac{e^{i\theta}-e^{-i\theta}}{2i}$.
\end{itemize}
\end{aretenir}

\coursec{De Moivre's Theorem and Trigonometric Applications}

In the previous section we discovered that a complex number of modulus $1$ can be written in the wonderfully compact form $e^{i\theta} = \cos\theta + i\sin\theta$, and that multiplication of such numbers simply \emph{adds} their arguments. If we multiply $\cos\theta + i\sin\theta$ by itself $n$ times, the argument $\theta$ is added to itself $n$ times, giving $n\theta$. This simple observation is the content of one of the most beautiful theorems of the Upper Sixth programme: \cle{De Moivre's theorem}. It will allow us to compute huge powers of complex numbers in seconds, to find roots of unity, and to transform powers of $\cos x$ and $\sin x$ into multiple angles (and vice versa) — a skill you will use constantly in integration.

\courssub{De Moivre's Theorem}

\begin{experience}[A motivating computation]
Let $z = \cos\theta + i\sin\theta$. Compute the first powers using the addition formulae:
\begin{align*}
z^2 &= (\cos\theta + i\sin\theta)^2 = (\cos^2\theta - \sin^2\theta) + i(2\sin\theta\cos\theta) = \cos 2\theta + i\sin 2\theta,\\
z^3 &= z^2 \cdot z = (\cos 2\theta + i\sin 2\theta)(\cos\theta + i\sin\theta) = \cos 3\theta + i\sin 3\theta,
\end{align*}
since multiplying two unit-modulus numbers adds their arguments. A pattern is emerging: $z^n = \cos n\theta + i\sin n\theta$. Let us prove it properly.
\end{experience}

\begin{propriete}[De Moivre's Theorem — proof examinable]
For every real number $\theta$ and every integer $n \in \mathbb{Z}$,
\[ (\cos\theta + i\sin\theta)^n = \cos n\theta + i\sin n\theta. \]
Equivalently, $\left(e^{i\theta}\right)^n = e^{in\theta}$.
\end{propriete}

\textbf{Proof for $n \in \mathbb{N}$ (by induction).}\par
Let $P(n)$ be the statement ``$(\cos\theta + i\sin\theta)^n = \cos n\theta + i\sin n\theta$''.

\emph{Base case $n = 0$:} $(\cos\theta + i\sin\theta)^0 = 1 = \cos 0 + i\sin 0$. So $P(0)$ holds. (The case $n=1$ is also trivially true.)

\emph{Inductive step:} assume $P(k)$ true for some $k \geq 0$. Then
\begin{align*}
(\cos\theta + i\sin\theta)^{k+1} &= (\cos k\theta + i\sin k\theta)(\cos\theta + i\sin\theta)\\
&= (\cos k\theta\cos\theta - \sin k\theta\sin\theta) + i(\sin k\theta\cos\theta + \cos k\theta\sin\theta)\\
&= \cos(k\theta + \theta) + i\sin(k\theta + \theta) = \cos(k+1)\theta + i\sin(k+1)\theta,
\end{align*}
using the addition formulae $\cos(A+B)=\cos A\cos B-\sin A\sin B$ and $\sin(A+B)=\sin A\cos B+\cos A\sin B$. Hence $P(k) \Rightarrow P(k+1)$, and by induction $P(n)$ holds for all $n \in \mathbb{N}$.

\textbf{Extension to negative integers.}\par
Let $n = -m$ with $m \in \mathbb{N}$, $m \geq 1$. Since $|\cos\theta + i\sin\theta| = 1$, its inverse equals its conjugate:
\[ (\cos\theta + i\sin\theta)^{-1} = \cos\theta - i\sin\theta = \cos(-\theta) + i\sin(-\theta). \]
Therefore
\[ (\cos\theta + i\sin\theta)^{n} = \bigl(\cos(-\theta) + i\sin(-\theta)\bigr)^{m} = \cos(-m\theta) + i\sin(-m\theta) = \cos n\theta + i\sin n\theta, \]
applying the (already proved) result for the positive integer $m$ to the angle $-\theta$. The theorem thus holds for all $n \in \mathbb{Z}$. $\blacksquare$

\begin{popfigure}
\begin{center}
\begin{tikzpicture}[scale=2.6]
\draw[->, popmuted] (-1.4,0) -- (1.5,0) node[below left, black]{\small Re};
\draw[->, popmuted] (0,-1.3) -- (0,1.5) node[below left, black]{\small Im};
\draw[popline] (0,0) circle (1);
\draw[->, thick, popblue] (0,0) -- (0.7071,0.7071) node[midway, below left, popblue]{\small $z$};
\draw[->, thick, poporange] (0,0) -- (-0.7071,0.7071) node[midway, above left, poporange]{\small $z^2$};
\draw[->, thick, popgreen] (0,0) -- (-1,0) node[midway, above, popgreen]{\small $z^3$};
\draw[->, poppurple] (0.25,0) arc (0:45:0.25) node[midway, right, poppurple]{\small $\theta$};
\draw[->, poppurple] (0.4,0) arc (0:135:0.4) node[midway, above, poppurple]{\small $3\theta$};
\fill[popblue] (0.7071,0.7071) circle (0.02);
\fill[poporange] (-0.7071,0.7071) circle (0.02);
\fill[popgreen] (-1,0) circle (0.02);
\end{tikzpicture}
\end{center}
\legende{De Moivre on the unit circle: with $\theta = 45^{\circ}$, each multiplication by $z = \cos\theta + i\sin\theta$ rotates by a further $\theta$, so $z^n$ has argument $n\theta$.}
\end{popfigure}

\begin{attention}[Two classic traps]
\begin{itemize}
\item $(\cos\theta + i\sin\theta)^n \neq \cos^n\theta + i\sin^n\theta$. De Moivre multiplies the \emph{angle} by $n$; it does not raise $\cos\theta$ and $\sin\theta$ to the power $n$.
\item In its simple form the theorem is stated for \cle{integer} $n$. For fractional powers (e.g. square roots, cube roots) several distinct values appear; this is the subject of the next section on $n$th roots.
\end{itemize}
\end{attention}

\courssub{Application: Computing Large Powers}

\begin{methode}[Computing $z^n$ via De Moivre]
\begin{enumerate}
\item Write $z$ in polar form: find $r = |z|$ and $\theta = \arg(z)$, so $z = r(\cos\theta + i\sin\theta)$.
\item Apply De Moivre: $z^n = r^n(\cos n\theta + i\sin n\theta)$.
\item Reduce $n\theta$ modulo $2\pi$ and convert back to Cartesian form if required.
\end{enumerate}
\end{methode}

\begin{exoresolu}[Computing $(1+i)^{20}$ and $(\sqrt{3}+i)^{12}$]
Compute (a) $(1+i)^{20}$ and (b) $(\sqrt{3}+i)^{12}$, giving exact answers in Cartesian form.
\tcblower
\textbf{(a)} For $z = 1+i$: $r = \sqrt{1^2+1^2} = \sqrt{2}$ and $\theta = \arctan 1 = \dfrac{\pi}{4}$ (first quadrant). Hence
\[ (1+i)^{20} = \left(\sqrt{2}\right)^{20}\left(\cos\tfrac{20\pi}{4} + i\sin\tfrac{20\pi}{4}\right) = 2^{10}(\cos 5\pi + i\sin 5\pi) = 1024(-1 + 0i) = \boxed{-1024}. \]
\textbf{(b)} For $w = \sqrt{3}+i$: $r = \sqrt{3+1} = 2$ and $\theta = \arctan\dfrac{1}{\sqrt{3}} = \dfrac{\pi}{6}$. Hence
\[ (\sqrt{3}+i)^{12} = 2^{12}\left(\cos\tfrac{12\pi}{6} + i\sin\tfrac{12\pi}{6}\right) = 4096(\cos 2\pi + i\sin 2\pi) = \boxed{4096}. \]
Notice how hopeless a direct binomial expansion of $(\sqrt3+i)^{12}$ would have been: polar form plus De Moivre reduces the problem to one line.
\end{exoresolu}

\begin{exemplebox}[Preview: roots of unity]
Solve $z^3 = 1$ in $\mathbb{C}$. Write $z = \cos\theta + i\sin\theta$ (any cube root of $1$ has modulus $1$). By De Moivre, $z^3 = \cos 3\theta + i\sin 3\theta = 1$, so $3\theta = 2k\pi$, i.e. $\theta = \dfrac{2k\pi}{3}$. Taking $k = 0, 1, 2$ gives the three \cle{cube roots of unity}:
\[ 1, \qquad \cos\tfrac{2\pi}{3} + i\sin\tfrac{2\pi}{3} = -\tfrac12 + \tfrac{\sqrt3}{2}i, \qquad \cos\tfrac{4\pi}{3} + i\sin\tfrac{4\pi}{3} = -\tfrac12 - \tfrac{\sqrt3}{2}i. \]
Further values of $k$ repeat these. We shall generalise this to $n$th roots in the next section.
\end{exemplebox}

\courssub{From Powers of cos and sin to Multiple Angles (Linearisation)}

Set $z = \cos x + i\sin x = e^{ix}$, so that $|z| = 1$ and $z^{-1} = \cos x - i\sin x = \overline{z}$. Adding and subtracting,
\[ z + \frac{1}{z} = 2\cos x, \qquad z - \frac{1}{z} = 2i\sin x. \]
Applying De Moivre with integer $n$, $z^n = \cos nx + i\sin nx$ and $z^{-n} = \cos nx - i\sin nx$, so:

\begin{propriete}[The key pairing identities]
With $z = \cos x + i\sin x$ and $n \in \mathbb{Z}$,
\[ z^n + \frac{1}{z^n} = 2\cos nx, \qquad z^n - \frac{1}{z^n} = 2i\sin nx. \]
(Proof: immediate from De Moivre; quote it in the examination.)
\end{propriete}

\begin{popfigure}
\begin{center}
\begin{tabular}{|c|c|c|}
\hline
\rowcolor{popblueL} $n$ & $z^n + z^{-n} = 2\cos nx$ & $z^n - z^{-n} = 2i\sin nx$ \\
\hline
$1$ & $z + z^{-1} = 2\cos x$ & $z - z^{-1} = 2i\sin x$ \\
\hline
$2$ & $z^2 + z^{-2} = 2\cos 2x$ & $z^2 - z^{-2} = 2i\sin 2x$ \\
\hline
$3$ & $z^3 + z^{-3} = 2\cos 3x$ & $z^3 - z^{-3} = 2i\sin 3x$ \\
\hline
\end{tabular}
\end{center}
\legende{The pairing identities for $n = 1, 2, 3$: conjugate pairs $z^{\pm n}$ recombine into a single cosine or sine of a multiple angle.}
\end{popfigure}

\begin{methode}[Linearising $\cos^n x$ or $\sin^n x$]
\begin{enumerate}
\item Write $\cos x = \dfrac{1}{2}\left(z + \dfrac{1}{z}\right)$ or $\sin x = \dfrac{1}{2i}\left(z - \dfrac{1}{z}\right)$.
\item Raise to the power $n$ and expand with the binomial theorem: $(A \pm B)^n = \sum_{k=0}^n \binom{n}{k} A^{n-k}(\pm B)^k$.
\item Regroup conjugate pairs $z^k + z^{-k}$ (or $z^k - z^{-k}$) and replace them by $2\cos kx$ (or $2i\sin kx$).
\item Simplify the numerical factors.
\end{enumerate}
\end{methode}

\begin{exoresolu}[Linearising $\cos^4 x$]
Express $\cos^4 x$ as a sum of cosines of multiple angles.
\tcblower
Let $z = \cos x + i\sin x$, so $\cos x = \dfrac{1}{2}\left(z + z^{-1}\right)$. Then
\begin{align*}
\cos^4 x &= \frac{1}{16}\left(z + \frac{1}{z}\right)^4 = \frac{1}{16}\left(z^4 + 4z^2 + 6 + \frac{4}{z^2} + \frac{1}{z^4}\right)\\
&= \frac{1}{16}\left[\left(z^4 + \frac{1}{z^4}\right) + 4\left(z^2 + \frac{1}{z^2}\right) + 6\right] = \frac{1}{16}\left[2\cos 4x + 8\cos 2x + 6\right].
\end{align*}
Hence $\boxed{\cos^4 x = \dfrac{1}{8}\cos 4x + \dfrac{1}{2}\cos 2x + \dfrac{3}{8}}$.
\end{exoresolu}

\begin{exoresolu}[Linearising $\sin^5 x$]
Express $\sin^5 x$ in terms of sines of multiple angles.
\tcblower
With $\sin x = \dfrac{1}{2i}\left(z - z^{-1}\right)$:
\begin{align*}
\sin^5 x &= \frac{1}{(2i)^5}\left(z - \frac{1}{z}\right)^5 = \frac{1}{32i}\left(z^5 - 5z^3 + 10z - \frac{10}{z} + \frac{5}{z^3} - \frac{1}{z^5}\right)\\
&= \frac{1}{32i}\left[\left(z^5 - \frac{1}{z^5}\right) - 5\left(z^3 - \frac{1}{z^3}\right) + 10\left(z - \frac{1}{z}\right)\right]\\
&= \frac{1}{32i}\left[2i\sin 5x - 10i\sin 3x + 20i\sin x\right].
\end{align*}
Hence $\boxed{\sin^5 x = \dfrac{1}{16}\sin 5x - \dfrac{5}{16}\sin 3x + \dfrac{5}{8}\sin x}$.
\end{exoresolu}

\begin{exemplebox}[Linearising $\sin^2 x \cos^2 x$]
$\sin^2 x \cos^2 x = \dfrac{1}{4}\sin^2 2x = \dfrac{1}{4}\cdot\dfrac{1 - \cos 4x}{2} = \dfrac{1}{8} - \dfrac{1}{8}\cos 4x$. Alternatively, directly with $z$: $\sin^2x\cos^2x = \dfrac{1}{(2i)^2\cdot 2^2}(z - z^{-1})^2(z+z^{-1})^2 = -\dfrac{1}{16}(z^2 - z^{-2})^2 = \dfrac{1}{8} - \dfrac{1}{8}\cos 4x$.
\end{exemplebox}

\begin{savaistu}[Why linearisation matters]
Integrals such as $\displaystyle\int \cos^4 x\, dx$ cannot be done by substitution, but once linearised they are immediate: $\int \cos^4 x\,dx = \dfrac{1}{32}\sin 4x + \dfrac{1}{4}\sin 2x + \dfrac{3}{8}x + C$. These expansions are examined regularly in the Pure and Further papers.
\end{savaistu}

\courssub{From Multiple Angles to Powers of cos and sin}

The reverse direction uses De Moivre together with the \cle{binomial theorem}: expand $(\cos x + i\sin x)^n$, then equate real and imaginary parts with $\cos nx + i\sin nx$.

\begin{methode}[Expanding $\cos nx$ and $\sin nx$]
\begin{enumerate}
\item Write $\cos nx + i\sin nx = (\cos x + i\sin x)^n$ (De Moivre).
\item Expand the right-hand side with the binomial theorem, writing $c = \cos x$, $s = \sin x$:
  \[ (c + is)^n = \sum_{k=0}^n \binom{n}{k} c^{n-k} (is)^k. \]
\item Equate real parts to obtain $\cos nx$, and imaginary parts to obtain $\sin nx$.
\item Use $i^2 = -1$ to simplify, and $\sin^2 x = 1 - \cos^2 x$ (or $\cos^2 x = 1 - \sin^2 x$) if a single-variable polynomial is wanted.
\end{enumerate}
\end{methode}

\begin{exoresolu}[The triple-angle formulae]
Prove that $\cos 3x = 4\cos^3 x - 3\cos x$ and $\sin 3x = 3\sin x - 4\sin^3 x$.
\tcblower
By De Moivre and the binomial theorem, with $c = \cos x$, $s = \sin x$:
\[ \cos 3x + i\sin 3x = (c + is)^3 = c^3 + 3c^2(is) + 3c(is)^2 + (is)^3 = (c^3 - 3cs^2) + i(3c^2s - s^3). \]
Equating real parts: $\cos 3x = c^3 - 3cs^2 = c^3 - 3c(1 - c^2) = 4c^3 - 3c$, i.e.
\[ \boxed{\cos 3x = 4\cos^3 x - 3\cos x}. \]
Equating imaginary parts: $\sin 3x = 3c^2s - s^3 = 3(1 - s^2)s - s^3 = 3s - 4s^3$, i.e.
\[ \boxed{\sin 3x = 3\sin x - 4\sin^3 x}. \qquad \blacksquare \]
\end{exoresolu}

\begin{exemplebox}[Obtaining $\tan nx$]
Dividing the expansions of $\sin nx$ and $\cos nx$ and then dividing numerator and denominator by $\cos^n x$ expresses $\tan nx$ in terms of $t = \tan x$. For $n = 3$:
\[ \tan 3x = \frac{3c^2s - s^3}{c^3 - 3cs^2} = \frac{3t - t^3}{1 - 3t^2}, \qquad t = \tan x. \]
The same technique gives $\tan 2x = \dfrac{2t}{1 - t^2}$, a formula you already know — now derived from De Moivre.
\end{exemplebox}

\begin{attention}[Common mistakes in expansions]
\begin{itemize}
\item Forgetting that $(is)^2 = -s^2$ and $(is)^3 = -is^3$: the signs of the real part alternate.
\item Mixing the two directions: to \emph{linearise} $\cos^n x$ start from $z + z^{-1}$; to \emph{expand} $\cos nx$ start from $(c + is)^n$.
\item When equating parts, remember $i\sin nx$ is the \emph{imaginary} part of the expansion, so $\sin nx$ equals the coefficient of $i$.
\end{itemize}
\end{attention}

\begin{exercice}[Practice]
\begin{enumerate}
\item Compute $(1 - i\sqrt{3})^{15}$ in Cartesian form.
\item Express $\cos^5 x$ as a sum of cosines of multiple angles.
\item Express $\sin^4 x$ as a sum of cosines of multiple angles, and hence evaluate $\displaystyle\int_0^{\pi/2} \sin^4 x\, dx$.
\item Prove that $\cos 4x = 8\cos^4 x - 8\cos^2 x + 1$ and find the corresponding expression for $\tan 4x$ in terms of $t = \tan x$.
\end{enumerate}
\end{exercice}

\begin{corrige}[Exercise 1]
\textbf{1.} $1 - i\sqrt{3}$ has modulus $2$ and argument $-\dfrac{\pi}{3}$. Hence $(1 - i\sqrt{3})^{15} = 2^{15}\left(\cos(-5\pi) + i\sin(-5\pi)\right) = 32768(-1) = \boxed{-32768}$.

\textbf{2.} $\cos^5 x = \dfrac{1}{32}\left(z + z^{-1}\right)^5 = \dfrac{1}{32}\left[\left(z^5 + z^{-5}\right) + 5\left(z^3 + z^{-3}\right) + 10\left(z + z^{-1}\right)\right] = \boxed{\dfrac{1}{16}\cos 5x + \dfrac{5}{16}\cos 3x + \dfrac{5}{8}\cos x}$.

\textbf{3.} $\sin^4 x = \dfrac{1}{16}\left(z - z^{-1}\right)^4 = \dfrac{1}{16}\left[\left(z^4 + z^{-4}\right) - 4\left(z^2 + z^{-2}\right) + 6\right] = \boxed{\dfrac{1}{8}\cos 4x - \dfrac{1}{2}\cos 2x + \dfrac{3}{8}}$. Then $\displaystyle\int_0^{\pi/2}\sin^4 x\,dx = \left[\tfrac{1}{32}\sin 4x - \tfrac14\sin 2x + \tfrac{3}{8}x\right]_0^{\pi/2} = \boxed{\dfrac{3\pi}{16}}$.

\textbf{4.} $(c + is)^4 = (c^4 - 6c^2s^2 + s^4) + i(4c^3s - 4cs^3)$. Real part: $\cos 4x = c^4 - 6c^2(1-c^2) + (1-c^2)^2 = 8c^4 - 8c^2 + 1$, i.e. $\boxed{\cos 4x = 8\cos^4 x - 8\cos^2 x + 1}$. Dividing imaginary by real parts and then by $c^4$: $\tan 4x = \dfrac{4t - 4t^3}{1 - 6t^2 + t^4}$, i.e. $\boxed{\tan 4x = \dfrac{4\tan x - 4\tan^3 x}{1 - 6\tan^2 x + \tan^4 x}}$.
\end{corrige}

\begin{aretenir}[Key points of the section]
\begin{itemize}
\item De Moivre: $(\cos\theta + i\sin\theta)^n = \cos n\theta + i\sin n\theta$ for all $n \in \mathbb{Z}$ (proved by induction, extended to negative $n$ via the conjugate).
\item Large powers: convert to polar form, apply De Moivre, reduce the angle modulo $2\pi$.
\item With $z = e^{ix}$: $z^n + z^{-n} = 2\cos nx$ and $z^n - z^{-n} = 2i\sin nx$ — the engine of linearisation.
\item $\cos nx$ and $\sin nx$ are obtained by expanding $(\cos x + i\sin x)^n$ binomially and equating real and imaginary parts; $\tan nx$ follows by division.
\end{itemize}
\end{aretenir}

\coursec{Roots of Complex Numbers: nth Roots and Cube Roots}

In the previous section, De Moivre's theorem taught us how to \emph{raise} a complex number to a power: $(re^{i\theta})^n = r^n e^{in\theta}$. We now ask the reverse question, which is far richer: given a complex number $w$, can we find \emph{all} the numbers $z$ such that $z^n = w$? In the real numbers, the equation $x^2 = 4$ has two solutions and $x^2 = -4$ has none. In $\mathbb{C}$, the situation is beautifully complete: \textbf{every} non-zero complex number has \textbf{exactly} $n$ distinct $n$th roots, and they are arranged with perfect symmetry in the Argand diagram. This section solves the problem completely, then studies in detail the most important special case, the cube roots, which appear constantly at the GCE Advanced Level.

\courssub{The nth Roots of a Complex Number}

\textbf{Setting up the problem.} Let $w$ be a non-zero complex number. We seek all $z \in \mathbb{C}$ such that $z^n = w$. Write both numbers in exponential form:
\[ w = R e^{i\varphi} \quad (R > 0), \qquad z = r e^{i\theta}. \]
Then $z^n = w$ becomes $r^n e^{in\theta} = R e^{i\varphi}$. Two complex numbers in exponential form are equal if and only if their moduli are equal and their arguments differ by a multiple of $2\pi$. Hence
\[ r^n = R \quad\text{and}\quad n\theta = \varphi + 2k\pi \quad (k \in \mathbb{Z}). \]
Since $r$ is a positive real number, $r = R^{1/n}$ (the ordinary positive real $n$th root). The arguments are $\theta = \dfrac{\varphi + 2k\pi}{n}$. Taking $k = 0, 1, \dots, n-1$ gives $n$ distinct values; any other integer $k$ repeats one of these, because increasing $k$ by $n$ adds $2\pi$ to $\theta$.

\begin{propriete}[The $n$th roots of a complex number]
Every non-zero complex number $w = R e^{i\varphi}$ has exactly $n$ distinct $n$th roots, given by
\[ z_k = R^{1/n}\, e^{\,i(\varphi + 2k\pi)/n}, \qquad k = 0, 1, \dots, n-1. \]
The derivation above is examinable: you may be asked to \emph{show} that these are the only solutions.
\end{propriete}

\begin{attention}[Do not forget the $2k\pi$ term]
The most frequent error at the GCE is to write only $\theta = \varphi/n$, which yields just \emph{one} root. The equation $z^n = w$ has $n$ solutions: the term $2k\pi$ in $n\theta = \varphi + 2k\pi$ is what produces them all.
\end{attention}

\begin{attention}[The modulus of the roots]
$R^{1/n}$ is the \emph{positive real} $n$th root of the modulus $R$ (for example $8^{1/3} = 2$). It is an ordinary real number, not a complex root.
\end{attention}

\begin{propriete}[Geometric property: regular polygon]
The $n$ roots $z_0, z_1, \dots, z_{n-1}$ all have the same modulus $R^{1/n}$, and consecutive roots differ in argument by $\dfrac{2\pi}{n}$. Therefore their images in the Argand diagram lie on the circle of centre $O$ and radius $R^{1/n}$, at the vertices of a \cle{regular $n$-gon}.
\end{propriete}

\begin{methode}[Solving $z^n = w$]
\begin{enumerate}
\item Write $w$ in exponential form $w = R e^{i\varphi}$ (modulus, then argument).
\item Compute the common modulus $R^{1/n}$ of the roots.
\item Write the $n$ arguments $\theta_k = \dfrac{\varphi + 2k\pi}{n}$ for $k = 0, 1, \dots, n-1$.
\item State the roots $z_k = R^{1/n} e^{i\theta_k}$, then convert each to the form requested (often $a + ib$).
\end{enumerate}
\end{methode}

\begin{exoresolu}[Fourth roots of $-16$]
Solve $z^4 = -16$, giving the roots in the form $a + ib$.
\tcblower
\textbf{Solution.} $-16 = 16 e^{i\pi}$, so $R = 16$, $\varphi = \pi$, and $R^{1/4} = 16^{1/4} = 2$. The arguments are
\[ \theta_k = \frac{\pi + 2k\pi}{4}, \quad k = 0,1,2,3 \;:\quad \frac{\pi}{4},\ \frac{3\pi}{4},\ \frac{5\pi}{4},\ \frac{7\pi}{4}. \]
Hence
\begin{align*}
z_0 &= 2e^{i\pi/4} = \sqrt{2} + i\sqrt{2}, &
z_1 &= 2e^{3i\pi/4} = -\sqrt{2} + i\sqrt{2},\\
z_2 &= 2e^{5i\pi/4} = -\sqrt{2} - i\sqrt{2}, &
z_3 &= 2e^{7i\pi/4} = \sqrt{2} - i\sqrt{2}.
\end{align*}
Check: $(\sqrt{2}+i\sqrt{2})^2 = 4i$, so $(\sqrt{2}+i\sqrt{2})^4 = -16$. $\checkmark$
\end{exoresolu}

\courssub{Roots of Unity}

\begin{definition}[Roots of unity]
The solutions of $z^n = 1$ are called the \cle{$n$th roots of unity}. Since $1 = e^{i\cdot 0}$, they are
\[ \omega_k = e^{2ik\pi/n}, \qquad k = 0, 1, \dots, n-1, \]
the vertices of a regular $n$-gon inscribed in the unit circle, with one vertex at $1$.
\end{definition}

\begin{popfigure}
\begin{center}
\begin{tikzpicture}[scale=2.6]
\draw[popline,->] (-1.5,0) -- (1.5,0) node[below left, popmuted] {\scriptsize Re};
\draw[popline,->] (0,-1.5) -- (0,1.5) node[below left, popmuted] {\scriptsize Im};
\draw[popblue, thick] (0,0) circle (1);
\coordinate (A0) at (1,0);
\coordinate (A1) at (0.309,0.951);
\coordinate (A2) at (-0.809,0.588);
\coordinate (A3) at (-0.809,-0.588);
\coordinate (A4) at (0.309,-0.951);
\draw[poporange, thick] (A0)--(A1)--(A2)--(A3)--(A4)--cycle;
\foreach \p/\lab/\pos in {A0/{$\omega_0=1$}/right, A1/{$\omega_1=e^{2i\pi/5}$}/{above right}, A2/{$\omega_2$}/{above left}, A3/{$\omega_3$}/{below left}, A4/{$\omega_4$}/{below right}}
  \fill[popdark] (\p) circle (0.025) node[\pos, popdark] {\scriptsize \lab};
\foreach \p in {A1,A2,A3,A4} \draw[popmuted, thin] (0,0)--(\p);
\draw[popgreen, thin] (0,0)--(A0);
\draw[popgold, thick, ->] (0.45,0) arc (0:72:0.45);
\node[popgold] at (0.62,0.28) {\scriptsize $\frac{2\pi}{5}$};
\fill (0,0) circle (0.02) node[below left, popmuted] {\scriptsize $O$};
\end{tikzpicture}
\end{center}
\legende{The five 5th roots of unity: vertices of a regular pentagon on the unit circle.}
\end{popfigure}

\begin{propriete}[Sum of the roots of unity]
The sum of the $n$th roots of unity is zero:
\[ 1 + \omega + \omega^2 + \dots + \omega^{n-1} = 0, \qquad \text{where } \omega = e^{2i\pi/n}. \]
\emph{Proof} (examinable): this is a geometric progression with ratio $\omega \neq 1$, so the sum is $\dfrac{1 - \omega^n}{1 - \omega} = \dfrac{1-1}{1-\omega} = 0$, since $\omega^n = 1$.
\end{propriete}

Geometrically this says the ``centre of mass'' of the regular polygon is the origin — a satisfying confirmation of the symmetry.

\courssub{The Cube Roots of a Complex Number}

The case $n = 3$ deserves a detailed study, both because it is the most frequent at the GCE and because it introduces the celebrated number $j$. Solving $z^3 = w$ with $w = Re^{i\varphi}$ gives three roots
\[ z_k = R^{1/3} e^{\,i(\varphi + 2k\pi)/3}, \qquad k = 0, 1, 2, \]
spaced by $\dfrac{2\pi}{3}$ on the circle of radius $R^{1/3}$: an equilateral triangle.

\begin{definition}[The number $j$]
The cube roots of unity are the solutions of $z^3 = 1$:
\[ 1, \qquad j = e^{2i\pi/3} = -\frac{1}{2} + i\frac{\sqrt{3}}{2}, \qquad j^2 = e^{4i\pi/3} = -\frac{1}{2} - i\frac{\sqrt{3}}{2} = \bar{j}. \]
They satisfy the fundamental identities
\[ j^3 = 1, \qquad 1 + j + j^2 = 0, \qquad j^2 = \bar{j} = \frac{1}{j}. \]
\end{definition}

\begin{popfigure}
\begin{center}
\begin{tikzpicture}[scale=2.6]
\draw[popline,->] (-1.6,0) -- (1.6,0) node[below left, popmuted] {\scriptsize Re};
\draw[popline,->] (0,-1.5) -- (0,1.5) node[below left, popmuted] {\scriptsize Im};
\draw[popblue, thick] (0,0) circle (1);
\coordinate (B0) at (1,0);
\coordinate (B1) at (-0.5,0.866);
\coordinate (B2) at (-0.5,-0.866);
\draw[poporange, thick] (B0)--(B1)--(B2)--cycle;
\draw[popmuted, thin] (0,0)--(B0); \draw[popmuted, thin] (0,0)--(B1); \draw[popmuted, thin] (0,0)--(B2);
\fill[popdark] (B0) circle (0.025) node[right, popdark] {\scriptsize $1$};
\fill[popdark] (B1) circle (0.025) node[above left, popdark] {\scriptsize $j=-\frac12+i\frac{\sqrt3}{2}$};
\fill[popdark] (B2) circle (0.025) node[below left, popdark] {\scriptsize $j^2=\bar j$};
\draw[popgold, thick, ->] (0.5,0) arc (0:120:0.5);
\node[popgold] at (0.1,0.62) {\scriptsize $\frac{2\pi}{3}$};
\fill (0,0) circle (0.02) node[below right, popmuted] {\scriptsize $O$};
\end{tikzpicture}
\end{center}
\legende{The cube roots of unity $1, j, j^2$: an equilateral triangle on the unit circle.}
\end{popfigure}

\begin{exemplebox}[Factorising with $j$]
Since $j^3 = 1$, we have $a^3 + b^3 = (a+b)(a + jb)(a + j^2 b)$. Indeed, expanding $(a+jb)(a+j^2b) = a^2 + (j + j^2)ab + j^3 b^2 = a^2 - ab + b^2$ (using $1 + j + j^2 = 0$ and $j^3 = 1$), and $(a+b)(a^2 - ab + b^2) = a^3 + b^3$. Similarly $a^3 - b^3 = (a - b)(a - jb)(a - j^2 b)$.
\end{exemplebox}

\begin{exoresolu}[Solving $z^3 = -8$ and $z^3 = i$]
(a) Solve $z^3 = -8$. \quad (b) Solve $z^3 = i$, expressing the roots in the form $a + ib$.
\tcblower
\textbf{Solution.} (a) $-8 = 8e^{i\pi}$, so $R^{1/3} = 2$ and $\theta_k = \dfrac{\pi + 2k\pi}{3}$, giving $\theta = \dfrac{\pi}{3}, \pi, \dfrac{5\pi}{3}$. Hence
\[ z_0 = 2e^{i\pi/3} = 1 + i\sqrt{3}, \qquad z_1 = 2e^{i\pi} = -2, \qquad z_2 = 2e^{5i\pi/3} = 1 - i\sqrt{3}. \]
Quick alternative: $z^3 = -8 \iff \left(\tfrac{z}{-2}\right)^3 = 1$, so $z = -2,\ -2j,\ -2j^2$, which are the same three numbers.

(b) $i = e^{i\pi/2}$, so $R = 1$ and $\theta_k = \dfrac{\pi/2 + 2k\pi}{3} = \dfrac{\pi}{6}, \dfrac{5\pi}{6}, \dfrac{3\pi}{2}$. Therefore
\[ z_0 = e^{i\pi/6} = \frac{\sqrt{3}}{2} + \frac{1}{2}i, \qquad z_1 = e^{5i\pi/6} = -\frac{\sqrt{3}}{2} + \frac{1}{2}i, \qquad z_2 = e^{3i\pi/2} = -i. \]
\end{exoresolu}

\begin{popfigure}
\begin{center}
\begin{tikzpicture}[scale=1.3]
\draw[popline,->] (-3,0) -- (3,0) node[below left, popmuted] {\scriptsize Re};
\draw[popline,->] (0,-3) -- (0,3) node[below left, popmuted] {\scriptsize Im};
\draw[popblue, thick] (0,0) circle (2);
\coordinate (C0) at (1.732,1);
\coordinate (C1) at (-1.732,1);
\coordinate (C2) at (0,-2);
\draw[poporange, thick] (C0)--(C1)--(C2)--cycle;
\draw[popmuted, thin] (0,0)--(C0); \draw[popmuted, thin] (0,0)--(C1); \draw[popmuted, thin] (0,0)--(C2);
\fill[popdark] (C0) circle (0.045) node[above right, popdark] {\scriptsize $\sqrt3 + i$};
\fill[popdark] (C1) circle (0.045) node[above left, popdark] {\scriptsize $-\sqrt3 + i$};
\fill[popdark] (C2) circle (0.045) node[below, popdark] {\scriptsize $-2i$};
\draw[popgold, thick, ->] (0.8,0.5) arc (30:150:0.94);
\node[popgold] at (0,1.3) {\scriptsize $\frac{2\pi}{3}$};
\fill (0,0) circle (0.04) node[below right, popmuted] {\scriptsize $O$};
\end{tikzpicture}
\end{center}
\legende{The three cube roots of $8i$: $\sqrt{3}+i$, $-\sqrt{3}+i$, $-2i$, on the circle of radius $2$.}
\end{popfigure}

\begin{exemplebox}[Cube roots of $8i$ and $-8i$]
For $8i = 8e^{i\pi/2}$: $R^{1/3} = 2$, $\theta_k = \dfrac{\pi/2 + 2k\pi}{3} = \dfrac{\pi}{6}, \dfrac{5\pi}{6}, \dfrac{3\pi}{2}$. The roots are
\[ 2e^{i\pi/6} = \sqrt{3} + i, \qquad 2e^{5i\pi/6} = -\sqrt{3} + i, \qquad 2e^{3i\pi/2} = -2i. \]
For $-8i = 8e^{-i\pi/2}$: $\theta_k = \dfrac{-\pi/2 + 2k\pi}{3} = -\dfrac{\pi}{6}, \dfrac{\pi}{2}, \dfrac{7\pi}{6}$. The roots are
\[ 2e^{-i\pi/6} = \sqrt{3} - i, \qquad 2e^{i\pi/2} = 2i, \qquad 2e^{7i\pi/6} = -\sqrt{3} - i. \]
Both families form an equilateral triangle of radius $2$; only the rotation differs. The figure above illustrates the roots of $8i$.
\end{exemplebox}

\courssub{The Complete Picture and Examination Advice}

The programme on roots is now complete. The square-root method of Lesson 32 (solving $z^2 = w$ algebraically) is exactly the case $n = 2$ of the general formula: the two square roots of $Re^{i\varphi}$ are $\pm\sqrt{R}\,e^{i\varphi/2}$, diametrically opposite on the circle of radius $\sqrt{R}$. Whatever the degree, the strategy is identical: exponential form, root formula, conversion.

\begin{attention}[GCE exam extra: convert to $a + ib$]
The GCE frequently asks, after the roots have been found in exponential form, to ``\emph{express the roots in the form $a + ib$}''. You must then evaluate the cosines and sines exactly (standard angles $\frac{\pi}{6}, \frac{\pi}{4}, \frac{\pi}{3}, \frac{2\pi}{3}$, etc.). Leaving answers as $2e^{i\pi/6}$ when $a + ib$ is requested costs marks.
\end{attention}

\begin{aretenir}[Key points]
\begin{itemize}
\item $w = Re^{i\varphi} \neq 0$ has exactly $n$ $n$th roots: $z_k = R^{1/n} e^{i(\varphi + 2k\pi)/n}$, $k = 0, \dots, n-1$.
\item They form a regular $n$-gon on the circle of radius $R^{1/n}$; consecutive roots differ by $\frac{2\pi}{n}$.
\item The $n$th roots of unity sum to $0$.
\item $j = e^{2i\pi/3} = -\tfrac12 + i\tfrac{\sqrt3}{2}$ satisfies $j^3 = 1$, $1 + j + j^2 = 0$, $j^2 = \bar j$.
\item $z^3 = a^3$ has solutions $a,\ aj,\ aj^2$ — a powerful shortcut.
\end{itemize}
\end{aretenir}

\begin{exercice}[Practice]
\begin{enumerate}
\item Solve $z^5 = 1$ and verify that the sum of the roots is $0$.
\item Solve $z^3 = 27$, then $z^3 = -27i$, giving all roots in the form $a + ib$.
\item Show that $(1 + j)^2 = j$ and deduce $(1+j)^6$.
\item Find the fourth roots of $-8 - 8\sqrt{3}\,i$ in exponential form, then in the form $a + ib$.
\end{enumerate}
\end{exercice}

\begin{corrige}[Exercise 1]
\textbf{1.} $z_k = e^{2ik\pi/5}$, $k = 0,\dots,4$. Sum $= \dfrac{1 - \omega^5}{1 - \omega} = 0$ with $\omega = e^{2i\pi/5}$.
\textbf{2.} $z^3 = 27$: $z = 3, 3j, 3j^2$, i.e.\ $3$, $-\tfrac32 + \tfrac{3\sqrt3}{2}i$, $-\tfrac32 - \tfrac{3\sqrt3}{2}i$. For $z^3 = -27i = 27e^{-i\pi/2}$: modulus $3$, arguments $-\tfrac{\pi}{6}, \tfrac{\pi}{2}, \tfrac{7\pi}{6}$: roots $\tfrac{3\sqrt3}{2} - \tfrac32 i$, $3i$, $-\tfrac{3\sqrt3}{2} - \tfrac32 i$.
\textbf{3.} $(1+j)^2 = 1 + 2j + j^2 = (1 + j + j^2) + j = 0 + j = j$ (since $1+j+j^2=0$). Hence $(1+j)^6 = \big((1+j)^2\big)^3 = j^3 = 1$.
\textbf{4.} $-8 - 8\sqrt3 i = 16 e^{-2i\pi/3}$ (modulus $16$, argument $-\tfrac{2\pi}{3}$). Roots: $2e^{i(-2\pi/3 + 2k\pi)/4}$, i.e.\ $2e^{-i\pi/6} = \sqrt3 - i$, $2e^{i\pi/3} = 1 + i\sqrt3$, $2e^{5i\pi/6} = -\sqrt3 + i$, $2e^{4i\pi/3} = -1 - i\sqrt3$.
\end{corrige}

\coursec{Loci in the Complex Plane and Synthesis}

In the previous section we learnt how to find all the $n$th roots of a complex number, and we saw a beautiful fact: those roots are always the vertices of a regular polygon inscribed in a circle. In other words, the equation $z^n = w$ describes a \emph{set of points} in the plane with a precise geometric shape. This final section generalises that idea. A condition on a complex number $z$ --- an equation or an inequality involving $|z - a|$ or $\arg(z - a)$ --- picks out a curve or a region of the Argand diagram. Such a set of points is called a \cle{locus} (plural: \emph{loci}). Loci questions are among the most frequent and most rewarding questions in the GCE Advanced Level paper, because they test simultaneously your algebra, your geometry and your understanding of the modulus and the argument.

\courssub{Loci Defined by the Modulus}

\textbf{The circle.} Let $a$ be a fixed complex number, represented by the point $A$, and let $z$ be a variable complex number, represented by the point $M$. The number $z - a$ is represented by the vector $\overrightarrow{AM}$, so $|z - a|$ is nothing other than the \emph{distance} $AM$. The condition $|z - a| = r$ therefore says: ``$M$ is at distance $r$ from $A$'' --- which is exactly the definition of a circle.

\begin{propriete}[Circle and perpendicular bisector]
Let $a, b \in \mathbb{C}$ with images $A, B$, and let $r > 0$. Then:
\begin{itemize}
\item $|z - a| = r \iff M(z)$ lies on the \cle{circle} of centre $A(a)$ and radius $r$;
\item $|z - a| = |z - b| \iff M(z)$ lies on the \cle{perpendicular bisector} of the segment $[AB]$.
\end{itemize}
\end{propriete}

The second statement is the classical fact that the points equidistant from $A$ and $B$ form the perpendicular bisector of $[AB]$, since $|z - a| = AM$ and $|z - b| = BM$.

\begin{exemplebox}[Reading a locus directly]
The equation $|z - (1 + 2i)| = 3$ describes the circle of centre $(1, 2)$ and radius $3$. The equation $|z - 1| = |z + i|$, rewritten as $|z - 1| = |z - (-i)|$, describes the perpendicular bisector of the segment joining $A(1, 0)$ and $B(0, -1)$.
\end{exemplebox}

\begin{popfigure}
\begin{tikzpicture}[scale=0.85]
\draw[->, popmuted] (-3.5,0) -- (4.5,0) node[below left] {Re};
\draw[->, popmuted] (0,-3.5) -- (0,4.5) node[below left] {Im};
% circle |z-(1+2i)|=3
\draw[popblue, thick] (1,2) circle (3);
\fill[popblue] (1,2) circle (1.5pt) node[above right, popdark] {$1+2i$};
\draw[popblue, dashed] (1,2) -- (3.12,3.5) node[midway, above, popdark] {$3$};
% perpendicular bisector of A(1,0), B(0,-1)
\fill[poporange] (1,0) circle (1.5pt) node[above right, popdark] {$A(1)$};
\fill[poporange] (0,-1) circle (1.5pt) node[below right, popdark] {$B(-i)$};
\draw[poporange, dashed] (1,0) -- (0,-1);
% Corrected bisector line y = -x
\draw[poporange, thick] (-3,3) -- (3,-3) node[below right, popdark, align=center, text width=2.6cm] {bisector $|z-1|=|z+i|$};
\end{tikzpicture}
\legende{The circle $|z-(1+2i)|=3$ and the perpendicular bisector $|z-1|=|z+i|$.}
\end{popfigure}

\begin{methode}[From a locus to a Cartesian equation]
To obtain the Cartesian equation of a locus:
\begin{enumerate}
\item Replace $z$ by $x + iy$ with $x, y \in \mathbb{R}$.
\item Compute each modulus using $|u + iv| = \sqrt{u^2 + v^2}$.
\item Square both sides to eliminate the square roots.
\item Simplify; recognise the curve (circle, line, etc.).
\end{enumerate}
\end{methode}

\begin{exoresolu}[Cartesian equation of a bisector]
Find the Cartesian equation of the locus $|z - 1| = |z + i|$ and identify it.
\tcblower
Set $z = x + iy$. Then $z - 1 = (x-1) + iy$ and $z + i = x + i(y+1)$, so
\[ \sqrt{(x-1)^2 + y^2} = \sqrt{x^2 + (y+1)^2}. \]
Squaring: $(x-1)^2 + y^2 = x^2 + (y+1)^2$, that is
\[ x^2 - 2x + 1 + y^2 = x^2 + y^2 + 2y + 1 \iff -2x = 2y \iff y = -x. \]
The locus is the straight line $y = -x$, the perpendicular bisector of the segment joining $(1,0)$ and $(0,-1)$ --- exactly as the figure shows.
\end{exoresolu}

\courssub{Loci Defined by the Argument}

The number $\arg(z - a)$ is the angle that the vector $\overrightarrow{AM}$ makes with the positive real axis. Fixing this angle forces $M$ onto a ray starting at $A$.

\begin{propriete}[Half-line locus]
Let $a \in \mathbb{C}$ with image $A$, and $\theta \in \mathbb{R}$. Then
\[ \arg(z - a) = \theta \iff M(z) \text{ lies on the half-line from } A \text{ making angle } \theta \text{ with the positive real axis,} \]
the point $A$ itself being \emph{excluded}.
\end{propriete}

\begin{attention}[The origin of the ray is excluded]
The point $a$ is never part of the locus $\arg(z - a) = \theta$, because at $z = a$ we would need $\arg(0)$, which is undefined. Mark $A$ with an open (hollow) dot on your diagram. Similarly, in $\arg\!\left(\dfrac{z-a}{z-b}\right) = \theta$, both $a$ and $b$ are excluded.
\end{attention}

\textbf{The Thales locus.} A quotient of complex numbers has argument
\[ \arg\!\left(\frac{z - a}{z - b}\right) = \arg(z-a) - \arg(z-b) = \text{angle } (\overrightarrow{BM}, \overrightarrow{AM}) = \widehat{AMB}. \]
Requiring this angle to be $\pm \dfrac{\pi}{2}$ means the angle $\widehat{AMB}$ is a right angle. By the theorem of Thales, $M$ lies on the circle of diameter $[AB]$.

\begin{propriete}[{Circle of diameter $[AB]$}]
\[ \arg\!\left(\frac{z - a}{z - b}\right) = \pm\frac{\pi}{2} \pmod{\pi} \iff M(z) \text{ lies on the circle of diameter } [AB], \ A \text{ and } B \text{ excluded.} \]
With a single sign, $\arg\!\left(\dfrac{z-a}{z-b}\right) = \dfrac{\pi}{2}$ (principal value), only one \emph{semicircle} is obtained.
\end{propriete}

\begin{popfigure}
\begin{tikzpicture}[scale=1.1]
\draw[->, popmuted] (-2.2,0) -- (2.4,0) node[below left] {Re};
\draw[->, popmuted] (0,-1.6) -- (0,1.9) node[below left] {Im};
\draw[poppurple, very thick] (1,0) arc (0:180:1);
\draw[poppurple, dashed] (-1,0) arc (180:360:1);
\fill[white, draw=poppurple, thick] (1,0) circle (1.6pt);
\fill[white, draw=poppurple, thick] (-1,0) circle (1.6pt);
\node[below, popdark] at (1,-0.08) {$1$};
\node[below, popdark] at (-1,-0.08) {$-1$};
\fill[popblue] (0,1) circle (1.5pt) node[above, popdark] {$M=i$};
\draw[popblue] (0,1) -- (1,0);
\draw[popblue] (0,1) -- (-1,0);
% Corrected right angle mark at M(0,1)
\draw[popdark] (0.07,0.93) -- (0,0.86) -- (-0.07,0.93);
\end{tikzpicture}
\legende{The locus $\arg\!\left(\frac{z-1}{z+1}\right)=\frac{\pi}{2}$: the upper semicircle of diameter $[-1,1]$ (solid), endpoints excluded. The dashed semicircle corresponds to $-\frac{\pi}{2}$.}
\end{popfigure}

\begin{exemplebox}[Checking the Thales locus]
Take $z = i$. Then $\dfrac{z-1}{z+1} = \dfrac{-1+i}{1+i} = \dfrac{(-1+i)(1-i)}{(1+i)(1-i)} = \dfrac{2i}{2} = i$, whose argument is $\dfrac{\pi}{2}$. Indeed $i$ lies on the circle of diameter $[-1, 1]$, since $|i| = 1$.
\end{exemplebox}

\courssub{Inequalities and Regions of the Plane}

An inequality involving the modulus describes a \emph{region} rather than a curve:
\begin{itemize}
\item $|z - a| \leq r$: the closed \cle{disc} of centre $a$, radius $r$ (boundary circle included, drawn solid);
\item $|z - a| < r$: the open disc (boundary excluded, drawn dashed);
\item $r_1 \leq |z| \leq r_2$: an \cle{annulus} centred at the origin;
\item $\alpha \leq \arg z \leq \beta$: a \cle{sector} (wedge) from the origin.
\end{itemize}
Conditions can be combined: the intersection of an annulus and a sector is a ``piece of cake'' region.

\begin{attention}[Solid versus dashed boundaries]
For strict inequalities ($<$, $>$), draw the boundary \textbf{dashed}: it is not part of the locus. For non-strict inequalities ($\leq$, $\geq$), draw it \textbf{solid}. Examiners award marks for this convention.
\end{attention}

\begin{methode}[Sketching a region]
\begin{enumerate}
\item Sketch the boundary curve(s) first, using the equality cases.
\item Choose one convenient test point (e.g. $z = 0$ or $z = 1$) and check whether it satisfies the inequality.
\item Shade the side containing the test point if it works, the other side otherwise.
\item Apply the solid/dashed convention and mark excluded points with hollow dots.
\end{enumerate}
\end{methode}

\begin{popfigure}
\begin{tikzpicture}[scale=1.5]
\fill[popgreenL] (0:1) arc (0:60:1) -- (60:2) arc (60:0:2) -- cycle;
\draw[popgreen, thick] (0:1) arc (0:60:1);
\draw[popgreen, thick] (0:2) arc (0:60:2);
\draw[popgreen, thick] (0:1) -- (0:2);
\draw[popgreen, thick] (60:1) -- (60:2);
\draw[popgreen, dashed] (60:2) arc (60:360:2);
\draw[popgreen, dashed] (60:1) arc (60:360:1);
\draw[->, popmuted] (-2.4,0) -- (2.6,0) node[below left] {Re};
\draw[->, popmuted] (0,-2.3) -- (0,2.4) node[below left] {Im};
\draw[popmuted] (0,0) -- (60:2.3);
\node[popdark] at (30:0.55) {$\frac{\pi}{3}$};
\draw[popdark] (0.7,0) arc (0:60:0.7);
\node[popdark, right] at (1.05,0) {$1$};
\node[popdark, right] at (2.05,0) {$2$};
\end{tikzpicture}
\legende{The region $1 \leq |z| \leq 2$ and $0 \leq \arg z \leq \frac{\pi}{3}$ (shaded): intersection of an annulus and a sector.}
\end{popfigure}

\courssub{Synthesis: GCE-Style Problems}

Examination questions typically chain several ideas: a locus condition, a conversion to Cartesian form, and a computation involving powers, roots or De Moivre's theorem.

\begin{exoresolu}[GCE-style synthesis]
Let $z = x + iy$.
\begin{enumerate}
\item Show that the locus $|z - 4i| = 2|z - 1|$ is a circle, and find its centre and radius.
\item Find the complex number(s) on this locus with $\arg z = \dfrac{\pi}{2}$.
\end{enumerate}
\tcblower
\textbf{1.} $|x + i(y-4)| = 2|(x-1) + iy|$ gives, after squaring,
\[ x^2 + (y-4)^2 = 4\big[(x-1)^2 + y^2\big]. \]
Expanding: $x^2 + y^2 - 8y + 16 = 4x^2 - 8x + 4 + 4y^2$, so
\[ 3x^2 + 3y^2 - 8x + 8y - 12 = 0 \iff x^2 + y^2 - \tfrac{8}{3}x + \tfrac{8}{3}y - 4 = 0. \]
Completing the squares:
\[ \left(x - \tfrac{4}{3}\right)^2 + \left(y + \tfrac{4}{3}\right)^2 = 4 + \tfrac{16}{9} + \tfrac{16}{9} = \tfrac{68}{9}. \]
This is a circle of centre $\dfrac{4}{3} - \dfrac{4}{3}i$ and radius $\dfrac{\sqrt{68}}{3} = \dfrac{2\sqrt{17}}{3}$.

\textbf{2.} $\arg z = \dfrac{\pi}{2}$ means $z = iy$ with $y > 0$. Substituting $x = 0$:
\[ y^2 + \tfrac{8}{3}y - 4 = 0 \iff 3y^2 + 8y - 12 = 0 \iff y = \frac{-8 \pm \sqrt{64 + 144}}{6} = \frac{-8 \pm 4\sqrt{13}}{6}. \]
Only the positive root is acceptable: $z = \dfrac{-4 + 2\sqrt{13}}{3}\, i$.
\end{exoresolu}

\begin{exoresolu}[Locus and roots combined]
\begin{enumerate}
\item Sketch the locus $|z| = 1$ and mark on it the cube roots of $-8$.
\item Verify that each cube root of $-8$ lies on the circle $|z| = 2$, not $|z| = 1$.
\end{enumerate}
\tcblower
\textbf{1.} $|z| = 1$ is the unit circle centred at $O$.

\textbf{2.} Write $-8 = 8e^{i\pi}$. The cube roots are
\[ z_k = 2\, e^{i(\pi + 2k\pi)/3}, \quad k = 0, 1, 2, \]
that is $z_0 = 2e^{i\pi/3} = 1 + i\sqrt{3}$, $z_1 = 2e^{i\pi} = -2$, $z_2 = 2e^{i5\pi/3} = 1 - i\sqrt{3}$. Each has modulus $2$, so all three lie on the circle $|z| = 2$ --- they form an equilateral triangle inscribed in that circle, not on the unit circle. This illustrates the general principle: the $n$th roots of $w$ lie on the circle $|z| = \sqrt[n]{|w|}$.
\end{exoresolu}

\courssub{Chapter Revision Map: Choosing the Right Form}

Throughout this chapter you have met three faces of a complex number. Choosing the right one is half the battle in the examination.

\begin{center}
\begin{tabularx}{\linewidth}{|l|X|X|}
\hline
\rowcolor{popblueL} \textbf{Form} & \textbf{Best for} & \textbf{Typical tasks} \\
\hline
Cartesian $z = x + iy$ & Addition, subtraction, conjugates; converting loci to Cartesian equations & Simplifying expressions, solving equations with real/imaginary parts, locus equations \\
\hline
Polar $r(\cos\theta + i\sin\theta)$ & Multiplication, division, arguments, loci involving $\arg$ & Products and quotients, argument conditions, Thales-type loci \\
\hline
Exponential $re^{i\theta}$ & Powers, roots, De Moivre, trigonometric identities & Computing $z^n$, $n$th roots, expressing $\cos^n\theta$ in multiple angles and conversely \\
\hline
\end{tabularx}
\end{center}

\begin{aretenir}[Key points of the section]
\begin{itemize}
\item $|z - a| = r$: circle, centre $a$, radius $r$. $|z - a| = |z - b|$: perpendicular bisector of $[AB]$.
\item $\arg(z - a) = \theta$: half-line from $a$ ($a$ excluded). $\arg\!\left(\dfrac{z-a}{z-b}\right) = \pm\dfrac{\pi}{2}$: circle of diameter $[AB]$.
\item Inequalities give regions: discs, annuli, sectors; test one interior point; dashed boundary for strict inequalities.
\item Cartesian equation: set $z = x + iy$, compute moduli, square, simplify.
\item Powers and roots: switch to exponential form and use De Moivre.
\end{itemize}
\end{aretenir}

\begin{exercice}[Practice]
\begin{enumerate}
\item Find and sketch the locus $|z + 2 - i| = 3$; give its Cartesian equation.
\item Show that $|z - 3| = |z + 3i|$ is a line and find its equation.
\item Sketch the region defined by $|z| < 2$ and $\dfrac{\pi}{4} \leq \arg z \leq \dfrac{3\pi}{4}$.
\item Describe the locus $\arg\!\left(\dfrac{z - 2i}{z - 2}\right) = \dfrac{\pi}{2}$.
\item Find the points of the circle $|z - 1| = 2$ where $\arg(z - 1) = \dfrac{\pi}{3}$, giving them in Cartesian form.
\end{enumerate}
\end{exercice}

\begin{corrige}[Exercise 1]
\begin{enumerate}
\item Circle of centre $-2 + i$, radius $3$: $(x+2)^2 + (y-1)^2 = 9$.
\item Bisector of $(3,0)$ and $(0,-3)$: squaring gives $-6x = 6y$, i.e. $y = -x$.
\item Open disc of radius $2$ (dashed circle) intersected with the sector between the rays at $\frac{\pi}{4}$ and $\frac{3\pi}{4}$ (solid rays, origin excluded).
\item Semicircle of diameter joining $2i$ and $2$, endpoints excluded.
\item $z = 1 + 2e^{i\pi/3} = 1 + 2\left(\frac{1}{2} + i\frac{\sqrt{3}}{2}\right) = 2 + i\sqrt{3}$.
\end{enumerate}
\end{corrige}

\begin{savaistu}[Why loci matter beyond the classroom]
Locus thinking with complex numbers is the seed of \emph{conformal mapping}, a tool engineers use to study airflow around a wing or the electric field around a conductor: complicated regions of the plane are transformed into discs or half-planes where the equations become easy. The same mathematics governs the stability criteria used when designing control systems --- a modern echo of the humble circle $|z - a| = r$ you have just mastered.
\end{savaistu}

\newpage

\coursec{Integration activity}

\courssub{Aims of the activity}
This integration activity is designed to consolidate your mastery of complex numbers by placing you in a realistic engineering context. By the end of this activity, you should be able to:
\begin{itemize}
    \item Perform algebraic operations (addition, multiplication, division) on complex numbers in Cartesian form and interpret the results physically.
    \item Convert fluently between Cartesian, polar (trigonometric), and exponential forms.
    \item Apply De Moivre's theorem to compute powers of complex numbers efficiently.
    \item Interpret the modulus as a magnitude (impedance magnitude) and the argument as a phase shift in an AC circuit.
    \item Sketch and describe loci in the Argand diagram, linking geometric descriptions to algebraic equations.
    \item Communicate mathematical reasoning clearly, justifying each step as required in the GCE Advanced Level examination.
\end{itemize}

\courssub{Situation}
\textbf{Context:} You are a team of junior engineers attached to the \emph{Direction de la Distribution} at ENEO (Energy of Cameroon) in Yaoundé. Your unit is analysing a section of the low-voltage alternating current (AC) network supplying the Mfoundi district. In AC circuit theory, opposition to current is not a simple resistance but a \emph{complex impedance} $Z = R + iX$, where $R$ is the resistance (real part) and $X$ is the reactance (imaginary part). The modulus $|Z|$ gives the magnitude of the impedance (in ohms, $\Omega$), and the argument $\arg(Z)$ gives the phase difference between voltage and current.

Two series branches in the network have been modelled with the following impedances:
\[
Z_1 = 3 + 4i\ \Omega, \qquad Z_2 = 1 - 2i\ \Omega.
\]
The supply voltage across the combination is given in exponential form as $V = 220\, e^{i\pi/3}$ volts (RMS value).

Your team must prepare a technical note for the senior engineer, answering the questions below. All calculations must be shown step by step, and final answers given in the requested forms. Where an Argand diagram is required, draw it accurately to scale on graph paper (or using geometry software) and include it in your report.

\courssub{Tasks}
\begin{enumerate}
    \item \textbf{Series and parallel combinations.}
    \begin{enumerate}
        \item Compute the total impedance $Z_{\text{series}} = Z_1 + Z_2$ when the two branches are connected in series. Express your answer in the form $a+ib$.
        \item Compute the equivalent impedance $Z_{\text{parallel}} = \dfrac{Z_1 Z_2}{Z_1 + Z_2}$ when the two branches are connected in parallel. Express your answer in the form $a+ib$ (simplify the fraction completely).
    \end{enumerate}

    \item \textbf{Polar and exponential forms of $Z_1$.}
    \begin{enumerate}
        \item Calculate the modulus $|Z_1|$ and the principal argument $\arg(Z_1)$ (in radians, to three significant figures).
        \item Write $Z_1$ in polar (trigonometric) form: $r(\cos\theta + i\sin\theta)$.
        \item Write $Z_1$ in exponential form: $r e^{i\theta}$.
        \item Interpret physically: what do $|Z_1|$ and $\arg(Z_1)$ represent in the AC circuit?
    \end{enumerate}

    \item \textbf{Current calculation using exponential form.}
    Given the supply voltage $V = 220\, e^{i\pi/3}$ V across $Z_1$ alone, use the exponential forms to compute the current $I = \dfrac{V}{Z_1}$. Give your answer in exponential form $I = I_0 e^{i\phi}$, stating the magnitude $I_0$ (in amperes) and the phase $\phi$ (in radians, to three significant figures). Explain why the exponential form makes this division particularly simple.

    \item \textbf{De Moivre's theorem and $Z_1^3$.}
    \begin{enumerate}
        \item Use De Moivre's theorem to compute $Z_1^3$ directly from the polar/exponential form of $Z_1$. Give the result in both polar and Cartesian ($a+ib$) forms.
        \item If the same voltage $V$ were applied across a hypothetical impedance $Z_1^3$, how would the magnitude of the current compare to that in Task 3? Explain without recalculating the current.
    \end{enumerate}

    \item \textbf{Locus on the Argand diagram.}
    \begin{enumerate}
        \item On an Argand diagram, sketch the locus of all points representing complex impedances $Z$ such that $|Z| = |Z_1|$. Use a scale of 1 cm per ohm. Clearly mark the positions of $Z_1$, $Z_2$, $Z_{\text{series}}$, and $Z_{\text{parallel}}$ on the same diagram.
        \item Describe the locus in words and give its Cartesian equation.
        \item Which of the four impedances ($Z_1$, $Z_2$, $Z_{\text{series}}$, $Z_{\text{parallel}}$) lie exactly on this locus? Justify your answer numerically.
    \end{enumerate}
\end{enumerate}

\courssub{Model answer}

\begin{exoresolu}[Task 1: Series and parallel combinations]
\textbf{1(a) Series impedance.} In series, impedances add. Add the real parts together and the imaginary parts together:
\[
Z_{\text{series}} = Z_1 + Z_2 = (3+4i) + (1-2i) = (3+1) + i(4-2) = 4 + 2i\ \Omega.
\]

\textbf{1(b) Parallel impedance.}
First compute the product $Z_1 Z_2$, expanding term by term:
\[
Z_1 Z_2 = (3+4i)(1-2i) = 3(1) + 3(-2i) + 4i(1) + 4i(-2i) = 3 - 6i + 4i - 8i^2.
\]
Since $i^2 = -1$, we have $-8i^2 = +8$. Hence
\[
Z_1 Z_2 = (3+8) + i(-6+4) = 11 - 2i.
\]
From part (a), $Z_1 + Z_2 = 4 + 2i$. Therefore
\[
Z_{\text{parallel}} = \frac{Z_1 Z_2}{Z_1 + Z_2} = \frac{11 - 2i}{4 + 2i}.
\]
To write this in the form $a+ib$, multiply numerator and denominator by the complex conjugate of the denominator, $\overline{4+2i} = 4 - 2i$:
\[
Z_{\text{parallel}} = \frac{(11 - 2i)(4 - 2i)}{(4 + 2i)(4 - 2i)}.
\]
\textbf{Denominator} (product of conjugates, giving a real number):
\[
(4+2i)(4-2i) = 4^2 + 2^2 = 16 + 4 = 20.
\]
\textbf{Numerator:}
\[
(11-2i)(4-2i) = 44 - 22i - 8i + 4i^2 = 44 - 30i - 4 = 40 - 30i.
\]
Thus
\[
Z_{\text{parallel}} = \frac{40 - 30i}{20} = \frac{40}{20} - \frac{30}{20}i = 2 - 1.5i\ \Omega.
\]
\end{exoresolu}

\begin{exoresolu}[Task 2: Polar and exponential forms of $Z_1$]
\textbf{2(a) Modulus and argument.}
\[
|Z_1| = \sqrt{3^2 + 4^2} = \sqrt{9+16} = \sqrt{25} = 5\ \Omega.
\]
Since $Z_1 = 3+4i$ lies in the first quadrant (both real and imaginary parts positive), the principal argument is
\[
\arg(Z_1) = \arctan\left(\frac{4}{3}\right) \approx 0.927\ \text{rad} \quad (\text{to three significant figures}).
\]

\textbf{2(b) Polar (trigonometric) form.}
\[
Z_1 = 5\left(\cos 0.927 + i\sin 0.927\right)\ \Omega.
\]
(Exact version: since $\cos\theta = \dfrac{3}{5}$ and $\sin\theta = \dfrac{4}{5}$, we may write $Z_1 = 5\left(\dfrac{3}{5} + i\,\dfrac{4}{5}\right)$, which is of course just $3+4i$ again.)

\textbf{2(c) Exponential form.}
\[
Z_1 = 5\, e^{0.927\, i}\ \Omega.
\]

\textbf{2(d) Physical interpretation.}
\begin{itemize}
    \item The modulus $|Z_1| = 5\ \Omega$ is the \emph{magnitude of the impedance}. It is the ratio of the RMS voltage across the branch to the RMS current through it: $|Z_1| = V_{\text{rms}}/I_{\text{rms}}$. It quantifies the total opposition to the flow of alternating current.
    \item The argument $\arg(Z_1) \approx 0.927\ \text{rad}$ (about $53.1^\circ$) is the \emph{phase angle} by which the voltage \emph{leads} the current in this branch. Because the reactance is positive ($X = 4\ \Omega$, inductive), the voltage reaches its maximum before the current does.
\end{itemize}
\end{exoresolu}

\begin{exoresolu}[Task 3: Current calculation using exponential form]
We have $V = 220\, e^{i\pi/3}$ V and $Z_1 = 5\, e^{0.927\, i}\ \Omega$. Using the quotient rule for exponential forms (divide the moduli, subtract the arguments):
\[
I = \frac{V}{Z_1} = \frac{220\, e^{i\pi/3}}{5\, e^{0.927\, i}} = \frac{220}{5}\, e^{i(\pi/3 - 0.927)}.
\]
\textbf{Magnitude:} $I_0 = \dfrac{220}{5} = 44$ A.

\textbf{Phase:} $\dfrac{\pi}{3} \approx 1.0472$ rad, so
\[
\phi = 1.0472 - 0.9273 = 0.1199 \approx 0.120\ \text{rad} \quad (\text{to three significant figures}).
\]
Hence
\[
I = 44\, e^{0.120\, i}\ \text{A}.
\]

\textbf{Why the exponential form simplifies the division:} in exponential form, division reduces to dividing the moduli and subtracting the arguments --- two ordinary real-number operations. This avoids the algebraic manipulation of multiplying numerator and denominator by the complex conjugate, which is required in Cartesian form. It is especially efficient when several impedances are combined, or when powers and roots are needed (as in Task 4).
\end{exoresolu}

\begin{exoresolu}[Task 4: De Moivre's theorem and $Z_1^3$]
\textbf{4(a) Computing $Z_1^3$.}
De Moivre's theorem states that for $n \in \mathbb{Z}$,
\[
\left(r e^{i\theta}\right)^n = r^n e^{in\theta}.
\]
With $Z_1 = 5\, e^{0.9273\, i}$ (keeping extra precision in the angle to avoid rounding error):
\[
Z_1^3 = 5^3\, e^{i(3 \times 0.9273)} = 125\, e^{2.782\, i}\ \Omega^3.
\]
Since $3 \times 0.9273 = 2.7819$ rad lies in $(-\pi, \pi]$, this is already the principal argument. The polar form is
\[
Z_1^3 = 125\left(\cos 2.782 + i\sin 2.782\right).
\]
Converting to Cartesian form with $\cos 2.782 \approx -0.9365$ and $\sin 2.782 \approx 0.3508$:
\[
Z_1^3 \approx 125(-0.9365 + 0.3508\, i) \approx -117 + 43.8\, i\ \Omega^3.
\]
\textbf{Exact check by direct multiplication} (recommended in the examination when the Cartesian form is simple):
\[
Z_1^2 = (3+4i)^2 = 9 + 24i + 16i^2 = -7 + 24i,
\]
\[
Z_1^3 = (-7+24i)(3+4i) = -21 - 28i + 72i + 96i^2 = -21 + 44i - 96 = -117 + 44i.
\]
So the exact answer is
\[
Z_1^3 = -117 + 44i\ \Omega^3 \quad \text{(Cartesian)}, \qquad Z_1^3 = 125\, e^{2.782\, i}\ \Omega^3 \quad \text{(exponential)}.
\]
(The small discrepancy between $44i$ and $43.8i$ comes from rounding the argument to 0.927; the exact Cartesian computation confirms $-117 + 44i$.)

\textbf{4(b) Current magnitude comparison.}
By De Moivre's theorem, $|Z_1^3| = |Z_1|^3 = 5^3 = 125\ \Omega$. With the same voltage applied,
\[
I' = \frac{|V|}{|Z_1^3|} = \frac{220}{125} = 1.76\ \text{A}.
\]
In Task 3 the current magnitude was $I_0 = 44$ A. Since the impedance magnitude has been multiplied by $\dfrac{|Z_1^3|}{|Z_1|} = \dfrac{125}{5} = 25$, the current magnitude is divided by 25: indeed $\dfrac{44}{25} = 1.76$ A. In general, the current magnitude scales as $\dfrac{1}{|Z_1|^n}$ when the impedance is raised to the power $n$.
\end{exoresolu}

\begin{exoresolu}[Task 5: Locus on the Argand diagram]
\textbf{5(a) Sketch.} Draw axes with the horizontal axis = Real part (resistance $R$) and the vertical axis = Imaginary part (reactance $X$), using a scale of 1 cm per ohm. The locus $|Z| = |Z_1| = 5$ is the circle centred at the origin with radius 5. Plot the four points:
\begin{align*}
Z_1 &= 3 + 4i \quad \rightarrow \quad (3, 4),\\
Z_2 &= 1 - 2i \quad \rightarrow \quad (1, -2),\\
Z_{\text{series}} &= 4 + 2i \quad \rightarrow \quad (4, 2),\\
Z_{\text{parallel}} &= 2 - 1.5i \quad \rightarrow \quad (2, -1.5).
\end{align*}

\begin{popfigure}
\begin{center}
\begin{tikzpicture}[scale=0.85]
    \draw[popline, very thin, step=1] (-5.5,-5.5) grid (5.5,5.5);
    \draw[->, thick] (-5.8,0) -- (5.8,0) node[below left] {$\mathrm{Re}$};
    \draw[->, thick] (0,-5.8) -- (0,5.8) node[below left] {$\mathrm{Im}$};
    \draw[popblue, thick] (0,0) circle (5);
    \draw[popblue, dashed] (0,0) -- (3,4);
    \fill[poporange] (3,4) circle (2.2pt) node[above right, popdark] {$Z_1=3+4i$};
    \fill[poppurple] (1,-2) circle (2.2pt) node[below right, popdark] {$Z_2=1-2i$};
    \fill[popgreen] (4,2) circle (2.2pt) node[above right, popdark] {$Z_{\text{series}}=4+2i$};
    \fill[poppink] (2,-1.5) circle (2.2pt) node[below right, popdark] {$Z_{\text{parallel}}=2-1.5i$};
    \node[popblue] at (-3.2,4.2) {$|Z|=5$};
\end{tikzpicture}
\end{center}
\legende{Argand diagram: the locus $|Z|=5$ (circle, centre $O$, radius 5) and the four impedances. Only $Z_1$ lies on the circle.}
\end{popfigure}

\textbf{5(b) Description and equation.}
The locus is the set of all points in the complex plane at a fixed distance 5 from the origin: a \textbf{circle centred at the origin with radius 5}. Writing $Z = x + iy$, the condition $|Z| = 5$ becomes $\sqrt{x^2+y^2} = 5$, i.e.
\[
x^2 + y^2 = 25.
\]

\textbf{5(c) Which impedances lie on the locus?}
A point lies on the locus if and only if its modulus equals 5. Computing each modulus:
\begin{align*}
|Z_1| &= \sqrt{3^2+4^2} = \sqrt{25} = 5 \quad \checkmark \ \text{(on the locus)}\\
|Z_2| &= \sqrt{1^2+(-2)^2} = \sqrt{5} \approx 2.24 \neq 5\\
|Z_{\text{series}}| &= \sqrt{4^2+2^2} = \sqrt{20} \approx 4.47 \neq 5\\
|Z_{\text{parallel}}| &= \sqrt{2^2+(-1.5)^2} = \sqrt{6.25} = 2.5 \neq 5
\end{align*}
Only $Z_1$ lies exactly on the locus; the other three impedances lie strictly inside the circle since their moduli are all less than 5.
\end{exoresolu}

\begin{attention}[Common errors to avoid in this type of problem]
\begin{itemize}
    \item Forgetting that $i^2 = -1$ when expanding products: a frequent sign error in $Z_1 Z_2$ and in $Z_1^3$.
    \item Dividing by a complex number without multiplying by the conjugate: $\dfrac{11-2i}{4+2i}$ is \emph{not} $\dfrac{11}{4} - \dfrac{2i}{2i}$.
    \item Rounding the argument too early: using $\theta = 0.93$ instead of $0.9273$ in Task 4 shifts the Cartesian result noticeably. Keep at least 4 significant figures in intermediate steps and round only at the end.
    \item Confusing the locus $|Z| = 5$ (a circle) with loci such as $|Z - Z_1| = |Z - Z_2|$ (a perpendicular bisector) or $\arg(Z - Z_1) = \theta$ (a half-line). Read the condition carefully before sketching.
    \item Giving the argument in degrees when the question asks for radians (or vice versa). State the unit explicitly.
\end{itemize}
\end{attention}

\begin{aretenir}[Key synthesis points for the examination]
\begin{itemize}
    \item \textbf{Cartesian form} ($a+ib$) is best for addition and subtraction (series/parallel combinations).
    \item \textbf{Polar/Exponential form} ($r e^{i\theta}$) is best for multiplication, division, powers and roots (De Moivre's theorem): multiply/divide the moduli, add/subtract the arguments.
    \item \textbf{Modulus} = magnitude (impedance magnitude, voltage/current ratio). \textbf{Argument} = phase shift (voltage leads current if $\arg(Z) > 0$).
    \item \textbf{Locus $|Z| = k$} is the circle centred at the origin with radius $k$; its Cartesian equation is $x^2 + y^2 = k^2$. To test whether a given complex number lies on a locus, compute its modulus (or substitute into the Cartesian equation).
    \item In GCE answers: show all steps, give exact values where possible (e.g. $\arg Z_1 = \arctan\frac{4}{3}$), and round only at the final step (usually to 3 significant figures for angles).
\end{itemize}
\end{aretenir}

\newpage

\coursec{Methods and worked examples}

\courssub{Skill 1: Performing Algebraic Operations in $\mathbb{C}$}

Every complex number can be written uniquely in the \cle{algebraic form} $z = a + ib$, where $a = \mathrm{Re}(z)$ is the \cle{real part} and $b = \mathrm{Im}(z)$ the \cle{imaginary part} (note: $b$ is a \emph{real} number). All the usual rules of algebra (commutativity, associativity, distributivity, remarkable identities) remain valid in $\mathbb{C}$; the only new ingredient is the fundamental relation $i^2 = -1$.

\begin{methode}[Adding, Subtracting and Multiplying Complex Numbers]
To compute with complex numbers written in \cle{algebraic form} $z = a + ib$ ($a, b \in \mathbb{R}$):
\begin{enumerate}
\item Treat $i$ as an ordinary letter and expand using the usual rules of algebra (distributivity, identities such as $(a+b)^2 = a^2 + 2ab + b^2$).
\item Replace every occurrence of $i^2$ by $-1$, then $i^3 = -i$, $i^4 = 1$, and use the cycle of powers of $i$ for higher exponents.
\item Group the \cle{real parts} together and the \cle{imaginary parts} together.
\item Write the final answer in the standard form $a + ib$, with $a, b \in \mathbb{R}$.
\end{enumerate}
\textbf{Reflex:} never leave $i^2$ in a final answer; for instance $2 + 3i^2$ must be simplified to $2 - 3 = -1$.
\end{methode}

\begin{aretenir}[The Cycle of Powers of $i$]
\[ i^1 = i, \qquad i^2 = -1, \qquad i^3 = -i, \qquad i^4 = 1, \qquad i^{n+4} = i^n. \]
To evaluate $i^n$, divide $n$ by $4$ and keep the remainder: $i^n = i^r$ where $n = 4q + r$, $r \in \{0,1,2,3\}$.
\end{aretenir}

\begin{exoresolu}[Computing with Complex Numbers]
Given $z_1 = 3 + 2i$ and $z_2 = 1 - 4i$, express each of the following in the form $a + ib$:
\begin{enumerate}
\item[(a)] $z_1 + z_2$ \qquad (b) $z_1 - 2z_2$ \qquad (c) $z_1 z_2$ \qquad (d) $z_1^2$ \qquad (e) $i^{2023}$.
\end{enumerate}
\tcblower
\textbf{Solution.}
\begin{enumerate}
\item[(a)] Add real parts and imaginary parts separately:
\[ z_1 + z_2 = (3 + 2i) + (1 - 4i) = (3 + 1) + (2 - 4)i = 4 - 2i. \]
\item[(b)] First compute $2z_2 = 2 - 8i$, then:
\[ z_1 - 2z_2 = (3 + 2i) - (2 - 8i) = (3 - 2) + (2 + 8)i = 1 + 10i. \]
\item[(c)] Expand as with ordinary brackets:
\begin{align*}
z_1 z_2 &= (3 + 2i)(1 - 4i) = 3(1) + 3(-4i) + 2i(1) + 2i(-4i) \\
&= 3 - 12i + 2i - 8i^2.
\end{align*}
Since $i^2 = -1$, we have $-8i^2 = +8$, so
\[ z_1 z_2 = 3 + 8 + (-12 + 2)i = 11 - 10i. \]
\item[(d)] Using the identity $(a+b)^2 = a^2 + 2ab + b^2$:
\[ z_1^2 = (3 + 2i)^2 = 9 + 12i + 4i^2 = 9 - 4 + 12i = 5 + 12i. \]
\item[(e)] The powers of $i$ cycle with period $4$. Dividing $2023$ by $4$: $2023 = 4 \times 505 + 3$, hence
\[ i^{2023} = i^{4 \times 505 + 3} = \left(i^4\right)^{505} \cdot i^3 = 1 \cdot (-i) = -i. \]
\end{enumerate}
\textbf{Check (c):} $(3+2i)(1-4i)$: real part $= 3 \times 1 - 2 \times (-4) = 3 + 8 = 11$ \checkmark; imaginary part $= 3 \times (-4) + 2 \times 1 = -10$ \checkmark.
\end{exoresolu}

\courssub{Skill 2: Dividing Complex Numbers Using the Conjugate}

Addition and multiplication of complex numbers behave exactly as in $\mathbb{R}$. Division is the first genuinely new operation: a quotient $\dfrac{z_1}{z_2}$ is not immediately in the form $a + ib$, because its denominator contains $i$. The remedy is the \cle{complex conjugate}.

\begin{definition}[Complex Conjugate]
The \cle{conjugate} of $z = a + ib$ ($a, b \in \mathbb{R}$) is the complex number
\[ \bar{z} = a - ib. \]
Geometrically, $\bar{z}$ is the reflection of $z$ in the real axis. Immediate properties: $z + \bar{z} = 2a = 2\mathrm{Re}(z)$, $z - \bar{z} = 2ib = 2i\,\mathrm{Im}(z)$, and
\[ z\,\bar{z} = (a + ib)(a - ib) = a^2 + b^2 = |z|^2, \]
which is always a \emph{non-negative real} number. Moreover $z$ is real if and only if $z = \bar{z}$, and $z$ is purely imaginary if and only if $z = -\bar{z}$ (with $z \neq 0$).
\end{definition}

\begin{methode}[Quotient of Two Complex Numbers]
To write $\dfrac{z_1}{z_2}$ (with $z_2 = c + id \neq 0$) in the form $a + ib$:
\begin{enumerate}
\item Write down the \cle{conjugate} of the denominator: $\overline{z_2} = c - id$ (change only the sign of the imaginary part).
\item Multiply \textbf{numerator and denominator} by $\overline{z_2}$:
\[ \frac{z_1}{z_2} = \frac{z_1 \, \overline{z_2}}{z_2 \, \overline{z_2}}. \]
\item The denominator becomes the \emph{real} number $z_2 \overline{z_2} = c^2 + d^2 = |z_2|^2$.
\item Expand the numerator, simplify, and split into real and imaginary parts.
\end{enumerate}
\textbf{Key formula:} $(c + id)(c - id) = c^2 + d^2$ (a sum of squares — always real and positive when $z_2 \neq 0$).
\end{methode}

\begin{exoresolu}[Division via the Conjugate]
Express $\dfrac{5 + i}{2 - 3i}$ in the form $a + ib$, where $a, b \in \mathbb{R}$.
\tcblower
\textbf{Solution.} The conjugate of the denominator $2 - 3i$ is $2 + 3i$. Multiply top and bottom by $2 + 3i$:
\[ \frac{5 + i}{2 - 3i} = \frac{(5 + i)(2 + 3i)}{(2 - 3i)(2 + 3i)}. \]
\textbf{Denominator:}
\[ (2 - 3i)(2 + 3i) = 2^2 + 3^2 = 4 + 9 = 13. \]
\textbf{Numerator:}
\[ (5 + i)(2 + 3i) = 10 + 15i + 2i + 3i^2 = 10 + 17i - 3 = 7 + 17i. \]
\textbf{Therefore:}
\[ \frac{5 + i}{2 - 3i} = \frac{7 + 17i}{13} = \frac{7}{13} + \frac{17}{13}i. \]
\textbf{Check:} multiply the answer by the denominator: $\left(\tfrac{7}{13} + \tfrac{17}{13}i\right)(2 - 3i) = \tfrac{1}{13}\left(14 - 21i + 34i - 51i^2\right) = \tfrac{1}{13}(14 + 51 + 13i) = \tfrac{65 + 13i}{13} = 5 + i$ \checkmark.
\end{exoresolu}

\begin{aretenir}[Algebraic Properties of the Conjugate]
For all $z_1, z_2 \in \mathbb{C}$ (with $z_2 \neq 0$ for the quotient):
\[ \overline{z_1 + z_2} = \bar{z}_1 + \bar{z}_2, \qquad \overline{z_1 z_2} = \bar{z}_1 \, \bar{z}_2, \qquad \overline{\left(\frac{z_1}{z_2}\right)} = \frac{\bar{z}_1}{\bar{z}_2}, \qquad \overline{z^n} = \bar{z}^{\,n}, \qquad \bar{\bar{z}} = z. \]
These properties are frequently examined: they justify, for example, why non-real roots of real polynomials come in conjugate pairs.
\end{aretenir}

\courssub{Skill 3: Solving Equations in $\mathbb{C}$}

One of the historical motivations for inventing $\mathbb{C}$ was precisely to give \emph{every} quadratic equation a solution. In $\mathbb{C}$, the quadratic formula remains valid even when the discriminant is negative.

\begin{methode}[Quadratic Equations with Negative Discriminant]
To solve $az^2 + bz + c = 0$ ($a, b, c \in \mathbb{R}$, $a \neq 0$) in $\mathbb{C}$:
\begin{enumerate}
\item Compute the discriminant $\Delta = b^2 - 4ac$.
\item If $\Delta < 0$, write $\Delta = -|\Delta|$ so that $\sqrt{\Delta} = i\sqrt{|\Delta|}$.
\item Apply the quadratic formula:
\[ z = \frac{-b \pm i\sqrt{|\Delta|}}{2a}. \]
\item The two roots are always \cle{conjugates} of each other when the coefficients are real.
\end{enumerate}
\textbf{Reflex:} for equations of the type $z = x + iy$ with conditions (e.g.\ $z^2$ real, or $\dfrac{z-1}{z+1}$ purely imaginary), substitute $z = x + iy$, expand, and \textbf{equate real and imaginary parts separately} to obtain simultaneous equations in $x$ and $y$. This is legitimate because two complex numbers are equal if and only if they have the same real part and the same imaginary part.
\end{methode}

\begin{exoresolu}[Solving a Quadratic and an Equation with Unknown Real and Imaginary Parts]
\begin{enumerate}
\item[(a)] Solve in $\mathbb{C}$: $z^2 - 4z + 13 = 0$.
\item[(b)] Find the complex number $z$ such that $z + 2\bar{z} = 6 - 3i$.
\end{enumerate}
\tcblower
\textbf{Solution (a).} Here $a = 1$, $b = -4$, $c = 13$:
\[ \Delta = (-4)^2 - 4(1)(13) = 16 - 52 = -36 < 0. \]
Hence $\sqrt{\Delta} = i\sqrt{36} = 6i$, and
\[ z = \frac{4 \pm 6i}{2} = 2 \pm 3i. \]
The solution set is $S = \{2 - 3i,\; 2 + 3i\}$. Note the roots are conjugates, as expected since the coefficients are real.
\smallskip

\textbf{Solution (b).} Write $z = x + iy$ with $x, y \in \mathbb{R}$; then $\bar{z} = x - iy$. Substituting:
\[ (x + iy) + 2(x - iy) = 6 - 3i \;\Longrightarrow\; 3x - iy = 6 - 3i. \]
Equating real parts: $3x = 6$, so $x = 2$. Equating imaginary parts: $-y = -3$, so $y = 3$.
\[ \boxed{z = 2 + 3i} \]
\textbf{Check:} $z + 2\bar{z} = (2 + 3i) + 2(2 - 3i) = 2 + 3i + 4 - 6i = 6 - 3i$ \checkmark.
\end{exoresolu}

\begin{attention}[Square Roots of a Complex Number]
To find the \cle{square roots} of a complex number $w$ (i.e.\ to solve $z^2 = w$), do \textbf{not} write ``$z = \sqrt{w}$'' — the symbol $\sqrt{\phantom{w}}$ is only defined for non-negative reals. Instead set $z = x + iy$, expand $z^2 = (x^2 - y^2) + 2ixy = w$, equate real and imaginary parts, and use the extra relation $x^2 + y^2 = |w|$ (obtained by taking moduli). You will always obtain exactly \textbf{two opposite} solutions $z$ and $-z$.
\end{attention}

\courssub{Skill 4: Modulus, Argument and Polar Form}

A complex number $z = a + ib$ is represented in the \cle{Argand diagram} by the point $M(a, b)$, called the \cle{affix} (or image) of $z$. The distance $OM = \sqrt{a^2 + b^2}$ is the \cle{modulus} $|z|$, and the angle $\theta$ between the positive real axis and $[OM)$ is an \cle{argument} of $z$, defined modulo $2\pi$. The \cle{principal argument}, written $\mathrm{Arg}(z)$, is the unique argument in $(-\pi, \pi]$.

\begin{methode}[Converting $z = a + ib$ to Polar Form $r(\cos\theta + i\sin\theta)$]
\begin{enumerate}
\item Compute the \cle{modulus}: $r = |z| = \sqrt{a^2 + b^2}$ (always $r \geq 0$).
\item Find an \cle{argument} $\theta$ from
\[ \cos\theta = \frac{a}{r}, \qquad \sin\theta = \frac{b}{r}. \]
Use \textbf{both} equations to fix the quadrant (a calculator gives only $\tan\theta = b/a$, which is ambiguous).
\item Write $z = r(\cos\theta + i\sin\theta)$, choosing the \cle{principal argument} in $(-\pi, \pi]$ if required.
\end{enumerate}
\textbf{Reflex:} sketch the point on the Argand diagram to confirm the quadrant. Remember: $|z_1 z_2| = |z_1||z_2|$, $\left|\dfrac{z_1}{z_2}\right| = \dfrac{|z_1|}{|z_2|}$, $\arg(z_1 z_2) = \arg z_1 + \arg z_2$ and $\arg\!\left(\dfrac{z_1}{z_2}\right) = \arg z_1 - \arg z_2$ (all modulo $2\pi$).
\end{methode}

\begin{exoresolu}[From Algebraic Form to Polar Form]
Given $z = -1 + i\sqrt{3}$:
\begin{enumerate}
\item[(a)] Find $|z|$ and the principal argument of $z$.
\item[(b)] Write $z$ in polar form and in exponential form.
\item[(c)] Hence write $\dfrac{z}{1 + i}$ in polar form.
\end{enumerate}
\tcblower
\textbf{Solution (a).} Modulus:
\[ r = |z| = \sqrt{(-1)^2 + (\sqrt{3})^2} = \sqrt{1 + 3} = \sqrt{4} = 2. \]
Argument: we need $\theta$ with
\[ \cos\theta = \frac{-1}{2}, \qquad \sin\theta = \frac{\sqrt{3}}{2}. \]
Cosine negative and sine positive place $\theta$ in the \textbf{second quadrant}. The reference angle is $\dfrac{\pi}{3}$, so
\[ \theta = \pi - \frac{\pi}{3} = \frac{2\pi}{3}. \]
\smallskip

\textbf{Solution (b).} Polar form and exponential form:
\[ z = 2\left(\cos\frac{2\pi}{3} + i\sin\frac{2\pi}{3}\right) = 2e^{i\frac{2\pi}{3}}. \]
\smallskip

\textbf{Solution (c).} For $w = 1 + i$: $|w| = \sqrt{1^2 + 1^2} = \sqrt{2}$, and $\arg w = \dfrac{\pi}{4}$ (first quadrant, $\cos\theta = \sin\theta = \tfrac{1}{\sqrt{2}}$). Using the quotient rules:
\[ \left|\frac{z}{w}\right| = \frac{|z|}{|w|} = \frac{2}{\sqrt{2}} = \sqrt{2}, \qquad \arg\frac{z}{w} = \frac{2\pi}{3} - \frac{\pi}{4} = \frac{8\pi - 3\pi}{12} = \frac{5\pi}{12}. \]
Therefore
\[ \frac{z}{1+i} = \sqrt{2}\left(\cos\frac{5\pi}{12} + i\sin\frac{5\pi}{12}\right). \]
Since $\tfrac{5\pi}{12} \in (-\pi, \pi]$, this is already the principal argument.
\end{exoresolu}

\courssub{Skill 5: Applying De Moivre's Theorem}

Computing $(a + ib)^n$ by repeated multiplication is hopeless for large $n$. The polar form turns multiplication into addition of angles, which is the content of De Moivre's theorem.

\begin{propriete}[De Moivre's Theorem]
For every real number $\theta$ and every integer $n \in \mathbb{Z}$:
\[ \bigl(\cos\theta + i\sin\theta\bigr)^n = \cos(n\theta) + i\sin(n\theta). \]
Consequently, if $z = r(\cos\theta + i\sin\theta)$, then $z^n = r^n\bigl(\cos(n\theta) + i\sin(n\theta)\bigr)$. (Proof by induction on $n \geq 0$ using the addition formulae; extension to negative $n$ via $z^{-1} = \bar{z}/|z|^2$. The proof is short and may be asked.)
\end{propriete}

\begin{methode}[Powers of Complex Numbers via De Moivre]
To compute $z^n$ for large $n$:
\begin{enumerate}
\item Convert $z$ to polar form $z = r(\cos\theta + i\sin\theta)$.
\item Apply \cle{De Moivre's theorem}:
\[ z^n = r^n\bigl(\cos(n\theta) + i\sin(n\theta)\bigr). \]
\item Reduce $n\theta$ modulo $2\pi$ to express the result with a standard angle, then convert back to algebraic form if required.
\end{enumerate}
\textbf{Reflex:} never expand $(a + ib)^n$ by the binomial theorem for large $n$ — De Moivre is always faster and safer.
\end{methode}

\begin{exoresolu}[Computing a High Power]
Let $z = 1 + i\sqrt{3}$. Compute $z^{12}$, giving your answer in the form $a + ib$.
\tcblower
\textbf{Solution.} \textbf{Step 1 — Polar form.}
\[ r = |z| = \sqrt{1 + 3} = 2; \qquad \cos\theta = \frac{1}{2},\; \sin\theta = \frac{\sqrt{3}}{2} \;\Longrightarrow\; \theta = \frac{\pi}{3}. \]
So $z = 2\left(\cos\dfrac{\pi}{3} + i\sin\dfrac{\pi}{3}\right)$.
\smallskip

\textbf{Step 2 — De Moivre.}
\[ z^{12} = 2^{12}\left(\cos\frac{12\pi}{3} + i\sin\frac{12\pi}{3}\right) = 4096\bigl(\cos 4\pi + i\sin 4\pi\bigr). \]
\smallskip

\textbf{Step 3 — Reduce the angle.} Since $4\pi$ is a multiple of $2\pi$:
\[ \cos 4\pi = 1, \qquad \sin 4\pi = 0. \]
Therefore
\[ z^{12} = 4096(1 + 0i) = 4096. \]
The result is a positive real number — consistent with the fact that $\arg z = \tfrac{\pi}{3}$ and $12 \times \tfrac{\pi}{3} = 4\pi \equiv 0 \pmod{2\pi}$.
\end{exoresolu}

\courssub{Skill 6: Trigonometric Applications of De Moivre}

De Moivre's theorem works in both directions: it converts \emph{powers} of $\cos x$ and $\sin x$ into \emph{multiple angles} (linearisation, useful for integration), and conversely it expands $\cos nx$ and $\sin nx$ as polynomials in $\cos x$ and $\sin x$.

\begin{methode}[Powers of $\cos x$ and $\sin x$ into Multiple Angles (and conversely)]
\textbf{(A) Powers $\to$ multiple angles.} Set $z = \cos x + i\sin x$, so $z + z^{-1} = 2\cos x$ and $z - z^{-1} = 2i\sin x$.
\begin{enumerate}
\item Write $\cos^n x = \left(\dfrac{z + z^{-1}}{2}\right)^n$ (or $\sin^n x = \left(\dfrac{z - z^{-1}}{2i}\right)^n$).
\item Expand with the binomial theorem and group symmetric terms $z^k + z^{-k} = 2\cos kx$.
\end{enumerate}
\textbf{(B) Multiple angles $\to$ powers.} Expand $(\cos x + i\sin x)^n$ by the binomial theorem and equate real parts (for $\cos nx$) or imaginary parts (for $\sin nx$).
\end{methode}

\begin{exoresolu}[Linearising $\cos^4 x$ and Expanding $\cos 3x$]
\begin{enumerate}
\item[(a)] Express $\cos^4 x$ as a sum of cosines of multiples of $x$.
\item[(b)] Express $\cos 3x$ in terms of powers of $\cos x$.
\end{enumerate}
\tcblower
\textbf{Solution (a).} Let $z = \cos x + i\sin x$, so $z + z^{-1} = 2\cos x$ and $z^k + z^{-k} = 2\cos kx$. Then
\[ (2\cos x)^4 = \left(z + z^{-1}\right)^4 = z^4 + 4z^2 + 6 + 4z^{-2} + z^{-4}. \]
Grouping symmetric terms:
\[ 16\cos^4 x = \left(z^4 + z^{-4}\right) + 4\left(z^2 + z^{-2}\right) + 6 = 2\cos 4x + 8\cos 2x + 6. \]
Dividing by $16$:
\[ \cos^4 x = \frac{1}{8}\cos 4x + \frac{1}{2}\cos 2x + \frac{3}{8}. \]
\textbf{Check} at $x = 0$: RHS $= \tfrac{1}{8} + \tfrac{1}{2} + \tfrac{3}{8} = 1 = \cos^4 0$ \checkmark.
\smallskip

\textbf{Solution (b).} By De Moivre and the binomial theorem:
\[ (\cos x + i\sin x)^3 = \cos^3 x + 3i\cos^2 x \sin x + 3i^2 \cos x \sin^2 x + i^3 \sin^3 x. \]
The real part is $\cos 3x$:
\[ \cos 3x = \cos^3 x - 3\cos x \sin^2 x = \cos^3 x - 3\cos x\left(1 - \cos^2 x\right). \]
Hence
\[ \cos 3x = 4\cos^3 x - 3\cos x. \]
\textbf{Check} at $x = \tfrac{\pi}{3}$: $\cos \pi = -1$ and $4\left(\tfrac{1}{2}\right)^3 - 3\left(\tfrac{1}{2}\right) = \tfrac{1}{2} - \tfrac{3}{2} = -1$ \checkmark.
\end{exoresolu}

\courssub{Skill 7: Finding the $n$th Roots of a Complex Number}

In $\mathbb{R}$, the equation $x^n = a$ has at most two solutions. In $\mathbb{C}$, the equation $z^n = w$ (with $w \neq 0$) always has \textbf{exactly $n$ distinct solutions}: this is one of the great strengths of the complex number system.

\begin{methode}[The $n$th Roots of $w = re^{i\theta}$]
\begin{enumerate}
\item Write the number in polar/exponential form: $w = r(\cos\theta + i\sin\theta)$.
\item The $n$ roots are given by
\[ z_k = r^{1/n}\left(\cos\frac{\theta + 2k\pi}{n} + i\sin\frac{\theta + 2k\pi}{n}\right), \qquad k = 0, 1, 2, \ldots, n-1. \]
\item List the roots for each value of $k$; stop at $k = n-1$ (further values repeat earlier roots).
\item \textbf{Geometry:} the roots lie on the circle of radius $r^{1/n}$, equally spaced at angles $\dfrac{2\pi}{n}$ — they are the vertices of a regular $n$-gon.
\end{enumerate}
\textbf{Special case (cube roots of unity):} the roots of $z^3 = 1$ are $1$, $\omega = e^{2i\pi/3} = -\tfrac{1}{2} + i\tfrac{\sqrt{3}}{2}$, $\omega^2 = \bar{\omega}$, with $1 + \omega + \omega^2 = 0$ and $\omega^3 = 1$.
\end{methode}

\begin{exoresolu}[Cube Roots of a Complex Number]
Find all the cube roots of $z = -8$, giving them in algebraic form, and represent them on an Argand diagram.
\tcblower
\textbf{Solution.} \textbf{Step 1 — Polar form.} The number $-8$ lies on the negative real axis:
\[ -8 = 8\bigl(\cos\pi + i\sin\pi\bigr), \qquad r = 8,\ \theta = \pi. \]
\textbf{Step 2 — Apply the formula} with $n = 3$, $r^{1/3} = 8^{1/3} = 2$:
\[ z_k = 2\left(\cos\frac{\pi + 2k\pi}{3} + i\sin\frac{\pi + 2k\pi}{3}\right), \quad k = 0, 1, 2. \]
\textbf{Step 3 — List the roots.}
\begin{align*}
k = 0:\quad & z_0 = 2\left(\cos\frac{\pi}{3} + i\sin\frac{\pi}{3}\right) = 2\left(\frac{1}{2} + i\frac{\sqrt{3}}{2}\right) = 1 + i\sqrt{3}.\\[2pt]
k = 1:\quad & z_1 = 2\left(\cos\pi + i\sin\pi\right) = -2.\\[2pt]
k = 2:\quad & z_2 = 2\left(\cos\frac{5\pi}{3} + i\sin\frac{5\pi}{3}\right) = 2\left(\frac{1}{2} - i\frac{\sqrt{3}}{2}\right) = 1 - i\sqrt{3}.
\end{align*}
\textbf{Check:} $(-2)^3 = -8$ \checkmark; $(1 + i\sqrt{3})^2 = 1 + 2i\sqrt{3} - 3 = -2 + 2i\sqrt{3}$, and $(1+i\sqrt{3})^3 = (-2 + 2i\sqrt{3})(1 + i\sqrt{3}) = -2 - 2i\sqrt{3} + 2i\sqrt{3} + 6i^2 = -2 - 6 = -8$ \checkmark.
\smallskip

The three roots have modulus $2$ and arguments $\tfrac{\pi}{3}$, $\pi$, $\tfrac{5\pi}{3}$: they form an \textbf{equilateral triangle} inscribed in the circle of radius $2$ centred at the origin.

\begin{popfigure}
\begin{tikzpicture}[scale=1.1]
\draw[->, gray] (-2.6,0) -- (2.8,0) node[right] {\small Re};
\draw[->, gray] (0,-2.4) -- (0,2.6) node[above] {\small Im};
\draw[popblue, thick] (0,0) circle (2);
\fill[poporange] (1,1.732) circle (2.2pt) node[above right, popdark] {\small $1+i\sqrt{3}$};
\fill[poporange] (-2,0) circle (2.2pt) node[above left, popdark] {\small $-2$};
\fill[poporange] (1,-1.732) circle (2.2pt) node[below right, popdark] {\small $1-i\sqrt{3}$};
\draw[popgreen, dashed] (1,1.732) -- (-2,0) -- (1,-1.732) -- cycle;
\draw[->, poppurple] (0.7,0) arc (0:60:0.7) node[midway, right] {\scriptsize $\frac{\pi}{3}$};
\end{tikzpicture}
\legende{The cube roots of $-8$: vertices of an equilateral triangle on the circle $|z| = 2$.}
\end{popfigure}
\end{exoresolu}

\courssub{Skill 8: Loci in the Complex Plane}

A \cle{locus} is the set of points $M(z)$ satisfying a given condition. Because $|z - a|$ measures the \emph{distance} between the points of affixes $z$ and $a$, and $\arg(z - a)$ measures a \emph{direction}, conditions on $z$ translate directly into geometry.

\begin{methode}[Identifying and Sketching Loci]
Translate the given condition on $z = x + iy$ into geometry:
\begin{enumerate}
\item $|z - a| = r$: \cle{circle}, centre $a$, radius $r$.
\item $|z - a| = |z - b|$: \cle{perpendicular bisector} of the segment joining $a$ and $b$ (points equidistant from $a$ and $b$).
\item $\arg(z - a) = \theta$: \cle{half-line} (ray) from $a$ (the point $a$ itself excluded) making angle $\theta$ with the positive real axis.
\item $|z - a| = k\,|z - b|$ with $k \neq 1$: circle of Apollonius — substitute $z = x + iy$ and rearrange to recognise centre and radius.
\end{enumerate}
\textbf{Reflex:} if the geometric meaning is unclear, set $z = x + iy$, square the moduli, and simplify to a Cartesian equation.
\end{methode}

\begin{exoresolu}[A Perpendicular Bisector and a Circle of Apollonius]
\begin{enumerate}
\item[(a)] Sketch the locus of points $z$ satisfying $|z - 2i| = |z - 4|$ and find its Cartesian equation.
\item[(b)] Show that the locus $|z| = 2|z - 3|$ is a circle; find its centre and radius.
\end{enumerate}
\tcblower
\textbf{Solution (a).} The condition says $z$ is equidistant from $A = 2i$ (point $(0,2)$) and $B = 4$ (point $(4,0)$): the locus is the \textbf{perpendicular bisector} of $[AB]$. Algebraically, with $z = x + iy$:
\[ |x + i(y-2)|^2 = |(x-4) + iy|^2 \;\Longrightarrow\; x^2 + (y-2)^2 = (x-4)^2 + y^2. \]
Expanding:
\[ x^2 + y^2 - 4y + 4 = x^2 - 8x + 16 + y^2 \;\Longrightarrow\; 8x - 4y = 12 \;\Longrightarrow\; 2x - y = 3. \]
The locus is the straight line $y = 2x - 3$ (it passes through the midpoint $(2,1)$ of $[AB]$ \checkmark).
\smallskip

\textbf{Solution (b).} Set $z = x + iy$ and square both sides:
\[ |z|^2 = 4|z - 3|^2 \;\Longrightarrow\; x^2 + y^2 = 4\left[(x-3)^2 + y^2\right]. \]
Expanding the right-hand side:
\[ x^2 + y^2 = 4x^2 - 24x + 36 + 4y^2 \;\Longrightarrow\; 3x^2 + 3y^2 - 24x + 36 = 0. \]
Divide by $3$:
\[ x^2 + y^2 - 8x + 12 = 0 \;\Longrightarrow\; (x - 4)^2 - 16 + y^2 + 12 = 0 \;\Longrightarrow\; (x-4)^2 + y^2 = 4. \]
This is a \textbf{circle of centre $(4, 0)$, i.e.\ the affix $4$, and radius $2$}.
\smallskip

\textbf{Check} with a point on the locus: $z = 2$ gives $|z| = 2$ and $2|z - 3| = 2 \times 1 = 2$ \checkmark; indeed $(2-4)^2 + 0 = 4$ \checkmark.
\end{exoresolu}

\begin{attention}[Frequent Errors at the GCE]
\begin{itemize}
\item Forgetting that $\sqrt{\Delta}$ for $\Delta < 0$ requires $i$: writing $\sqrt{-36} = -6$ is wrong; $\sqrt{-36} = 6i$.
\item Using only $\tan\theta = b/a$ to find the argument: this confuses opposite quadrants. Always use \textbf{both} $\cos\theta = a/r$ and $\sin\theta = b/r$.
\item Giving only one $n$th root: there are always exactly $n$ distinct roots, for $k = 0, 1, \ldots, n-1$.
\item In locus problems with $\arg(z - a) = \theta$, the point $a$ itself is \textbf{excluded} (the argument of $0$ is undefined).
\item Mixing degrees and radians: at Advanced Level, arguments are expected in \textbf{radians}.
\end{itemize}
\end{attention}

\newpage

\coursec{Exercises}

\courssub{I check my knowledge}

\begin{exercice}[Multiple choice questions]
For each question, only one answer is correct. Write down the letter of the correct answer.
\begin{enumerate}
\item The value of $i^{2023}$ is:
\begin{enumerate}
\item[(a)] $i$ \quad (b) $-i$ \quad (c) $1$ \quad (d) $-1$
\end{enumerate}
\item If $z = 3 - 4i$, then $|z|$ equals:
\begin{enumerate}
\item[(a)] $5$ \quad (b) $7$ \quad (c) $1$ \quad (d) $25$
\end{enumerate}
\item The conjugate of $z = -2 + 5i$ is:
\begin{enumerate}
\item[(a)] $2 - 5i$ \quad (b) $-2 - 5i$ \quad (c) $2 + 5i$ \quad (d) $5 - 2i$
\end{enumerate}
\item The exponential form of $z = 2\left(\cos\dfrac{\pi}{3} + i\sin\dfrac{\pi}{3}\right)$ is:
\begin{enumerate}
\item[(a)] $2e^{i\pi/3}$ \quad (b) $e^{2i\pi/3}$ \quad (c) $2e^{3i\pi}$ \quad (d)] $2e^{-i\pi/3}$
\end{enumerate}
\end{enumerate}
\end{exercice}

\begin{exercice}[True or false]
State whether each statement is true or false, justifying your answer briefly.
\begin{enumerate}
\item For every complex number $z$, $z + \bar{z}$ is a real number.
\item If $z = i$, then $\arg(z) = 0$.
\item The product of two complex conjugates is always a non-negative real number.
\item The equation $z^2 + 4 = 0$ has no solution in $\mathbb{C}$.
\item If $|z| = 0$, then $z = 0$.
\end{enumerate}
\end{exercice}

\begin{exercice}[Definitions and vocabulary]
\begin{enumerate}
\item Define the \cle{modulus} and the \cle{argument} of a non-zero complex number $z = a + ib$.
\item State \cle{De Moivre's theorem}.
\item What is meant by the \cle{principal argument} of a complex number? In which interval is it taken?
\item Explain what an \cle{Argand diagram} is and state what the axes represent.
\end{enumerate}
\end{exercice}

\begin{exercice}[Direct calculations]
\begin{enumerate}
\item Write each of the following in the form $a + ib$, where $a, b \in \mathbb{R}$:
\begin{enumerate}
\item $(2 + 3i) + (5 - 7i)$;
\item $(1 - 2i)(3 + 4i)$;
\item $\dfrac{3 + i}{1 - 2i}$.
\end{enumerate}
\item Find the modulus and the principal argument of $z = 1 - i$.
\item Write $z = -2i$ in polar form and in exponential form.
\item Compute $(1 + i)^2$ and deduce the square roots of $2i$.
\end{enumerate}
\end{exercice}

\courssub{I practise}

\begin{exercice}[Algebra of complex numbers]
Given $z_1 = 2 + i$ and $z_2 = 1 - 3i$, compute each of the following in the form $a + ib$:
\begin{enumerate}
\item $z_1 z_2$;
\item $\dfrac{z_1}{z_2}$;
\item $z_1^2 - 2\bar{z}_2$;
\item $\dfrac{z_1 + z_2}{z_1 - z_2}$.
\end{enumerate}
\end{exercice}

\begin{exercice}[Quadratic equations in $\mathbb{C}$]
\begin{enumerate}
\item Solve in $\mathbb{C}$ the equation $z^2 - 4z + 13 = 0$.
\item Represent the two roots $z_1$ and $z_2$ on an Argand diagram.
\item State the modulus and the principal argument of each root.
\item Verify that $z_1 + z_2 = 4$ and $z_1 z_2 = 13$.
\end{enumerate}
\end{exercice}

\begin{exercice}[Square roots of a complex number]
\begin{enumerate}
\item Find, in the form $a + ib$, the two square roots of $5 + 12i$.
\item Hence solve in $\mathbb{C}$ the equation
\[ z^2 - (3 + i)z + (2 - 2i) = 0. \]
\end{enumerate}
\end{exercice}

\begin{exercice}[Loci in the complex plane]
The complex number $z$ is represented by the point $M$ in the Argand diagram.
\begin{enumerate}
\item Describe and sketch the locus of $M$ in each case:
\begin{enumerate}
\item $|z - 2 + i| = 3$;
\item $|z - 1| = |z + 3i|$;
\item $\arg(z - 2) = \dfrac{\pi}{4}$.
\end{enumerate}
\item Find the Cartesian equation of the locus in part (b).
\end{enumerate}
\end{exercice}

\courssub{I prepare for the GCE}

\begin{exercice}[Polynomial with a complex root — GCE style]
Consider the polynomial $P(z) = z^3 - 5z^2 + 8z - 6$.
\begin{enumerate}
\item Show that $1 + i$ is a root of $P(z)$. \hfill (3 marks)
\item Write down, without further calculation, another complex root of $P(z)$, justifying your answer. \hfill (2 marks)
\item Hence find all the roots of $P(z) = 0$. \hfill (4 marks)
\item Represent the three roots on an Argand diagram and show that they form an isosceles triangle. \hfill (4 marks)
\end{enumerate}
\end{exercice}

\begin{exercice}[Polar form, De Moivre and powers — GCE style]
Let $z = -1 + i\sqrt{3}$.
\begin{enumerate}
\item Find the modulus and the principal argument of $z$, and write $z$ in polar form and in exponential form. \hfill (4 marks)
\item Using De Moivre's theorem, compute $z^{15}$ in the form $a + ib$. \hfill (4 marks)
\item Use De Moivre's theorem to prove that
\[ \cos 3x = 4\cos^3 x - 3\cos x. \] \hfill (4 marks)
\item Express $\cos^4 x$ as a sum of cosines of multiple angles. \hfill (4 marks)
\end{enumerate}
\end{exercice}

\begin{exercice}[Roots of unity and loci — GCE style]
Let $j = e^{2i\pi/3}$.
\begin{enumerate}
\item Prove that $j^3 = 1$ and that $1 + j + j^2 = 0$. \hfill (4 marks)
\item Hence solve in $\mathbb{C}$ the equation $z^3 = 8i$, giving all the roots in the form $a + ib$. \hfill (5 marks)
\item Find the four 4th roots of $-16$ in exponential form, plot them on an Argand diagram, and show that they are the vertices of a square. \hfill (5 marks)
\item On the same diagram, shade the region defined by
\[ 1 \leq |z| \leq 3 \qquad \text{and} \qquad -\frac{\pi}{4} \leq \arg z \leq \frac{\pi}{4}, \]
and find the exact area of this region. \hfill (4 marks)
\end{enumerate}
\end{exercice}

\newpage

\coursec{Solutions to the exercises}

\begin{corrige}[Exercise 1 — Multiple choice questions]
\begin{enumerate}
\item \textbf{(b) $-i$}. The powers of $i$ are periodic with period $4$: $i^1=i$, $i^2=-1$, $i^3=-i$, $i^4=1$, then the cycle repeats. Dividing $2023$ by $4$: $2023 = 4\times 505 + 3$, so $i^{2023} = i^{4\times 505}\cdot i^3 = (i^4)^{505}\cdot i^3 = 1^{505}\cdot(-i) = -i$.
\item \textbf{(a) $5$}. For $z = 3-4i$, $|z| = \sqrt{3^2 + (-4)^2} = \sqrt{9+16} = \sqrt{25} = 5$.
\item \textbf{(b) $-2-5i$}. The conjugate of $z = a+ib$ is $\bar{z} = a-ib$: only the sign of the \emph{imaginary} part changes. Here $a=-2$ and $b=5$, so $\bar{z} = -2 - 5i$.
\item \textbf{(a) $2e^{i\pi/3}$}. The polar form $r(\cos\theta + i\sin\theta)$ corresponds, by Euler's formula, to the exponential form $re^{i\theta}$. Here $r=2$ and $\theta = \pi/3$, so $z = 2e^{i\pi/3}$.
\end{enumerate}
\end{corrige}

\begin{corrige}[Exercise 2 — True or false]
\begin{enumerate}
\item \textbf{True}. Let $z = a+ib$ with $a,b\in\mathbb{R}$. Then $\bar{z} = a-ib$, so $z+\bar{z} = (a+ib)+(a-ib) = 2a \in \mathbb{R}$. In fact $z+\bar{z} = 2\operatorname{Re}(z)$.
\item \textbf{False}. $z = i$ corresponds to the point $(0,1)$ on the positive imaginary axis. Its modulus is $|i| = 1$ and its argument is $\frac{\pi}{2}$ (i.e. $90^\circ$), not $0$. An argument of $0$ characterises positive \emph{real} numbers.
\item \textbf{True}. If $z = a+ib$, then $\bar{z} = a-ib$ and
\[
z\bar{z} = (a+ib)(a-ib) = a^2 - (ib)^2 = a^2 + b^2 = |z|^2 \ge 0,
\]
a non-negative real number.
\item \textbf{False}. The equation $z^2+4=0$ gives $z^2 = -4$, hence $z = \pm\sqrt{-4} = \pm 2i$. These are perfectly valid \emph{complex} solutions: in $\mathbb{C}$ every quadratic equation has solutions.
\item \textbf{True}. $|z| = \sqrt{a^2+b^2} = 0 \iff a^2+b^2=0$. Since a sum of squares of real numbers is zero only when each term is zero, this is equivalent to $a=0$ and $b=0$, i.e. $z=0$.
\end{enumerate}
\end{corrige}

\begin{corrige}[Exercise 3 — Definitions and vocabulary]
\begin{enumerate}
\item For a non-zero complex number $z = a+ib$ ($a,b\in\mathbb{R}$):
\begin{itemize}
\item The \cle{modulus} of $z$, denoted $|z|$, is the distance from the origin $O$ to the image point $M(a,b)$ in the Argand diagram: $|z| = \sqrt{a^2+b^2}$.
\item An \cle{argument} of $z$, denoted $\arg(z)$, is any angle (in radians) that the vector $\overrightarrow{OM}$ makes with the positive real axis, measured anticlockwise. It is defined only up to the addition of multiples of $2\pi$.
\end{itemize}
\item \cle{De Moivre's theorem}: for any integer $n$ and any real number $\theta$,
\[
(\cos\theta + i\sin\theta)^n = \cos(n\theta) + i\sin(n\theta).
\]
Equivalently, if $z = r(\cos\theta + i\sin\theta)$, then $z^n = r^n\bigl(\cos(n\theta) + i\sin(n\theta)\bigr)$.
\item The \cle{principal argument} of a non-zero complex number $z$, denoted $\operatorname{Arg}(z)$, is the unique value of $\arg(z)$ lying in the interval $(-\pi, \pi]$.
\item An \cle{Argand diagram} is a geometric representation of complex numbers as points of a plane equipped with an orthonormal frame. The horizontal axis (the \emph{real axis}) carries the real part $\operatorname{Re}(z)$ and the vertical axis (the \emph{imaginary axis}) carries the imaginary part $\operatorname{Im}(z)$. The complex number $z = a+ib$ is represented by the point $M(a,b)$, called the \emph{image} of $z$, or equivalently by the vector $\overrightarrow{OM}$.
\end{enumerate}
\end{corrige}

\begin{corrige}[Exercise 4 — Direct calculations]
\begin{enumerate}
\item \begin{enumerate}
\item $(2+3i)+(5-7i) = (2+5) + i(3-7) = 7 - 4i$.
\item $(1-2i)(3+4i) = 3 + 4i - 6i - 8i^2 = 3 - 2i - 8(-1) = 11 - 2i$.
\item Multiply numerator and denominator by the conjugate of the denominator:
\[
\dfrac{3+i}{1-2i} = \dfrac{(3+i)(1+2i)}{(1-2i)(1+2i)} = \dfrac{3+6i+i+2i^2}{1^2+2^2} = \dfrac{3+7i-2}{5} = \dfrac{1+7i}{5} = \frac{1}{5} + \frac{7}{5}i.
\]
\end{enumerate}
\item For $z = 1-i$: $|z| = \sqrt{1^2+(-1)^2} = \sqrt{2}$. The image point $(1,-1)$ lies in the fourth quadrant, and $\tan\theta = \frac{-1}{1} = -1$, so the principal argument is $\operatorname{Arg}(z) = -\frac{\pi}{4}$.
\item $z = -2i$: modulus $|z| = \sqrt{0^2+(-2)^2} = 2$. The image point $(0,-2)$ lies on the negative imaginary axis, so $\operatorname{Arg}(z) = -\frac{\pi}{2}$. Polar form: $z = 2\left(\cos\left(-\frac{\pi}{2}\right) + i\sin\left(-\frac{\pi}{2}\right)\right)$. Exponential form: $z = 2e^{-i\pi/2}$.
\item $(1+i)^2 = 1 + 2i + i^2 = 1 + 2i - 1 = 2i$. Hence $1+i$ is a square root of $2i$, and the two square roots of $2i$ are $\pm(1+i)$, i.e. $1+i$ and $-1-i$.
\end{enumerate}
\end{corrige}

\begin{corrige}[Exercise 5 — Algebra of complex numbers]
Given $z_1 = 2+i$ and $z_2 = 1-3i$.
\begin{enumerate}
\item $z_1z_2 = (2+i)(1-3i) = 2 - 6i + i - 3i^2 = 2 - 5i + 3 = 5 - 5i$.
\item $\dfrac{z_1}{z_2} = \dfrac{2+i}{1-3i} = \dfrac{(2+i)(1+3i)}{(1-3i)(1+3i)} = \dfrac{2+6i+i+3i^2}{1^2+3^2} = \dfrac{2+7i-3}{10} = \dfrac{-1+7i}{10} = -0.1 + 0.7i$.
\item $z_1^2 = (2+i)^2 = 4 + 4i + i^2 = 3 + 4i$. Also $\bar{z}_2 = \overline{1-3i} = 1+3i$, so $2\bar{z}_2 = 2+6i$. Therefore
\[
z_1^2 - 2\bar{z}_2 = (3+4i) - (2+6i) = 1 - 2i.
\]
\item $z_1+z_2 = (2+i)+(1-3i) = 3-2i$ and $z_1-z_2 = (2+i)-(1-3i) = 1+4i$. Then
\[
\frac{z_1+z_2}{z_1-z_2} = \frac{3-2i}{1+4i} = \frac{(3-2i)(1-4i)}{(1+4i)(1-4i)} = \frac{3-12i-2i+8i^2}{1^2+4^2} = \frac{3-14i-8}{17} = \frac{-5-14i}{17} = -\frac{5}{17} - \frac{14}{17}i.
\]
\end{enumerate}
\end{corrige}

\begin{corrige}[Exercise 6 — Quadratic equations in $\mathbb{C}$]
\begin{enumerate}
\item Solve $z^2 - 4z + 13 = 0$. The discriminant is
\[
\Delta = (-4)^2 - 4(1)(13) = 16 - 52 = -36 = (6i)^2.
\]
Since $\Delta < 0$, the equation has two complex conjugate roots:
\[
z = \frac{-(-4) \pm 6i}{2(1)} = \frac{4 \pm 6i}{2} = 2 \pm 3i.
\]
So $z_1 = 2+3i$ and $z_2 = 2-3i$.
\item The images are the points $A(2,3)$ and $B(2,-3)$. They have the same real part and opposite imaginary parts, so they are symmetric with respect to the real axis — as expected for conjugate roots.
\begin{popfigure}
\begin{center}
\begin{tikzpicture}[scale=0.9]
\draw[->,popline] (-0.5,0) -- (4,0) node[below left] {\small Re};
\draw[->,popline] (0,-3.8) -- (0,3.8) node[below left] {\small Im};
\foreach \x in {1,2,3} \draw (\x,0.08) -- (\x,-0.08) node[below] {\scriptsize $\x$};
\foreach \y in {-3,-2,-1,1,2,3} \draw (0.08,\y) -- (-0.08,\y) node[left] {\scriptsize $\y$};
\draw[popblue,thick,->] (0,0) -- (2,3);
\draw[poporange,thick,->] (0,0) -- (2,-3);
\fill[popblue] (2,3) circle (2pt) node[above right] {\small $A(2,3)$};
\fill[poporange] (2,-3) circle (2pt) node[below right] {\small $B(2,-3)$};
\draw[dashed,popmuted] (2,3) -- (2,-3);
\end{tikzpicture}
\end{center}
\legende{Images of the conjugate roots $z_1 = 2+3i$ and $z_2 = 2-3i$: symmetry about the real axis.}
\end{popfigure}
\item For $z_1 = 2+3i$: $|z_1| = \sqrt{2^2+3^2} = \sqrt{13}$. The point $(2,3)$ is in the first quadrant, so $\operatorname{Arg}(z_1) = \arctan\!\left(\frac{3}{2}\right) \approx 0.983$ rad. For $z_2 = 2-3i$: $|z_2| = \sqrt{13}$ and, by conjugation, $\operatorname{Arg}(z_2) = -\arctan\!\left(\frac{3}{2}\right) \approx -0.983$ rad. Note that $|z_1| = |z_2|$ and $\operatorname{Arg}(z_2) = -\operatorname{Arg}(z_1)$, as always for conjugates.
\item Sum and product of the roots:
\[
z_1+z_2 = (2+3i)+(2-3i) = 4, \qquad z_1z_2 = (2+3i)(2-3i) = 2^2 + 3^2 = 13.
\]
These match the coefficients of the equation: for $z^2 - 4z + 13 = 0$ the sum of roots is $4$ and the product is $13$. Verified.
\end{enumerate}
\end{corrige}

\begin{corrige}[Exercise 7 — Square roots of a complex number]
\begin{enumerate}
\item Find the square roots of $5+12i$. Let $x+iy$ ($x,y\in\mathbb{R}$) satisfy $(x+iy)^2 = 5+12i$. Expanding, $x^2 - y^2 + 2ixy = 5 + 12i$, and equating real and imaginary parts:
\[
x^2 - y^2 = 5 \quad (1), \qquad 2xy = 12 \;\Rightarrow\; xy = 6 \quad (2).
\]
Taking moduli of $(x+iy)^2 = 5+12i$ gives $x^2+y^2 = |5+12i| = \sqrt{25+144} = 13 \quad (3)$.
Adding $(1)$ and $(3)$: $2x^2 = 18 \Rightarrow x^2 = 9 \Rightarrow x = \pm 3$. Subtracting: $2y^2 = 8 \Rightarrow y^2 = 4 \Rightarrow y = \pm 2$. Since $xy = 6 > 0$, $x$ and $y$ have the \emph{same} sign. The two square roots are
\[
3+2i \quad \text{and} \quad -3-2i.
\]
Check: $(3+2i)^2 = 9 + 12i + 4i^2 = 5 + 12i$. \checkmark
\item Solve $z^2 - (3+i)z + (2-2i) = 0$. The discriminant is
\[
\Delta = (3+i)^2 - 4(1)(2-2i) = (9 + 6i + i^2) - 8 + 8i = (8 + 6i) - 8 + 8i = 14i.
\]
We need a square root of $\Delta = 14i$. Observe that $(1+i)^2 = 2i$, so
\[
\bigl(\sqrt{7}(1+i)\bigr)^2 = 7(1+i)^2 = 7(2i) = 14i.
\]
Thus $\sqrt{\Delta} = \pm\sqrt{7}(1+i) = \pm(\sqrt{7} + i\sqrt{7})$. The roots are
\[
z = \frac{(3+i) \pm (\sqrt{7}+i\sqrt{7})}{2},
\]
that is
\[
z_1 = \frac{3+\sqrt{7}}{2} + i\,\frac{1+\sqrt{7}}{2}, \qquad z_2 = \frac{3-\sqrt{7}}{2} + i\,\frac{1-\sqrt{7}}{2}.
\]
\end{enumerate}
\end{corrige}

\begin{corrige}[Exercise 8 — Loci in the complex plane]
\begin{enumerate}
\item \begin{enumerate}
\item $|z - 2 + i| = 3 \iff |z - (2-i)| = 3$. The distance from the image of $z$ to the point $(2,-1)$ is constantly $3$: the locus is the \textbf{circle} with centre $2-i$ (point $(2,-1)$) and radius $3$.
\item $|z-1| = |z+3i| \iff |z-1| = |z-(-3i)|$. The image of $z$ is equidistant from the points $A(1,0)$ and $B(0,-3)$: the locus is the \textbf{perpendicular bisector} of the segment $[AB]$. Its midpoint is $\left(\frac{1}{2}, -\frac{3}{2}\right)$; the gradient of $[AB]$ is $\frac{-3-0}{0-1} = 3$, so the perpendicular gradient is $-\frac{1}{3}$. Equation: $y + \frac{3}{2} = -\frac{1}{3}\left(x - \frac{1}{2}\right)$.
\item $\arg(z-2) = \frac{\pi}{4}$. Writing $w = z - 2$, the condition $\arg w = \frac{\pi}{4}$ describes a \textbf{half-line (ray)} starting from the point $2$ (i.e. $(2,0)$), making an angle $\frac{\pi}{4}$ with the positive real axis. The point $(2,0)$ itself is \emph{excluded}, because $\arg(0)$ is undefined.
\end{enumerate}
\item Cartesian equation for (b): let $z = x+iy$. Squaring both sides of $|z-1| = |z+3i|$:
\[
(x-1)^2 + y^2 = x^2 + (y+3)^2.
\]
Expanding: $x^2 - 2x + 1 + y^2 = x^2 + y^2 + 6y + 9$. Cancelling $x^2 + y^2$:
\[
-2x + 1 = 6y + 9 \;\Rightarrow\; -2x - 6y = 8 \;\Rightarrow\; x + 3y + 4 = 0.
\]
This is indeed the straight line found geometrically in (b): substituting the midpoint $\left(\frac{1}{2}, -\frac{3}{2}\right)$ gives $\frac{1}{2} - \frac{9}{2} + 4 = 0$. \checkmark
\end{enumerate}
\end{corrige}

\begin{corrige}[Exercise 9 — Polynomial with a complex root — GCE style]
\begin{enumerate}
\item \textbf{Show that $1+i$ is a root.} Substitute $z = 1+i$ into $P(z) = z^3 - 5z^2 + 8z - 6$. First compute the powers:
\[
(1+i)^2 = 1 + 2i + i^2 = 2i, \qquad (1+i)^3 = (1+i)(1+i)^2 = (1+i)(2i) = 2i + 2i^2 = -2 + 2i.
\]
Then
\[
P(1+i) = (-2+2i) - 5(2i) + 8(1+i) - 6 = (-2 + 8 - 6) + (2i - 10i + 8i) = 0 + 0i = 0.
\]
Hence $1+i$ is a root of $P$. \hfill (3 marks)
\item \textbf{Another root.} The coefficients of $P$ are all real, so by the \cle{conjugate root theorem} the complex roots occur in conjugate pairs. Therefore $\overline{1+i} = 1-i$ is also a root. \hfill (2 marks)
\item \textbf{Find all the roots.} From the two roots $1+i$ and $1-i$ we build a real quadratic factor:
\[
\bigl(z-(1+i)\bigr)\bigl(z-(1-i)\bigr) = z^2 - \bigl((1+i)+(1-i)\bigr)z + (1+i)(1-i) = z^2 - 2z + 2.
\]
Since $P$ is monic of degree $3$, the remaining factor is linear: $P(z) = (z^2 - 2z + 2)(z - a)$. Comparing constant terms: $-6 = 2\times(-a)$, so $a = 3$. Check by expansion:
\[
(z^2 - 2z + 2)(z - 3) = z^3 - 3z^2 - 2z^2 + 6z + 2z - 6 = z^3 - 5z^2 + 8z - 6. \checkmark
\]
The three roots are $z = 1+i$, $z = 1-i$ and $z = 3$. \hfill (4 marks)
\item \textbf{Argand diagram and isosceles triangle.} The images are $A(1,1)$, $B(1,-1)$ and $C(3,0)$.
\begin{popfigure}
\begin{center}
\begin{tikzpicture}[scale=1.0]
\draw[->,popline] (-0.5,0) -- (4,0) node[below left] {\small Re};
\draw[->,popline] (0,-2) -- (0,2) node[below left] {\small Im};
\foreach \x in {1,2,3} \draw (\x,0.07) -- (\x,-0.07) node[below] {\scriptsize $\x$};
\foreach \y in {-1,1} \draw (0.07,\y) -- (-0.07,\y) node[left] {\scriptsize $\y$};
\draw[popblue,thick] (1,1) -- (1,-1) -- (3,0) -- cycle;
\fill[popblue] (1,1) circle (2pt) node[above left] {\small $A(1,1)$};
\fill[popblue] (1,-1) circle (2pt) node[below left] {\small $B(1,-1)$};
\fill[poporange] (3,0) circle (2pt) node[above right] {\small $C(3,0)$};
\end{tikzpicture}
\end{center}
\legende{Triangle formed by the images of the three roots of $P$.}
\end{popfigure}
Distances:
\[
AB = |1-(-1)| = 2, \quad AC = \sqrt{(3-1)^2 + (0-1)^2} = \sqrt{5}, \quad BC = \sqrt{(3-1)^2 + (0-(-1))^2} = \sqrt{5}.
\]
Since $AC = BC = \sqrt{5}$, the triangle $ABC$ is isosceles with apex $C$. \hfill (4 marks)
\end{enumerate}
\textbf{Mark scheme / Examiner expectations:}
\begin{itemize}
\item Part (a): correct values of $(1+i)^2$ and $(1+i)^3$ (2 marks), substitution and simplification to $0$ (1 mark).
\item Part (b): explicit statement of the conjugate root theorem (1 mark), root $1-i$ (1 mark).
\item Part (c): quadratic factor $z^2 - 2z + 2$ (1 mark), correct division or comparison of coefficients (2 marks), third root $z = 3$ (1 mark).
\item Part (d): correct plot (1 mark), correct distance computations (2 marks), conclusion isosceles (1 mark).
\end{itemize}
\end{corrige}

\begin{corrige}[Exercise 10 — Polar form, De Moivre and powers — GCE style]
Let $z = -1 + i\sqrt{3}$.
\begin{enumerate}
\item \textbf{Modulus and argument.} $|z| = \sqrt{(-1)^2 + (\sqrt{3})^2} = \sqrt{1+3} = 2$. The image point $(-1, \sqrt{3})$ lies in the \emph{second} quadrant (negative real part, positive imaginary part). We have $\cos\theta = \frac{-1}{2}$ and $\sin\theta = \frac{\sqrt{3}}{2}$, so the principal argument is $\theta = \pi - \frac{\pi}{3} = \frac{2\pi}{3}$. Polar form:
\[
z = 2\left(\cos\frac{2\pi}{3} + i\sin\frac{2\pi}{3}\right), \qquad \text{exponential form: } z = 2e^{2i\pi/3}.
\]
\hfill (4 marks)
\item \textbf{Compute $z^{15}$ using De Moivre's theorem.}
\[
z^{15} = 2^{15}\left(\cos\left(15 \times \frac{2\pi}{3}\right) + i\sin\left(15 \times \frac{2\pi}{3}\right)\right) = 32768\bigl(\cos(10\pi) + i\sin(10\pi)\bigr).
\]
Since $10\pi$ is an even multiple of $\pi$, $\cos(10\pi) = 1$ and $\sin(10\pi) = 0$. Therefore $z^{15} = 32768$, a positive real number. \hfill (4 marks)
\item \textbf{Prove that $\cos 3x = 4\cos^3 x - 3\cos x$.} By De Moivre's theorem with $n = 3$:
\[
\cos 3x + i\sin 3x = (\cos x + i\sin x)^3.
\]
Expanding the right-hand side with the binomial theorem:
\[
(\cos x + i\sin x)^3 = \cos^3 x + 3i\cos^2 x \sin x + 3i^2\cos x \sin^2 x + i^3\sin^3 x = \cos^3 x + 3i\cos^2 x\sin x - 3\cos x \sin^2 x - i\sin^3 x.
\]
Equating the real parts:
\[
\cos 3x = \cos^3 x - 3\cos x \sin^2 x = \cos^3 x - 3\cos x(1 - \cos^2 x) = 4\cos^3 x - 3\cos x. \qquad \blacksquare
\]
\hfill (4 marks)
\item \textbf{Express $\cos^4 x$ as a sum of cosines of multiple angles.} From Euler's formulae, $2\cos x = e^{ix} + e^{-ix}$. Raising to the fourth power and expanding binomially:
\[
2^4\cos^4 x = (e^{ix} + e^{-ix})^4 = e^{4ix} + 4e^{2ix} + 6 + 4e^{-2ix} + e^{-4ix}.
\]
Grouping conjugate terms and using $e^{ikx} + e^{-ikx} = 2\cos kx$:
\[
16\cos^4 x = \bigl(e^{4ix} + e^{-4ix}\bigr) + 4\bigl(e^{2ix} + e^{-2ix}\bigr) + 6 = 2\cos 4x + 8\cos 2x + 6.
\]
Hence
\[
\cos^4 x = \frac{1}{8}\cos 4x + \frac{1}{2}\cos 2x + \frac{3}{8}.
\]
\hfill (4 marks)
\end{enumerate}
\textbf{Mark scheme / Examiner expectations:}
\begin{itemize}
\item Part (a): correct modulus (1 mark), correct quadrant identified and argument $\frac{2\pi}{3}$ (2 marks), correct polar and exponential forms (1 mark).
\item Part (b): correct application of De Moivre (2 marks), reduction of the angle $10\pi$ (1 mark), final value $32768$ (1 mark).
\item Part (c): De Moivre stated for $n = 3$ (1 mark), correct binomial expansion (1 mark), equating real parts and using $\sin^2 x = 1 - \cos^2 x$ (2 marks).
\item Part (d): use of $2\cos x = e^{ix} + e^{-ix}$ (1 mark), correct expansion of $(e^{ix} + e^{-ix})^4$ (2 marks), simplification to the required form (1 mark).
\end{itemize}
\end{corrige}

\begin{corrige}[Exercise 11 — Roots of unity and loci — GCE style]
Let $j = e^{2i\pi/3}$.
\begin{enumerate}
\item \textbf{Prove that $j^3 = 1$ and $1 + j + j^2 = 0$.} By the laws of exponents:
\[
j^3 = \left(e^{2i\pi/3}\right)^3 = e^{2i\pi} = \cos 2\pi + i\sin 2\pi = 1.
\]
Since $j \neq 1$ (indeed $j = -\frac{1}{2} + i\frac{\sqrt{3}}{2}$), factorising $j^3 - 1 = 0$ gives
\[
j^3 - 1 = (j - 1)(j^2 + j + 1) = 0 \;\Rightarrow\; j^2 + j + 1 = 0,
\]
because the factor $j - 1$ is non-zero. Hence $1 + j + j^2 = 0$. \hfill (4 marks)
\item \textbf{Solve $z^3 = 8i$.} Write the right-hand side in exponential form: $8i = 8e^{i\pi/2}$. The three cube roots are
\[
z_k = \sqrt[3]{8}\; e^{i(\pi/2 + 2k\pi)/3} = 2\, e^{i(\pi/6 + 2k\pi/3)}, \qquad k = 0, 1, 2.
\]
Computing each root in Cartesian form:
\begin{align*}
k=0:&\quad z_0 = 2e^{i\pi/6} = 2\left(\cos\frac{\pi}{6} + i\sin\frac{\pi}{6}\right) = 2\left(\frac{\sqrt{3}}{2} + \frac{1}{2}i\right) = \sqrt{3} + i,\\
k=1:&\quad z_1 = 2e^{i5\pi/6} = 2\left(\cos\frac{5\pi}{6} + i\sin\frac{5\pi}{6}\right) = 2\left(-\frac{\sqrt{3}}{2} + \frac{1}{2}i\right) = -\sqrt{3} + i,\\
k=2:&\quad z_2 = 2e^{i3\pi/2} = 2\left(\cos\frac{3\pi}{2} + i\sin\frac{3\pi}{2}\right) = 2(0 - i) = -2i.
\end{align*}
The three roots are $\sqrt{3}+i$, $-\sqrt{3}+i$ and $-2i$. \hfill (5 marks)
\item \textbf{Fourth roots of $-16$.} Write $-16 = 16e^{i\pi}$. The four fourth roots are
\[
z_k = \sqrt[4]{16}\; e^{i(\pi + 2k\pi)/4} = 2\, e^{i(\pi/4 + k\pi/2)}, \qquad k = 0, 1, 2, 3,
\]
that is $z_0 = 2e^{i\pi/4}$, $z_1 = 2e^{i3\pi/4}$, $z_2 = 2e^{i5\pi/4}$, $z_3 = 2e^{i7\pi/4}$.
\begin{popfigure}
\begin{center}
\begin{tikzpicture}[scale=1.1]
\draw[->,popline] (-2.8,0) -- (2.8,0) node[below left] {\small Re};
\draw[->,popline] (0,-2.8) -- (0,2.8) node[below left] {\small Im};
\draw[popmuted] (0,0) circle (2);
\draw[popblue,thick] (1.414,1.414) -- (-1.414,1.414) -- (-1.414,-1.414) -- (1.414,-1.414) -- cycle;
\fill[popblue] (1.414,1.414) circle (2pt) node[above right] {\small $z_0$};
\fill[popblue] (-1.414,1.414) circle (2pt) node[above left] {\small $z_1$};
\fill[popblue] (-1.414,-1.414) circle (2pt) node[below left] {\small $z_2$};
\fill[popblue] (1.414,-1.414) circle (2pt) node[below right] {\small $z_3$};
\draw[poporange,->] (0.8,0) arc (0:45:0.8) node[midway,right] {\scriptsize $\pi/4$};
\end{tikzpicture}
\end{center}
\legende{The four fourth roots of $-16$: vertices of a square inscribed in the circle of radius $2$.}
\end{popfigure}
All four points lie on the circle of radius $2$ centred at the origin, and consecutive arguments differ by exactly $\frac{\pi}{2}$. Hence the four images are the vertices of a \textbf{square} (a regular quadrilateral) inscribed in that circle. \hfill (5 marks)
\item \textbf{Region defined by $1 \le |z| \le 3$ and $-\frac{\pi}{4} \le \arg z \le \frac{\pi}{4}$.} The condition $1 \le |z| \le 3$ describes the \emph{annulus} between the circles of radii $1$ and $3$ centred at the origin; the condition on the argument restricts to the \emph{sector} between the rays $\arg z = -\frac{\pi}{4}$ and $\arg z = \frac{\pi}{4}$. The locus is therefore a sector of an annulus, of angle $\frac{\pi}{4} - \left(-\frac{\pi}{4}\right) = \frac{\pi}{2}$.
\begin{popfigure}
\begin{center}
\begin{tikzpicture}[scale=0.85]
\draw[->,popline] (-1.5,0) -- (4,0) node[below left] {\small Re};
\draw[->,popline] (0,-2.6) -- (0,2.6) node[below left] {\small Im};
\fill[popblueL] (0.707,0.707) arc (45:-45:1) -- (2.121,-2.121) arc (-45:45:3) -- cycle;
\draw[popblue,thick] (0.707,0.707) arc (45:-45:1);
\draw[popblue,thick] (2.121,2.121) arc (45:-45:3);
\draw[popblue,thick] (0.707,0.707) -- (2.121,2.121);
\draw[popblue,thick] (0.707,-0.707) -- (2.121,-2.121);
\draw[dashed,popmuted] (0,0) -- (3.4,3.4) node[above] {\scriptsize $\arg z = \pi/4$};
\draw[dashed,popmuted] (0,0) -- (3.4,-3.4) node[below] {\scriptsize $\arg z = -\pi/4$};
\node[popdark] at (1.9,0) {\scriptsize region};
\end{tikzpicture}
\end{center}
\legende{Sector of an annulus: $1 \le |z| \le 3$ and $|\arg z| \le \pi/4$.}
\end{popfigure}
The area of a sector of angle $\alpha$ of an annulus of radii $r$ and $R$ is $\frac{1}{2}\alpha(R^2 - r^2)$. Here:
\[
\text{Area} = \frac{1}{2} \times \frac{\pi}{2}\times(3^2 - 1^2) = \frac{\pi}{4}(9 - 1) = 2\pi \text{ square units}.
\]
\hfill (4 marks)
\end{enumerate}
\textbf{Mark scheme / Examiner expectations:}
\begin{itemize}
\item Part (a): $j^3 = 1$ obtained from the laws of exponents (2 marks), factorisation of $j^3 - 1$ and conclusion $1 + j + j^2 = 0$ with the justification $j \neq 1$ (2 marks).
\item Part (b): $8i$ written in exponential form (1 mark), correct formula for the cube roots (2 marks), the three roots given in $a + ib$ form (2 marks).
\item Part (c): $-16$ written in exponential form (1 mark), the four roots in exponential form (2 marks), correct plot (1 mark), justification of the square (equal moduli, arguments differing by $\frac{\pi}{2}$) (1 mark).
\item Part (d): region correctly described (1 mark), correct sketch (1 mark), area formula for a sector of an annulus (1 mark), correct value $2\pi$ (1 mark).
\end{itemize}
\end{corrige}

\newpage

\coursec{Summary sheet}

\begin{popfigure}
\begin{center}\tcbox[colback=popdark,colframe=popdark,coltext=white,arc=9pt,boxsep=4pt,left=6pt,right=6pt]{\bfseries\small Complex Numbers}\end{center}
\vspace{-2pt}
\begin{tcbraster}[raster columns=3,raster equal height=rows,raster column skip=6pt,raster row skip=6pt,enhanced,blankest]
\begin{tcolorbox}[enhanced,colback=popblue,colframe=popblue,coltext=white,boxrule=0.9pt,arc=3pt,left=4pt,right=4pt,top=3pt,bottom=3pt,boxsep=1pt,halign=left]\bfseries\scriptsize Birth of $\mathbb{C}$; $i^2=-1$; algebra of $a+ib$\end{tcolorbox}
\begin{tcolorbox}[enhanced,colback=popgreen,colframe=popgreen,coltext=white,boxrule=0.9pt,arc=3pt,left=4pt,right=4pt,top=3pt,bottom=3pt,boxsep=1pt,halign=left]\bfseries\scriptsize Conjugate $\bar z$; division; properties\end{tcolorbox}
\begin{tcolorbox}[enhanced,colback=poporange,colframe=poporange,coltext=white,boxrule=0.9pt,arc=3pt,left=4pt,right=4pt,top=3pt,bottom=3pt,boxsep=1pt,halign=left]\bfseries\scriptsize Equations; square roots in $\mathbb{C}$\end{tcolorbox}
\begin{tcolorbox}[enhanced,colback=poppurple,colframe=poppurple,coltext=white,boxrule=0.9pt,arc=3pt,left=4pt,right=4pt,top=3pt,bottom=3pt,boxsep=1pt,halign=left]\bfseries\scriptsize Argand diagram; $|z|$, $\arg z$\end{tcolorbox}
\begin{tcolorbox}[enhanced,colback=poppink,colframe=poppink,coltext=white,boxrule=0.9pt,arc=3pt,left=4pt,right=4pt,top=3pt,bottom=3pt,boxsep=1pt,halign=left]\bfseries\scriptsize Polar $r(\cos\theta+i\sin\theta)$; $re^{i\theta}$\end{tcolorbox}
\begin{tcolorbox}[enhanced,colback=popgold,colframe=popgold,coltext=white,boxrule=0.9pt,arc=3pt,left=4pt,right=4pt,top=3pt,bottom=3pt,boxsep=1pt,halign=left]\bfseries\scriptsize De Moivre; trig.\ applications; $n$th roots; loci\end{tcolorbox}
\end{tcbraster}

\legende{Mind map of the chapter \emph{Complex Numbers}.}
\end{popfigure}

\begin{bilan}
\textbf{Key definitions}
\begin{itemize}
\item A \cle{complex number} is $z=a+ib$, $a,b\in\mathbb{R}$, where $i^2=-1$; $a=\mathrm{Re}(z)$, $b=\mathrm{Im}(z)$. If $b=0$, $z$ is real; if $a=0$, $z$ is pure imaginary.
\item \cle{Conjugate}: $\bar z=a-ib$; reflection of $z$ in the real axis.
\item \cle{Modulus}: $|z|=\sqrt{a^2+b^2}$ (distance from $O$). \cle{Argument}: $\arg z=\theta$ with $\cos\theta=\dfrac{a}{r}$, $\sin\theta=\dfrac{b}{r}$; principal value $-\pi<\theta\le\pi$.
\item \cle{Polar form}: $z=r(\cos\theta+i\sin\theta)$; \cle{exponential form}: $z=re^{i\theta}$.
\end{itemize}

\textbf{Laws, formulas and theorems}
\begin{itemize}
\item Powers of $i$ cycle: $i^1=i$, $i^2=-1$, $i^3=-i$, $i^4=1$.
\item $z\bar z=|z|^2$; $\overline{z_1\pm z_2}=\bar z_1\pm\bar z_2$; $\overline{z_1z_2}=\bar z_1\bar z_2$; $\overline{\left(\dfrac{z_1}{z_2}\right)}=\dfrac{\bar z_1}{\bar z_2}$.
\item $z+\bar z=2\mathrm{Re}(z)$, $z-\bar z=2i\,\mathrm{Im}(z)$; $z$ real $\iff z=\bar z$.
\item Division: multiply numerator and denominator by the conjugate of the denominator.
\item $|z_1z_2|=|z_1||z_2|$; $\left|\dfrac{z_1}{z_2}\right|=\dfrac{|z_1|}{|z_2|}$; $\arg(z_1z_2)=\arg z_1+\arg z_2$; $\arg\!\left(\dfrac{z_1}{z_2}\right)=\arg z_1-\arg z_2$ (mod $2\pi$).
\item \cle{De Moivre}: $(\cos\theta+i\sin\theta)^n=\cos n\theta+i\sin n\theta$, $n\in\mathbb{Z}$ (proof by induction for $n\in\mathbb{N}$ is examinable).
\item \cle{$n$th roots}: $z^{1/n}=r^{1/n}\!\left(\cos\dfrac{\theta+2k\pi}{n}+i\sin\dfrac{\theta+2k\pi}{n}\right)$, $k=0,1,\dots,n-1$; the roots lie on a circle, vertices of a regular $n$-gon.
\item Euler: $\cos\theta=\dfrac{e^{i\theta}+e^{-i\theta}}{2}$, $\sin\theta=\dfrac{e^{i\theta}-e^{-i\theta}}{2i}$ (to express $\cos^n\theta$, $\sin^n\theta$ in multiple angles, and conversely).
\item Quadratic with real coefficients: non-real roots occur in conjugate pairs; $z=\dfrac{-b\pm i\sqrt{4ac-b^2}}{2a}$ when $\Delta<0$.
\end{itemize}

\textbf{Methods}
\begin{itemize}
\item \textbf{Argument}: use $\tan\theta=\dfrac{b}{a}$ \emph{plus the quadrant} of $(a,b)$; sketch the point first.
\item \textbf{Square roots of $a+ib$}: set $(x+iy)^2=a+ib$, equate parts, use $x^2+y^2=\sqrt{a^2+b^2}$.
\item \textbf{Multiple angles}: expand $(\cos\theta+i\sin\theta)^n$ by the binomial theorem, then equate real/imaginary parts.
\item \textbf{Powers to multiple angles}: write $z=e^{i\theta}$, expand $\left(z+\dfrac1z\right)^n$ or $\left(z-\dfrac1z\right)^n$.
\item \textbf{Loci}: $|z-a|=r$ circle; $|z-a|=|z-b|$ perpendicular bisector; $\arg(z-a)=\theta$ half-line; $|z-a|=k|z-b|$ circle of Apollonius.
\end{itemize}

\textbf{Common mistakes}
\begin{itemize}
\item Writing $\arg z=\tan^{-1}(b/a)$ without checking the quadrant (e.g.\ $-1+i$ has $\arg=\tfrac{3\pi}{4}$, not $-\tfrac{\pi}{4}$).
\item Forgetting the $2k\pi$ in $n$th roots, hence finding only one root.
\item Dividing without multiplying by the conjugate of the \emph{whole} denominator.
\item Confusing $|z|^2=z\bar z$ with $z^2$.
\item Mixing degrees and radians; GCE answers are expected in radians, exact form.
\end{itemize}

\textbf{GCE tips}
\begin{itemize}
\item State answers in exact form ($\pi/3$, $\sqrt2$, $e^{i\pi/6}$); no decimals unless asked.
\item For loci questions, always \emph{identify, describe and sketch} the locus on an Argand diagram.
\item Show the binomial expansion step in De Moivre applications --- method marks are awarded for it.
\item Check roots by substitution or by verifying the product/sum relations.
\end{itemize}
\end{bilan}

\findecours

\end{document}
